\section{The row package, automorphic counting and period entries}\label{sec:pf:rows}

\providecommand{\rowsRow}{\sfname{Row}}
\providecommand{\rowsRowZ}{\sfname{Row}_0}
\providecommand{\rowsAtAll}{\mathrm{All}}
\providecommand{\rowsEvP}{\mathrm{Ev}_{\mathrm{par}}}
\providecommand{\rowsEvC}{\mathcal{E}}
\providecommand{\rowsLBC}{\mathrm{LB}}
\providecommand{\rowsLg}{\mathrm{Lg}}
\providecommand{\rowsLinked}{\sfname{LinkedAt}}
\providecommand{\rowsClosed}{\sfname{Closed}}
\providecommand{\rowsBuilt}{\sfname{IsBuiltE}}
\providecommand{\rowsGood}{\sfname{IsGoodE}}
\providecommand{\rowsKPx}{\mathrm{KP}_x}
\providecommand{\rowsDiii}{\mathcal{D}_3}
\providecommand{\rowsWi}{\mathcal{W}_1}
\providecommand{\rowsWiii}{\mathcal{W}_3}
\providecommand{\rowsLk}{\mathrm{Lk}}
\providecommand{\rowsGA}{G_{\A}}
\providecommand{\rowsGinf}{G_\infty}
\providecommand{\rowsGf}{G_{\mathrm f}}
\providecommand{\rowsDst}{D_*}
\providecommand{\rowsDc}{D}
\providecommand{\rowsCH}{\mathrm{CH}}
\providecommand{\rowsSWone}{\mathrm{SW}_1}
\providecommand{\rowsgv}{\operatorname{gvol}}
\providecommand{\rowsvU}{v_{\Uni}}
\providecommand{\rowsmr}{\operatorname{mr}}
\providecommand{\rowsrf}{\operatorname{rf}}
\providecommand{\rowsmc}{\operatorname{mc}}
\providecommand{\rowsleast}{\operatorname{least}}
\providecommand{\rowssize}{\operatorname{size}}
\providecommand{\rowssizeQ}{\operatorname{size}}
\providecommand{\rowsdHS}{d_{\mathrm{HS}}}
\providecommand{\rowsdOp}{d_{\mathrm{op}}}
\providecommand{\rowskap}{\kappa}
\providecommand{\rowsdens}{\operatorname{dens}}
\providecommand{\rowsdimN}{\nu}
\providecommand{\rowsdeep}{\mathcal{X}_{\mathrm{deep}}}
\providecommand{\rowsmod}{\mathcal{X}_{\mathrm{mod}}}
\providecommand{\rowsLv}{\Lambda}
\providecommand{\rowsRes}{\rho_E}

This section proves the row package (\cref{prop:families:allrows}) together with three results of Part~I which
carry much of its weight and whose formal proofs belong to the same part of the development: the volume
recursion obtained from the Siegel--Weil formula (item (iv) of \cref{prop:families:volratio}), the automorphic
counting theorem with exponentially small error (\cref{thm:counting}), and the statements on period entries and
level valuations (\cref{prop:families:periods}).

\paragraph{What is proved.}
The row package $\sfname{AllRowsAlt}$ is a conjunction of twenty-six \emph{linked rows} (\cref{def:families:rows}).
Six of them are hypotheses of the theorem, passed through unchanged, and several others hold for every slot and
every linkage. The two largest inputs, the local rows coming from the Siegel--Weil formula (loop rows, Theorem~S and
the trace frame) and the service row coming from global Artin reciprocity, are proved in other sections of Part~II:
we use them in the form \lref{GlobalTheory.LocalRows.lg1_sLp_sO_hTF} and \lref{GlobalTheory.T.hServ_of_linkR}.
The remaining substance of the row package is proved here:
\begin{enumerate}[label=(\alph*)]
\item the counting row, which rests on the automorphic counting theorem (\cref{sec:pf:rows:counting,sec:pf:rows:countrow});
\item six rows on period entries, ties at the real places and valuations at the level primes, which hold at all
points lying over true column choices (\cref{sec:pf:rows:periods}); for the valuations at the level primes we give the proof
down to the stability of lattice column chains and the $p$-adic moments of CM theta functions (\cref{sec:pf:rows:level});
\item the definite rows at stage zero and the arithmetic row at the generator places (\cref{sec:pf:rows:rrows});
\item a number of smaller rows (\cref{sec:pf:rows:small}).
\end{enumerate}

\paragraph{The counting theorem.}
The mathematical centre of the section is \cref{thm:counting}. Let $I$ be a unitary group of rank $n$ over a
CM field $E$ with totally real subfield $F$, $f=[F:\Q]$, of signature $(p,q)$ at every real place, with a
self-dual lattice, and let $\sigma_\infty$ be a discrete series representation of $\prod_v\Uni(p,q)$ whose
Harish-Chandra parameters have a large gap. The number $\rowsDc$ to be counted is the dimension of the space of
intertwiners of $\sigma_\infty$ into $L^2(\Gamma\backslash \rowsGA)$, cut down by an idempotent $e$ of the Hecke algebra at
a level $g$ which is unipotent modulo an auxiliary prime $g_0$. The proof is a pre-trace formula argument in the
tradition of the limit multiplicity theorems for discrete series \cite{DeGeorgeWallach1978,Clozel1986,Savin1989},
but with explicit and uniform constants. The test function is the normalised matrix coefficient of the lowest
$K$-type vector tensored with $e$; the corresponding convolution operator is the orthogonal projection onto the
space to be counted, so that $\rowsDc$ is the integral of the diagonal of the automorphic kernel. The identity term
gives the main term $\rowsDst$, and the other terms are made exponentially small relative to $\rowsDst$ by three
mechanisms:
\begin{enumerate}[label=(\roman*)]
\item a large gap in the Harish-Chandra parameters makes the lowest-$K$-type matrix coefficient decay like
$\exp(-h\,d^2/8)$ in the Riemannian distance $d$ to the maximal compact subgroup (a result proved elsewhere in Part~II);
\item a rational element $\gamma\ne1$ which stays close to the maximal compact subgroup at all infinite places and is
in the support of the idempotent is unipotent (\cref{pf:rows:nearunip}): all traces of its powers are congruent to
$n$ modulo $g_0>4n$ and bounded by $3n$ at all archimedean places, hence equal to $n$;
\item reduction theory for the hermitian lattices $\Lambda_x$ attached to adelic points, by filtrations of the kind of the
canonical filtrations of Stuhler and Grayson \cite{Stuhler1976,Grayson1984}, shows that points carrying a destabilising
(necessarily isotropic) flag are rare: the volume of the set of points of a given flag type is computed by unfolding
along isotropic flags and bounded through the volume recursion (\cref{pf:rows:flagmeasure}), and turns out to be at
most $e^{-\frac13 nf\log n}$ times the total volume.
\end{enumerate}
Points very deep in the cusp, whose lattice has a sublattice of size below $\tau_*=(nfr_dgb)^{-11q}$, are treated
separately with a pointwise bound for cusp forms there (\cref{pf:rows:deepkernel}). The volume recursion itself
(\cref{pf:rows:volstep}) is deduced from the Siegel--Weil formula for one vector by unfolding a Gaussian theta series
over the orbit of a unit vector.

\paragraph{Organisation.}
\Cref{sec:pf:rows:framework} recalls rows and linked rows and proves the formal properties used later.
\Cref{sec:pf:rows:main} states the row package and proves it modulo the results of the later subsections.
\Cref{sec:pf:rows:volume} proves the volume recursion, \cref{sec:pf:rows:counting} the counting theorem and
\cref{sec:pf:rows:countrow} the counting row. \Cref{sec:pf:rows:periods} treats period entries and level
valuations, and \cref{sec:pf:rows:level} the valuations at the level primes; \cref{sec:pf:rows:rrows} treats the definite rows
and the arithmetic row at generator places, and \cref{sec:pf:rows:small} the remaining rows. \Cref{sec:pf:rows:aux} proves
auxiliary results on Weierstrass curves over ultrametric fields, singular $j$-invariants and approximation in unitary groups,
which are used for the integrality of the CM kernels. \Cref{sec:pf:rows:remarks} collects remarks on the formal statements.

% ---------------------------------------------------------------------------
\subsection{Rows and linked rows}\label{sec:pf:rows:framework}

We use the notions of \cref{def:families:closed,def:families:rows}, which we now make precise. Throughout, a
\emph{family} $\mathcal F\in\sfname{FamP}(\eps)$ has threshold $x_1=x_1(\mathcal F)$, index sets $I_x$, members
$(D_{x,i},L_{x,i})$ and, for $x>x_1$, objects $O_{x,i}$ over $D_{x,i}$ (size data, lattice instances, constants,
divisor data and the extra data $\mathrm{ex}(x,i)$). For a class $\mathcal A$ we write $\rowsWi(\mathcal A)$,
$\rowsWiii(\mathcal A)$ and $\rowsDiii(\mathcal A)=D3(3,\mathcal A)$ for the three standard subclasses of members of
$\mathcal A$ satisfying, respectively, the divisibility, tameness, box, row and cylinder conditions with smoothness
and good twisting; these together with a chosen exceptional set and the size conditions $\sfname{TuSize}$,
$\sfname{DGamSize}$; and the designed families of \cref{def:families:closed}. Every member of $\rowsDiii(\mathcal A)$ is a member of
$\rowsWiii(\mathcal A)$ (\lref{GlobalTheory.Fc.working_classes_frozen}).

\begin{definition}[Rows]\label{pf:rows:defrow}
Let $\mathrm{Bd}$ be a \emph{bundle type} (a type $\mathrm{Bd}(D,O,c)$ of bundles for every datum $D$, object $O$ over
$D$ and $c\in\R$), $\rowsLk$ a \emph{linkage} for $\mathrm{Bd}$ (a relation $\rowsLk(K,\beta,c,b)$ between a context $K$,
a bundle point $\beta$ of $K$, a real number $c$ and a bundle $b\in\mathrm{Bd}(\beta.D,\beta.O,c)$) and $Q$ a slot
(a predicate on cores; we write $\beta\models Q$ if the core of $\beta$ satisfies $Q$). A \emph{raw selection}
$\sigma$ over a class $\mathcal A$ chooses a bundle $\sigma(\mathcal F,x,i)$ over $(D_{x,i},O_{x,i})$ for every
$\mathcal F\in\mathcal A$, $x>x_1$ and $i\in I_x$.
\begin{enumerate}[label=(\alph*)]
\item A \emph{parameter} $\kappa$ consists of an exponent, a constant system and five real constants; it is
\emph{positive} if three of them are positive, it \emph{governs} a context $K$ if $K$ has its exponent and constants
and satisfies a goodness condition, and it is \emph{linked} to a family $\mathcal F$ if it is positive, standard, and
has the exponent and the constant lists of $\mathcal F$. For a property $P(\kappa,K)$, $\rowsEvP(P)$ means: for small
$\eps$ and staged large constants (the filter of \cref{def:families:closed} on parameters), every positive parameter
$\kappa$ admits $x_0$ such that $P(\kappa,K)$ holds for all contexts $K$ governed by $\kappa$ with scale $x\ge x_0$.
\item A property $E(\kappa,D,O,c,b)$ of bundles holds \emph{eventually at linked bundles} (for $Q$ and $\rowsLk$) if
$\rowsEvP\bigl(\kappa,K\mapsto \forall\beta\models Q\ \forall c\ \forall b:\ \rowsLk(K,\beta,c,b)\Rightarrow
E(\kappa,\beta.D,\beta.O,c,b)\bigr)$.
\item A bundle $b$ over $(D,O,c)$ is \emph{built} for $\kappa$ if $(D,O,b)=(\beta.D,\beta.O,b')$ for a bundle point $\beta$
of a context governed by $\kappa$ with $\beta\models Q$ and $\rowsLk(K,\beta,c,b')$.
\item A raw selection $\sigma$ over $\mathcal A$ is an \emph{$E$-selection} if every $\mathcal F\in\mathcal A$ has a parameter
$\kappa$ linked to $\mathcal F$ such that $\mathcal F$ satisfies the seven duties and every $\sigma(\mathcal F,x,i)$ is built for
$\kappa$ and satisfies $E(\kappa,\cdot)$.
\item For a property $H(\mathcal A,\sigma)$ of classes with raw selections, $\rowsRow(Q,\rowsLk,H)$ means: there is a
property $E$ which holds eventually at linked bundles such that $H(\mathcal A,\sigma)$ holds for every class $\mathcal A$ and
every $E$-selection $\sigma$ over $\mathcal A$.
\item $\rowsAtAll(R)(\mathcal A,\sigma)$ means $R(\sigma(\mathcal F,x,i))$ for all $\mathcal F\in\mathcal A$, $x>x_1$, $i\in I_x$; for a
property $R(D,O)$ of data with objects, $\rowsEvC(R)(\mathcal A)$ is
$\rowsClosed\bigl(\rowsWiii(\mathcal A),\ \mathcal F\mapsto \exists x_0\ \forall x\ge x_0,\,x>x_1\ \forall i\in I_x:\
R(D_{x,i},O_{x,i})\bigr)$; and $\rowsRowZ(Q,H)$ is the row for the trivial bundle type and the trivial linkage
with requirement $(\mathcal A,\sigma)\mapsto H(\mathcal A)$.
\end{enumerate}
\end{definition}

\begin{definition}[Linked rows]\label{pf:rows:deflinked}
Let $\mathrm{KP}$ be a predicate on bases with frames and $Q,R$ slots. $\rowsLBC(\mathrm{KP},Q,R)$ is the statement:
eventually in the scale and for large twisting data (the filter $\sfname{EvLgG}$ of \cref{def:families:twist}), for
every base $B$ with frame $f$ satisfying $\mathrm{KP}$, every stage tuple $X$ over $B$ and every core $c$, $Q$ implies
$R$ at $c$. For a property $r$ of slots,
\[
\rowsLg(\mathrm{KP},Q,Q',r)\ :\Longleftrightarrow\ \exists R:\ \rowsLBC(\mathrm{KP},Q,R)\wedge r(Q'\wedge R).
\]
For a parameter set $p$ (\cref{def:families:rows}) with slot system $p.s$ and bundle linkage $\rowsLk_p$,
$\rowsLinked(p,r)=\rowsLg\bigl(\rowsKPx,\rowsBuilt(p.s),\rowsGood(p.s),S\mapsto r(S,\rowsLk_p)\bigr)$, where
$\rowsKPx(B,f)$ says that the coefficient field of $B$ is exact for some field datum relative to $f$ and that the
pre-base of $B$ is pinned.
\end{definition}

We use the following formal properties proved elsewhere. The statement $\rowsLBC(\mathrm{KP},Q,\mathrm{True})$ holds
(\lref{GlobalTheory.Fc.largeBaseProperty_true}: take all exceptional sets empty and all thresholds $0$).
If $H(\mathcal A,\sigma)$ holds for all classes and raw selections, then $\rowsRow(Q,\rowsLk,H)$
(\lref{GlobalTheory.Fc.Row.of_sel}: take $E=\mathrm{True}$). A row for $H$ is a row for every $H'$ implied by
$H$ on selections (\lref{GlobalTheory.Fc.Row.imp}). Two rows for the same slot and linkage give a row for the
conjunction (\lref{GlobalTheory.Fc.Row.and}: take $E\wedge E'$; the filters are closed under finite intersections),
and the same holds for rows under $\rowsLg$ (\lref{GlobalTheory.Fc.RowL.and}, intersecting the auxiliary
slots). A row for a larger slot is a row for a smaller one
(\lref{GlobalTheory.Fc.Row.strengthen}). Rows can be transported along maps of bundle types compatible with the
linkages (\lref{GlobalTheory.Fc.Row.map}). A row for $\rowsGood(p.s)$ and $\rowsLk_p$ gives a linked row
(\lref{GlobalTheory.Fc.RL.LinkedAt.of_row}: take the auxiliary slot $\mathrm{True}$). Finally, a property which holds
for all families of a class holds closedly on the class (\lref{Formalism.MainNotions.Pk.Closed.of_class}).

\begin{lemma}[Calculus of linked rows]\label{pf:rows:calculus}
Let $\mathrm{KP}$ be a base predicate and $Q,Q'$ slots.
\begin{enumerate}[label=(\alph*)]
\item If $r(S)$ holds for every slot $S$, then $\rowsLg(\mathrm{KP},Q,Q',r)$.
\item If $\rowsLg(\mathrm{KP},Q,Q',r)$ and $r(S)\Rightarrow r'(S)$ for every slot $S$, then $\rowsLg(\mathrm{KP},Q,Q',r')$.
\item For a linkage $\rowsLk$ of the bundle type of \cref{def:families:rows}, if $\rowsLg(\mathrm{KP},Q,Q',S\mapsto
\rowsRow(S,\rowsLk,H))$ and $\rowsLg(\mathrm{KP},Q,Q',S\mapsto\rowsRow(S,\rowsLk,H'))$, then
$\rowsLg(\mathrm{KP},Q,Q',S\mapsto\rowsRow(S,\rowsLk,H\wedge H'))$.
\item If $\rowsRowZ(Q,H)$, then $\rowsRow(Q,\rowsLk,(\mathcal A,\sigma)\mapsto H(\mathcal A))$ for every bundle type and linkage.
\item \emph{(Transfer.)} Let $R(D,O)$ be a property of data with objects such that
$\rowsEvP\bigl(\kappa,K\mapsto\forall\beta\models Q:\ R(\beta.D,\beta.O)\bigr)$. Then $\rowsRowZ(Q,\rowsEvC(R))$.
\item If a bundle $b$ is linked to a bundle point $\beta$ by the bundle linkage of a parameter set $p$, then the record
$b.R$ of $b$ is linked to $\beta$ by the reciprocity linkage of $p$.
\end{enumerate}
\formallabel{GlobalTheory.Fc.AllAtMemberStage0.eventually_coreProperty_of_forall}
\formallabel{GlobalTheory.Fc.AllAtMemberStage0.eventually_coreProperty_mono}
\formallabel{GlobalTheory.Fc.AllAtMemberStage0.forall_linked_bundles_and_alt}
\formallabel{GlobalTheory.Fc.Row0.row}
\formallabel{GlobalTheory.Fc.validForConstructedPointClasses_eventuallyWorkingClass}
\formallabel{GlobalTheory.Fc.AllAtMemberStage0.linkR_of_isLinked'}
\end{lemma}

\begin{proof}
(a) Take the auxiliary slot $R=\mathrm{True}$, for which $\rowsLBC(\mathrm{KP},Q,R)$ holds. (b) Keep the auxiliary slot.
(c) This is \lref{GlobalTheory.Fc.RowL.and}. (d) Transport along the map sending every bundle to the unique
trivial bundle (\lref{GlobalTheory.Fc.Row.map}); the trivial linkage is implied by any linkage.
(e) Take $E(\kappa,D,O,c,b)=R(D,O)$; that $E$ holds eventually at linked bundles is the hypothesis. If $\sigma$ is an
$E$-selection over $\mathcal A$, then every member $\mathcal F$ of $\mathcal A$, and therefore every member of $\rowsWiii(\mathcal A)$,
satisfies $R(D_{x,i},O_{x,i})$ for all $x>x_1$ and $i\in I_x$; hence $\rowsEvC(R)(\mathcal A)$ holds with $\eps_1=1$,
an arbitrary threshold and $x_0=0$.
(f) Unpack the witnessed point $q$, the bundle $b'$ over it and the bundle witness $w$ given by the bundle linkage; then
the point $q$, the record $b'.R$ and the record pin contained in $w$ witness the reciprocity linkage, the required equality of triples $(D,O,\text{record})$
being the image of the given equality of triples $(D,O,\text{bundle})$ under $b\mapsto b.R$.
\end{proof}

% ---------------------------------------------------------------------------
\subsection{The row package}\label{sec:pf:rows:main}

\paragraph{The inputs.}
Besides the hypothesis $\sfname{RecTExists}$ (\cref{sec:galois}), which serves only to form the parameter sets and the
bundle linkage, the theorem has nine propositional hypotheses: global Artin reciprocity (class field theory in the
normalisation in which uniformisers go to geometric Frobenius elements \cite{ArtinTate1967,Tate1967gcft,Neukirch1999};
it is used only by the service row), the Siegel--Weil formula in the convergent range (\cref{lit:families:sw}), the five
statements of \cref{pf:rows:arch} below, and $\sfname{SurjectiveOntoX}$ and $\sfname{LambdaLowerBound}_2$
(\cref{prop:families:points}). In the application all of them are theorems of the development
(\cref{rem:families:supply}).

\begin{littheorem}[Real unitary groups and their automorphic forms {\cite{HarishChandra1966,HarishChandra1968,HechtSchmid1975,Casselman1989,Wallach1992,BorelJacquet1979}}]\label{pf:rows:arch}
Let $\Uni(\eps)=\{g\in\GL_n(\C):g^*J_\eps g=J_\eps\}$, $J_\eps=\mathrm{diag}(\eps_1,\dots,\eps_n)$, $\eps_a=\pm1$.
\begin{enumerate}[label=(\arabic*)]
\item \emph{(Discrete series cover.)} The discrete series of $\Uni(\eps)$ is exhausted, up to unitary equivalence, by
a family indexed by Harish-Chandra parameters, with prescribed $K$-types \cite{HarishChandra1966}.
\item \emph{(Harish-Chandra modules.)} The $K$-finite vectors of an irreducible unitary representation are smooth and
form an admissible irreducible $(\mathfrak g,K)$-module, and infinitesimally equivalent irreducible unitary
representations are unitarily equivalent \cite{HarishChandra1966,Wallach1992}.
\item \emph{(Rapid decay.)} Smooth cuspidal automorphic functions of uniformly moderate growth on unitary groups are
rapidly decreasing on Siegel domains \cite{HarishChandra1968,BorelJacquet1979}.
\item \emph{(Casselman--Wallach.)} An embedding of the Harish-Chandra module of an irreducible unitary representation
into a smooth Fr\'echet representation of moderate growth extends uniquely to a continuous equivariant map of smooth
vectors \cite{Casselman1989,Wallach1992}.
\item \emph{(Classification.)} For every strictly decreasing $\lambda\in(\frac{n-1}2+\Z)^n$ there is an irreducible
square-integrable representation with infinitesimal character $\lambda$ and a unique lowest $K$-type, the other
$K$-types differing from it by non-negative combinations of positive non-compact roots; every irreducible
square-integrable representation is of this form \cite{HarishChandra1966,HechtSchmid1975}.
\end{enumerate}
\end{littheorem}

The formal statements are specialisations of these results to the groups above; they are fields of the bundle $P$ of
\cref{sec:families:theta}.

\begin{theorem}[\cref{prop:families:allrows}]\label{pf:rows:main}
Assume global Artin reciprocity, the Siegel--Weil formula in the convergent range, the five statements of
\cref{pf:rows:arch}, $\sfname{SurjectiveOntoX}$ and $\sfname{LambdaLowerBound}_2$, and let
$h_L:\sfname{RecTExists}$. Let $N$ be a norm-choice extension, $S$ an $S$-extension, $\Gamma$ a geometric predicate,
$\mathcal N$ a predicate on family classes, and let $p=p(N,S,\Gamma)$ be the parameter set formed by the fixed row
position $R_X$ and $(N,S,\Gamma)$ (\cref{def:families:rows}). Assume
\begin{enumerate}[label=(\roman*)]
\item the twelve open parameter statements for $p$ (\cref{thm:families:open}): the linked
rows $\mathsf{hPA}$, $\mathsf{hLamP}$, $\mathsf{hDivG}$, $\mathsf{hDcPtG}$, $\mathsf{selJ}$ of \cref{tab:pf:rows:fields},
the comparison row $\mathsf{tie}'$, the label row $\mathsf{lab}$, the goodness row $\mathsf{good}$, the geometric row
$\mathsf{geo}$, the multiplicity row $\mathsf{mi}$, the $\kappa$-row $\mathsf{kap}$ and the defect row $\mathsf{df}$;
\item $\rowsLinked(p,\cdot)$ for the covolume row (eventually, for all $C\ge C_0$, the covolume bound $\sfname{L3}$ of the
realisation of the selected bundles), for the functional row (the same for the slope bound $\sfname{FUN}$, for members
satisfying $\sfname{DIVinN}$) and for the norm row of $\mathcal N$ (the requirement $\mathcal N(\rowsDiii(\mathcal A))$).
\end{enumerate}
Then $\sfname{AllRowsAlt}\bigl(\rowsKPx,\rowsBuilt(p.s),\rowsGood(p.s),\rowsLk_p,\mathcal N,\sfname{DefiniteSideCondition}\bigr)$:
each of the twenty-six requirements of \cref{tab:pf:rows:fields} holds as a row $\rowsLg(\rowsKPx,\rowsBuilt(p.s),
\rowsGood(p.s),S\mapsto\rowsRow(S,\rowsLk_p,H))$ (for $\mathsf{c30}$, as the statement of \cref{pf:rows:countrow}).
\formallabel{GlobalTheory.Fc.AllAtMemberStage0.allRows_of_open0_alt}
\end{theorem}

{\small
\begin{longtable}{>{\raggedright\arraybackslash}p{0.09\textwidth}>{\raggedright\arraybackslash}p{0.6\textwidth}>{\raggedright\arraybackslash}p{0.22\textwidth}}
\caption{The twenty-six rows of $\sfname{AllRowsAlt}$. Here $b=\sigma(\mathcal F,x,i)$, and ``eventually'' means: for
$x$ beyond some $x_0$ and all $i\in I_x$.}\label{tab:pf:rows:fields}\\
\toprule
field & requirement $H(\mathcal A,\sigma)$ & proved in\\
\midrule
\endfirsthead
\toprule
field & requirement & proved in\\
\midrule
\endhead
\bottomrule
\endfoot
$\mathsf{hn4}$ & $\rowsClosed(\rowsWi(\mathcal A))$: at all points, the class number formula for the words whose level prime
splits & \cref{pf:rows:small}(a)\\
$\mathsf{jOn}$ & $\rowsClosed(\rowsWiii(\mathcal A))$: eventually $\iota_{\mu_0}$ maps $V_0(\mu_0)$ onto the whole space & \cref{pf:rows:small}(d)\\
$\mathsf{hInstR}$ & at all points: volume--signature independence and the class items in the rank of $Y_0$
& \cref{pf:rows:small}(b)\\
$\mathsf{aC2}$ & $\rowsClosed(\rowsWiii(\mathcal A))$: the second-complement inputs & \cref{pf:rows:small}(c)\\
$\mathsf{aP10}$ & $\rowsEvC$: for every word $u$, the sharp valuation bound at the level prime and the
twisted column bound at $T_u$ & \cref{pf:rows:periods}\\
$\mathsf{aAtG}$ & $\rowsClosed(\rowsWiii(\mathcal A))$: for some $C_1$, eventually the used tuples are bounded by $C_1$ at the
generator places, and the size bound at the type prime & \cref{pf:rows:atg}\\
$\mathsf{aPt0}$ & $\rowsEvC$: the reference-chain evaluation statement for every word & \cref{pf:rows:periods}\\
$\mathsf{aPtTw}$ & $\rowsEvC$: integrality of the twisted period points & \cref{pf:rows:periods}\\
$\mathsf{aPer0}$ & $\rowsEvC$: the periods at $\mu_0$ are values of elements of $\mathit{LL}$ with bounded denominators
& \cref{pf:rows:periods}\\
$\mathsf{aPerTw}$ & $\rowsEvC$: the same for twisted periods & \cref{pf:rows:periods}\\
$\mathsf{tie}$ & $\rowsClosed(\rowsWiii(\mathcal A))$: the real-place second-layer conditions & \cref{pf:rows:periods}\\
$\mathsf{hRB}$ & $\rowsAtAll$: the realisation of $b$ has bounded rows & \cref{pf:rows:small}(e)\\
$\mathsf{hTw}$ & $\rowsAtAll$: the twisting records of $b$ are tied & \cref{pf:rows:small}(f)\\
$\mathsf{hPA}$, $\mathsf{hLamP}$, $\mathsf{hDivG}$, $\mathsf{hDcPtG}$, $\mathsf{selJ}$ & $\rowsAtAll$: primitive rows and separated
types; $\lambda$-rows; division rows; decomposition points; existence of a join & hypothesis (i)\\
$\mathsf{sLp}$, $\mathsf{sO}$, $\mathsf{hTF}$ & loop rows, Theorem~S and the trace frame for $\rowsDiii(\mathcal A)$ and the spectral
selection (\cref{thm:theorem-S}) & \lref{GlobalTheory.LocalRows.lg1_sLp_sO_hTF}\\
$\mathsf{hServ}$ & the selection serves $\rowsDiii(\mathcal A)$ & \lref{GlobalTheory.T.hServ_of_linkR}\\
$\mathsf{zL}$ & $\rowsClosed(\rowsDiii(\mathcal A))$: eventually the level projections $\pi_k$ are surjective &
\cref{pf:rows:small}(g)\\
$\mathsf{c30}$ & counting constants: the counting row and the irreducibility count row & \cref{pf:rows:countrow}\\
$\mathsf{rRows}$ & $\rowsClosed(\rowsDiii(\mathcal A))$: for some $C$, eventually, the constant bounds of the realisation and
$\sfname{DefiniteSideCondition}$ & \cref{pf:rows:rrows}\\
$\mathsf{nRest}$ & $\mathcal N(\rowsDiii(\mathcal A))$ & hypothesis (ii)\\
\end{longtable}
}

\begin{proof}
Write $Q=\rowsBuilt(p.s)$, $Q'=\rowsGood(p.s)$, $\rowsLk=\rowsLk_p$ and $o$ for the open statements. We establish the
fields of $\sfname{AllRowsAlt}$ one group at a time; ``for every slot'' means that the requirement is proved as
$\rowsRow(S,\rowsLk,H)$ for every slot $S$, so that it holds under $\rowsLg$ by \cref{pf:rows:calculus}(a).
\begin{enumerate}[label=\arabic*.]
\item \emph{Passed through.} $\mathsf{hPA}$, $\mathsf{hLamP}$, $\mathsf{hDivG}$, $\mathsf{hDcPtG}$ and $\mathsf{selJ}$ are the
corresponding fields of $o$, and $\mathsf{nRest}$ is the norm row of hypothesis (ii).
\item \emph{Rows for every slot.} $\mathsf{hn4}$ and $\mathsf{aC2}$ hold for every slot by \cref{pf:rows:small}(a),(c).
By \lref{GlobalTheory.VolRatioFromSiegelWeil.volumeRatio}, the Siegel--Weil formula implies the mass ratio
identity (the hermitian volume ratio of the datum, whenever its volumes are tied) for every datum $D$, and then
$\mathsf{hInstR}$ holds for every slot by \cref{pf:rows:small}(b). $\mathsf{hRB}$ holds for every slot by
\cref{pf:rows:small}(e), using (1)--(4) of \cref{pf:rows:arch}.
\item \emph{Period rows.} $\mathsf{aPt0}$, $\mathsf{aPtTw}$, $\mathsf{aPer0}$, $\mathsf{aPerTw}$, $\mathsf{tie}$ and $\mathsf{aP10}$ are
\cref{pf:rows:periods} applied to the slot system $p.s$ and the linkage $\rowsLk$. $\mathsf{aAtG}$ is \cref{pf:rows:atg},
applicable because the row position of $p=p(N,S,\Gamma)$ is the fixed position $R_X$ by definition.
\item \emph{Local rows.} By \lref{GlobalTheory.LocalRows.lg1_sLp_sO_hTF}, the Siegel--Weil formula implies, for
every base predicate, all slots and every linkage of bundles, the three rows $\mathsf{sLp}$, $\mathsf{sO}$ and $\mathsf{hTF}$
under $\rowsLg$; we apply it to $(\rowsKPx,Q,Q',\rowsLk)$.
\item \emph{Rows deduced from the open statements.} $\mathsf{jOn}$ and $\mathsf{zL}$ follow from $o.\mathsf{hPA}$ and
$\mathsf{hTw}$ from $o.\mathsf{tie}'$, by \cref{pf:rows:calculus}(b) and \cref{pf:rows:small}(d),(g),(f).
\item \emph{Counting row.} $\mathsf{c30}$ follows from $o.\mathsf{lab}$ and the Siegel--Weil formula by \cref{pf:rows:countrow}.
\item \emph{Service row.} By \cref{pf:rows:calculus}(f), every bundle linked to a point by $\rowsLk_p$ has a record linked
to the point by the reciprocity linkage. By \lref{GlobalTheory.T.hServ_of_linkR}, global Artin reciprocity
implies that for every parameter set $p$ and every linkage whose bundles have reciprocity-linked records, the
selections built from good points of $p$ serve their classes (the service row for the slot $\rowsGood(p.s)$); and
\lref{GlobalTheory.Fc.RL.LinkedAt.of_row} turns this row into the linked row $\mathsf{hServ}$.
\item \emph{Definite rows.} By \cref{pf:rows:calculus}(c) applied repeatedly, the nine linked rows $o.\mathsf{hPA}$,
$o.\mathsf{hDivG}$, $o.\mathsf{good}$, $o.\mathsf{geo}$, the covolume row, the functional row, $o.\mathsf{mi}$,
$o.\mathsf{kap}$ and $o.\mathsf{df}$ give a single linked row of their conjunction. For every slot at which this
conjunction holds, \cref{pf:rows:rrows} (which uses \cref{pf:rows:arch}, $\sfname{SurjectiveOntoX}$ and
$\sfname{LambdaLowerBound}_2$) gives the row $\mathsf{rRows}$ for the definite-side condition
$\sfname{DefiniteSideCondition}$; conclude by \cref{pf:rows:calculus}(b).\qedhere
\end{enumerate}
\end{proof}

% ---------------------------------------------------------------------------
\subsection{The volume recursion}\label{sec:pf:rows:volume}

\paragraph{Automorphic instances.}
The counting theorem and the volume recursion are statements about the following objects.

\begin{definition}[Automorphic instances]\label{pf:rows:instance}
An \emph{automorphic instance} $I$ consists of:
\begin{enumerate}[label=(\alph*)]
\item a totally real number field $F$, $f=[F:\Q]$, and a quadratic extension $E/F$ with an automorphism
$a\mapsto\bar a$ of $E$ whose fixed field is $F$, such that $E$ is totally complex, $\overline{\phi(a)}=\phi(\bar a)$
for every embedding $\phi:E\to\C$, and $E/F$ is unramified at all finite places (for every maximal ideal $\mathfrak P$
of $\OO_E$ there is $\theta\in\OO_E$ with $\theta-\bar\theta\notin\mathfrak P$); we write $d_E$ for the discriminant,
$r_d=|d_E|^{1/(2f)}$ for the root discriminant and $v_E=r_d/2$;
\item an integer $n\ge0$, a hermitian matrix $J\in M_n(E)$ with $\det J\ne0$, and a finitely generated
$\OO_E$-submodule $L\subseteq E^n$ spanning $E^n$ which is self-dual: $x\in L$ if and only if $y^*Jx\in\OO_E$ for all
$y\in L$;
\item for every real place $v$ of $F$, a vector of signs $\eps_v\in\{\pm1\}^n$, an embedding $\phi_v:E\to\C$ above $v$
and $P_v\in\GL_n(\C)$ with $P_v^*J^{\phi_v}P_v=J_{\eps_v}=\mathrm{diag}(\eps_v)$;
\item a group topology on $\rowsGf=\{g\in\GL_n(E\otimes_F\A_{F,\mathrm f}):g^*Jg=J\}$ for which the principal congruence
subgroups $K(g)=\{x\in K:(x-1)\hat L\subseteq g\hat L\}$, $g\ge1$, are open and form a basis of neighbourhoods of $1$; here
$\hat L$ is the closure of $L$ (its span over the integral adeles) and $K=K(1)$ is the stabiliser of $\hat L$.
\end{enumerate}
We put $\Uni(\eps)=\{g\in\GL_n(\C):g^*J_\eps g=J_\eps\}$, $\rowsGinf=\prod_v\Uni(\eps_v)$, $\rowsGA=\rowsGinf\times\rowsGf$, and let
$\Gamma\subseteq\rowsGA$ be the image of $\{\gamma\in\GL_n(E):\gamma^*J\gamma=J\}$ under
$\gamma\mapsto((P_v^{-1}\gamma^{\phi_v}P_v)_v,\gamma)$. The signature of $I$ at $v$ is $(p_v,q_v)$, the numbers of
$+1$ and $-1$ in $\eps_v$; $I$ has signature $(p,q)$ if $(p_v,q_v)=(p,q)$ for all $v$. On $\Uni(\eps_v)$ we fix the
Haar measure $\mu_v$ of the development; for a Haar measure $\mu_{\mathrm f}$ on $\rowsGf$ we put
$\mu_{\A}=\prod_v\mu_v\times\mu_{\mathrm f}$. A \emph{fundamental domain} is a measurable set $\mathcal F\subseteq\rowsGA$
with $\Gamma\mathcal F=\rowsGA$ and $\mu_{\A}(\gamma\mathcal F\cap\mathcal F)=0$ for $\gamma\in\Gamma\smallsetminus\{1\}$.
\end{definition}

For $x\in\rowsGA$ let $\rowsLv_x$ be the $\OO_E$-span of $\{y\in E^n: x_{\mathrm f}^{-1}y\in\hat L\}$, a lattice in $E^n$,
with the hermitian metrics $\|y\|_{x,v}=\|x_v^{-1}P_v^{-1}y^{\phi_v}\|$ (standard norm of $\C^n$) at the places of $E$
above $v$ (and their conjugates). Thus $\rowsLv_{\gamma x}=\gamma\rowsLv_x$ for $\gamma\in\Gamma$, and $x\mapsto(\rowsLv_x,\|\cdot\|_x)$
identifies $\Gamma\backslash\rowsGA/K$ with a set of isometry classes of hermitian $\OO_E$-lattices.

\paragraph{Gauge volumes and constants.}
For every $\eps$ there is a unique $c>0$ such that
$c\,\mu(\{g\in\Uni(\eps):\|g-1\|_{\mathrm{HS}}\le r\})\sim \vol(B^{n^2})\,r^{n^2}$ as $r\to0^+$
(\lref{GlobalTheory.GaugeConst.exists_unique_isGauge}); $c\mu$ is the Riemannian measure of the
Hilbert--Schmidt metric. Such constants $c=(c_v)_v$ are called \emph{gauge constants} of $I$. The \emph{gauge volume}
is
\[
\rowsgv(I)=\prod_vc_v\cdot\frac{\covol_\Gamma(\mu_{\A})}{\mu_{\mathrm f}(K)},\qquad
\covol_\Gamma(\mu_{\A})=\inf\{\mu_{\A}(D):D\text{ measurable},\ \Gamma D=\rowsGA\};
\]
it does not depend on $\mu_{\mathrm f}$, and $\covol_\Gamma(\mu_{\A})=\mu_{\A}(\mathcal F)$ for every fundamental domain
(\lref{GlobalTheory.T.leftQuotientCovolume_eq_fundamentalDomain_inter}). We use the explicit constants
\begin{gather*}
\rowsvU(i)=2^{i(i-1)/2}\prod_{l=1}^{i}\frac{2\pi^l}{(l-1)!},\qquad
L_\eta(M)=\prod_{\mathfrak P}e_M(\mathfrak P),\\
\rowsmr(E,M)=\frac{\bigl((M-1)!\,(v_E/\pi)^M\bigr)^f}{(2v_E)^{f/2}}L_\eta(M),\qquad
\rowsrf(E,M)=(2v_E)^{f(M-\frac12)}L_\eta(M),\\
\rowsmc(E,M)=\frac{(2v_E)^{f/2}}{v_E^{Mf}L_\eta(M)},
\end{gather*}
where $e_M(\mathfrak P)$ is the Euler factor of \cref{prop:families:mass}; thus $\rowsmr(E,M)$ is the right-hand side of
the Siegel mass ratio of \cref{prop:families:mass}.

\begin{definition}[Volume recursion]\label{pf:rows:defvolstep}
$\sfname{VolumeStep}$ is the statement: for every automorphic instance $I$ of signature $(P,Q)$ with $n\ge3$ and
every sign $s\in\{\pm1\}$ with $P\ge1$ if $s=1$ and $Q\ge1$ if $s=-1$, there are an instance $I'$, a field isomorphism
$e:E'\to E$ commuting with the conjugations, $x\in E^n$ and $A\in M_{n\times(n-1)}(E)$ such that $n'=n-1$,
$x^*Jx=s$, $x^*JA=0$, $A^*JA=J'^{\,e}$, $I'$ has signature $(P-1,Q)$ if $s=1$ and $(P,Q-1)$ if $s=-1$, and for all
gauge constants and Haar measures of $I$ and $I'$
\[
\frac{\rowsgv(I)}{\rowsvU(n)^f}=\frac{\rowsgv(I')}{\rowsvU(n-1)^f}\cdot\rowsmr(E,n).
\]
\end{definition}

\begin{proposition}[\cref{prop:families:volratio}(iv)]\label{pf:rows:volstep}
The Siegel--Weil formula in the convergent range (\cref{lit:families:sw}) implies $\sfname{VolumeStep}$.
\formallabel{GlobalTheory.GaussMean.volumeStep_of_siegelWeil}
\end{proposition}

The other items of \cref{prop:families:volratio} are proved elsewhere: the mass ratio identity (i) is
\lref{GlobalTheory.VolRatioFromSiegelWeil.volumeRatio}, the identities $\sfname{SWXS}$ (ii) are
\lref{GlobalTheory.St.swxs_of_siegelWeil}, and the non-generic bound (iii) is
\lref{GlobalTheory.rowNGP_of_siegelWeil}.

The proof of \cref{pf:rows:volstep} has four steps: the Siegel--Weil formula for one vector, its extension from
compactly supported to Gaussian test functions, the unfolding of the Gaussian theta series over the orbit of a unit
vector, and a comparison of the constants. For $s\in\{\pm1\}$ and $x\in\rowsGA$ put
\[
\Theta_s(x)=\sum_{y\in\rowsLv_x,\ y^*Jy=s}\ \prod_v\exp\bigl(-\|x_v^{-1}P_v^{-1}y^{\phi_v}\|^2\bigr)\in[0,\infty],
\]
and for a sign vector $\eps$ let $q_\eps(y)=\sum_a\eps_a|y_a|^2$ on $\C^n$ and
$g_\eps(t)=\frac{d}{dt}\int_{q_\eps(y)\le t}e^{-|y|^2}dy$ (the density of the Gaussian measure along $q_\eps$).

\paragraph{Step 1: one vector.}
Let $(\rowsSWone)$ be the following statement: for a CM extension $E/F$ and a non-degenerate hermitian $J$ of size $M$,
the quotient $U(J)(F)\backslash U(J)(\A_F)$ has finite volume, and if $M\ge3$, $\beta\in F^\times$, $\Phi_\infty$ is smooth and
compactly supported and $\Phi_{\mathrm f}$ is a pure Schwartz--Bruhat function on $(\A_E)^M$, the mean over
$U(J)(F)\backslash U(J)(\A_F)$ of $h\mapsto\sum_{x\in E^M,\ h_J(x,x)=\beta}\Phi(h^{-1}x)$ equals
$F_\infty(\beta)\prod_vF_v(\beta)$, with continuous densities $F_\infty,F_v$ of the push-forwards of $\Phi_\infty,\Phi_v$
under $x\mapsto h_J(x,x)$ and a convergent product, for self-dual measures. By
\lref{GlobalTheory.GaussMean.oneVectorCc_of_siegelWeil}, the Siegel--Weil formula (\cref{lit:families:sw}) implies
$(\rowsSWone)$: it is the case of one vector, $n=1$ and $T=(\beta)$.

\begin{lemma}[Siegel means with compactly supported test functions]\label{pf:rows:siegelcc}
Assume $(\rowsSWone)$. Let $I$ be an instance with $n\ge3$, $s\in\{\pm1\}$, $\mu_{\mathrm f}$ a Haar measure
and $\mathcal F$ a fundamental domain. Then $\mu_{\A}(\mathcal F)<\infty$, and for smooth compactly supported functions
$\varphi_v\ge0$ on $\C^n$ with densities $D_v$ along $q_{\eps_v}$ (that is, $D_v$ is continuous on $\R\smallsetminus\{0\}$
and $\int_{q_{\eps_v}(y)\in S}\varphi_v=\int_SD_v$ for compact $S\subseteq\R\smallsetminus\{0\}$),
\[
\int_{\mathcal F}\ \sum_{y\in\rowsLv_x,\ y^*Jy=s}\ \prod_v\varphi_v\bigl(x_v^{-1}P_v^{-1}y^{\phi_v}\bigr)\,d\mu_{\A}(x)
=\mu_{\A}(\mathcal F)\cdot\rowsmc(E,n)\cdot\prod_vD_v(s).
\]
\formallabel{GlobalTheory.OneVectorMean.siegelMeanCc_of_oneVectorCc}
\end{lemma}

\begin{proof}
The instance is transported to the adelic model of $\Uni(J)$ used in $(\rowsSWone)$
(\lref{GlobalTheory.T.exists_adelic_unitaryGroup_models_equiv} and \lref{GlobalTheory.SiegelWeilMean.exists_transport}:
there is a measure-preserving identification of $\rowsGA$, $\Gamma$ and $\mathcal F$ with the corresponding adelic objects);
the first assertion is the finiteness of the volume contained in $(\rowsSWone)$. If $\prod_vD_v(s)=0$, then some
$\varphi_v$ vanishes on the quadric $q_{\eps_v}=s$ (\lref{GlobalTheory.OneVectorMean.eq_zero_of_isDens}: a
non-negative continuous function whose density vanishes at $s\ne0$ vanishes on the quadric), and both sides vanish.
Otherwise apply $(\rowsSWone)$ with $M=n$, $\beta=s$, $\Phi_\infty=\prod_v\varphi_v$ composed with the
coordinates $y\mapsto x_v^{-1}P_v^{-1}y^{\phi_v}$, and $\Phi_{\mathrm f}$ the indicator function of $\hat L$, which is
Schwartz--Bruhat and the product of the local indicators, equal to the indicator of the standard lattice at almost all
places. The theta sum becomes the lattice sum on the left. The archimedean density at $s$ is $\prod_vD_v(s)$
(\lref{GlobalTheory.T.archimedean_density_product_testFunction_at_sign}, which identifies the density of the
product test function with the product of the densities). At the finite places, $E/F$ is unramified and $\hat L$ is
self-dual, so the local densities are computed by an Euler product, and with the self-dual normalisations of the
measures their product is $\rowsmc(E,n)$
(\lref{GlobalTheory.T.cm_euler_product_discr_unramified} and \lref{GlobalTheory.SiegelWeilMean.const_eq}).
\end{proof}

\begin{lemma}[Gaussian Siegel means]\label{pf:rows:siegelgauss}
Under the conclusion of \cref{pf:rows:siegelcc} (for all instances), for every instance $I$ with $n\ge3$, every
$s\in\{\pm1\}$, every Haar measure $\mu_{\mathrm f}$ and fundamental domain $\mathcal F$, $\mu_{\A}(\mathcal F)<\infty$ and
\[
\int_{\mathcal F}\Theta_s\,d\mu_{\A}=\mu_{\A}(\mathcal F)\cdot\rowsmc(E,n)\cdot\prod_vg_{\eps_v}(s).
\]
\formallabel{GlobalTheory.GaussMean.Limit.siegelMean_of_cc}
\end{lemma}

\begin{proof}
For every $\eps$ there are smooth compactly supported $\varphi_k\ge0$ on $\C^n$, increasing in $k$ with limit $e^{-|y|^2}$
pointwise, with densities $D_k$ along $q_\eps$ such that $D_k(t)\to g_\eps(t)$ for every $t\ne0$
(\lref{GlobalTheory.T.exists_truncated_gaussians_with_densities}). Apply \cref{pf:rows:siegelcc} to
$(\varphi_{v,k})_v$. The left-hand sides are integrals of countable sums of non-negative terms which increase with $k$
to the terms of $\Theta_s$; by monotone convergence they converge to $\int_{\mathcal F}\Theta_s$. The right-hand sides
converge to $\mu_{\A}(\mathcal F)\rowsmc(E,n)\prod_vg_{\eps_v}(s)$, because $\mu_{\A}(\mathcal F)<\infty$.
\end{proof}

\begin{lemma}[Orthogonal complements of unit vectors]\label{pf:rows:perp}
Let $I$ have signature $(P,Q)$, $n\ge3$, and let $s\in\{\pm1\}$ with $P\ge1$ if $s=1$, $Q\ge1$ if $s=-1$. There are an
instance $I'$, an isomorphism $e:E'\to E$ commuting with the conjugations, $x\in E^n$, $A\in M_{n\times(n-1)}(E)$ and
$x_{\mathrm f}\in\rowsGf$ such that $n'=n-1$, $x^*Jx=s$, $x^*JA=0$, $A^*JA=J'^{\,e}$, $x\in\rowsLv_{(1,x_{\mathrm f})}$,
$y\in L'$ if and only if $Ay^e\in\rowsLv_{(1,x_{\mathrm f})}$, and $I'$ has signature $(P-1,Q)$ resp.\ $(P,Q-1)$.
\end{lemma}

This is \lref{GlobalTheory.T.exists_unit_vector_orthogonal_instance_rank_three}; its proof constructs $I'$ from
the signs and the unimodular local structure of $L$ (\lref{GlobalTheory.PerpInst.exists_isPerpInst}), and we do
not reproduce it. We call such data a \emph{perpendicular instance} for $(I,s)$.

\begin{lemma}[Unfolding the Gaussian theta series]\label{pf:rows:unfold}
Let $(I',e,x,A,x_{\mathrm f})$ be a perpendicular instance for $(I,s)$ with $n\ge3$, let $c,c'$ be gauge constants of
$I,I'$, $\mu_{\mathrm f},\mu'_{\mathrm f}$ Haar measures and $\mathcal F$ a fundamental domain for $I$. Then
\[
\prod_vc_v\cdot\frac{\int_{\mathcal F}\Theta_s\,d\mu_{\A}}{\mu_{\mathrm f}(K)}=\rowsgv(I')\cdot 2^{nf}\prod_vg_{\eps_v}(s).
\]
\formallabel{GlobalTheory.T.gaussianSum_stabiliserUnfolding}
\formallabel{GlobalTheory.StabUnfold.gaugeUnfold}
\end{lemma}

\begin{proof}[Proof (structure)]
The formal proof (\lref{GlobalTheory.StabUnfold.unfold_sib}) is the classical unfolding. (1) By Witt's theorem
$\Gamma$ acts transitively on the vectors $y\in E^n$ with $y^*Jy=s$, so
$\Theta_s(g)=\sum_i\Phi(r_ig)$ for representatives $r_i$ of $\Gamma_x\backslash\Gamma$, where
$\Phi(g)=\prod_v\exp(-\|g_v^{-1}P_v^{-1}x^{\phi_v}\|^2)\cdot\one[x\in\rowsLv_{(1,g_{\mathrm f})}]$
(\lref{GlobalTheory.T.latticeGaussianSumOfLength_eq_tsum_stabilizer_cosets}); hence
$\int_{\mathcal F}\Theta_s=\covol_{\Gamma_x}(\Phi\,\mu_{\A})$
(\lref{GlobalTheory.T.lintegral_tsum_cosets_fundamentalDomain_eq_covol}). (2) The stabiliser of $x$ in $\rowsGA$
is the image of the adelic unitary group of $I'$ under a closed embedding $j$ which maps $\Gamma'$ onto $\Gamma_x$
(\lref{GlobalTheory.T.image_rationalOrthComplementUnitary_eq_rationalStabilizer}) and the congruence subgroups
of $I'$ into those of $I$ (\lref{GlobalTheory.T.mem_congruence_level_iff_orthEmbedding}). (3) For a left
$j$-invariant density $\Phi$ and a cut-off $b$ with $\int b(j(h)g)\,dh=1$ for all $g$,
$\covol_{j(\Gamma')}(\Phi\mu)=\covol_{\Gamma'}(\mu')\int\Phi b\,d\mu$
(\lref{GlobalTheory.T.leftQuotientCovolume_withDensity_unfold_subgroup}). (4) At the finite places the stabiliser
of $x$ acts transitively on $\{g: x\in\rowsLv_{(1,g)}\}$ modulo $K$
(\lref{GlobalTheory.T.stabilizer_transitive_adelic_unit_length_vectors}), which allows a cut-off with
$\int\one[x\in\rowsLv_{(1,g)}]b_{\mathrm f}(g)\,d\mu_{\mathrm f}=\mu_{\mathrm f}(K)/\mu'_{\mathrm f}(K')$
(\lref{GlobalTheory.T.exists_massOne_cutoff_doubleCoset_integral}). (5) At a real place, the orbit of $x_v$ under
$\Uni(\eps_v)$ is the quadric $q_{\eps_v}=s$, and with gauge-normalised measures the integral of the Gaussian against the
cut-off equals $(c'_v/c_v)\,2^n g_{\eps_v}(s)$ (\lref{GlobalTheory.StabUnfold.real_quot}). Multiplying over all
places gives the formula.
\end{proof}

\begin{lemma}[Constants]\label{pf:rows:constants}
For $M\ge1$: $\rowsrf(E,M)\,\rowsmc(E,M)=2^{Mf}$ and $\rowsmr(E,M)\,\rowsvU(M)^f=\rowsrf(E,M)\,\rowsvU(M-1)^f$.
\formallabel{GlobalTheory.GaussMean.Step.rankFac_mul_meanConst}
\formallabel{GlobalTheory.T.mass_ratio_rank_factor}
\formallabel{GlobalTheory.FlagMeas.mass_ratio_mul_unitary_volume}
\end{lemma}

\begin{proof}
$\rowsrf\cdot\rowsmc=(2v_E)^{f(M-\frac12)+\frac f2}v_E^{-Mf}=2^{Mf}$. Since
$\rowsvU(M)/\rowsvU(M-1)=2^{M-1}\cdot2\pi^M/(M-1)!$,
$\rowsmr(E,M)\rowsvU(M)^f=\rowsvU(M-1)^f\,(2v_E)^{-f/2}L_\eta(M)\,(2^Mv_E^M)^f=\rowsrf(E,M)\rowsvU(M-1)^f$.
\end{proof}

\begin{lemma}[From Gaussian means to the volume recursion]\label{pf:rows:gausstostep}
If the conclusion of \cref{pf:rows:siegelgauss} holds for all instances, then $\sfname{VolumeStep}$ holds.
\formallabel{GlobalTheory.GaussMean.Step.volumeStep_of_siegelMean}
\end{lemma}

\begin{proof}
Let $I,P,Q,s$ be as in \cref{pf:rows:defvolstep}, and take a perpendicular instance $(I',e,x,A,x_{\mathrm f})$
(\cref{pf:rows:perp}); it has all the properties required in $\sfname{VolumeStep}$. Let $c,c',\mu_{\mathrm f},\mu'_{\mathrm f}$
be given and choose a fundamental domain $\mathcal F$
(\lref{GlobalTheory.T.automorphicInstance_exists_fundamental_domain}). Put $d=\prod_vg_{\eps_v}(s)$; then $d>0$,
because the sign $s$ occurs in every $\eps_v$ and $g_\eps(s)>0$ in that case
(\lref{GlobalTheory.T.gaussianQuadricDensity_exists_pos_continuous}); also $L_\eta(n)>0$
(\lref{GlobalTheory.T.neg_log_prod_lValEta_le_and_pos}). Combining the hypothesis with \cref{pf:rows:unfold} and
$\covol_\Gamma(\mu_{\A})=\mu_{\A}(\mathcal F)$,
\[
\rowsgv(I)\cdot\rowsmc(E,n)\,d=\prod_vc_v\frac{\mu_{\A}(\mathcal F)}{\mu_{\mathrm f}(K)}\,\rowsmc(E,n)\,d
=\rowsgv(I')\cdot2^{nf}d .
\]
By \cref{pf:rows:constants}, $2^{nf}=\rowsrf(E,n)\rowsmc(E,n)$, so cancelling $\rowsmc(E,n)d\in(0,\infty)$ gives
$\rowsgv(I)=\rowsgv(I')\rowsrf(E,n)$ (in $[0,\infty]$), and the second identity of \cref{pf:rows:constants} turns this
into the formula of $\sfname{VolumeStep}$ (\lref{GlobalTheory.GaussMean.Step.ennreal_volume_ratio_step_cancel}).
\end{proof}

\begin{proof}[Proof of \cref{pf:rows:volstep}]
Step 1, \cref{pf:rows:siegelcc}, \cref{pf:rows:siegelgauss} and \cref{pf:rows:gausstostep}, applied in this order.
\end{proof}

% ---------------------------------------------------------------------------
\subsection{Automorphic counting}\label{sec:pf:rows:counting}

\paragraph{Counting data.}
Let $I$ be an automorphic instance (\cref{pf:rows:instance}). \emph{Counting data} $B$ for $I$ consist of a Haar measure
$\mu_{\mathrm f}$ on $\rowsGf$ and a fundamental domain $\mathcal F$; for every real place $v$ a vector $\lambda_v\in\R^n$, a
Hilbert space $H_v$ with a representation $\sigma_v$ of $\Uni(\eps_v)$ and a vector $w_v\in H_v$; a Hilbert space $H_s$
with a representation $\sigma_s$ of $\rowsGinf$ and a map $\otimes:\prod_vH_v\to H_s$ realising $H_s$ as the completed
tensor product ($\langle\otimes x,\otimes y\rangle=\prod_v\langle x_v,y_v\rangle$, $\sigma_s(g)(\otimes x)=\otimes(\sigma_v(g_v)x_v)$,
dense span); the right regular representation $R$ of $\rowsGA$ on $L^2(\mathcal F,\mu_{\A})=L^2(\Gamma\backslash\rowsGA)$;
integers $g,g_0$; an \emph{admissible idempotent} $e:\rowsGf\to\C$ of level $g$ and type $g_0$, that is, $e$ is locally
constant, $e*e=e$ for convolution with respect to $\mu_{\mathrm f}$, $e(x^{-1})=\overline{e(x)}$, $e$ is of positive type,
$e$ is supported in $K$, $e$ is bi-invariant under $K(g)$, and $e(x)\ne0$ implies $(x-1)^n\hat L\subseteq g_0\hat L$; and the
operator $R_e$ with $\langle F,R_eF'\rangle=\int e(y)\langle F,R(1,y)F'\rangle\,d\mu_{\mathrm f}(y)$. The \emph{count} is
\[
\rowsDc=\dim_\C\bigl\{R_e\circ T:\ T:H_s\to L^2(\Gamma\backslash\rowsGA)\text{ bounded},\ T\sigma_s(g)=R(g,1)T\ (g\in\rowsGinf)\bigr\}.
\]
For $d=(d_v)_v$, the data are \emph{discrete series data} if for every $v$: $\lambda_v$ is a Harish-Chandra parameter
($\lambda_v$ strictly decreasing, $\lambda_v(a)-\frac{n-1}2\in\Z$), $\sigma_v$ is a discrete series representation of
$\Uni(\eps_v)$ (non-zero, unitary, irreducible, square-integrable, with a $K$-type of highest weight
$\lambda_v+\rho_n-\rho_c$ which minimises the norm among its $K$-types), $w_v$ is a unit vector of that lowest $K$-type,
and $d_v>0$ is the formal degree with respect to $\mu_v$:
$\int\langle x,\sigma_v(g)y\rangle\overline{\langle x',\sigma_v(g)y'\rangle}\,d\mu_v=\langle x',x\rangle\langle y,y'\rangle/d_v$.
The \emph{main term} and the \emph{test function} are
\[
\rowsDst=\Bigl(\prod_vd_v\Bigr)\,\mu_{\A}(\mathcal F)\,\Re e(1),\qquad
\varphi(x)=\Bigl(\prod_vd_v\Bigr)\overline{\prod_v\langle w_v,\sigma_v(x_v)w_v\rangle}\;e(x_{\mathrm f})\quad(x\in\rowsGA),
\]
and the automorphic kernel is $K_\varphi(x,y)=\sum_{\gamma\in\Gamma}\varphi(x^{-1}\gamma y)$.

\begin{definition}[Counting hypotheses]\label{pf:rows:counthyp}
Let $C_g\in\R$, $n_0\in\N$, $p,q\in\N$, $b\in\R$. The data satisfy the \emph{counting hypotheses}
$\rowsCH(C_g,n_0,d,p,q,b)$ if they are discrete
series data, $I$ has signature $(p,q)$, and
\begin{gather*}
n\ge n_0,\quad q^2\le n,\quad \log r_d\le e^n,\quad \log f\le\sqrt n,\quad b\ge2,\quad \log b\le e^{\frac1{16}nf\log n},\\
|\lambda_v(a)|\le b^2,\quad 1\le g\le b^{C_g},\quad g_0\ \text{prime},\quad 4n<g_0,\quad g_0\ \text{splits completely in }E,\\
|\lambda_v(a)-\lambda_v(a')|\ge h\ \text{whenever}\ \eps_v(a)\ne\eps_v(a'),
\end{gather*}
for all $v,a,a'$, where $h=8\cdot10^{14}(C_g+4)\,n^{19}f^5\log(b\,r_d)$.
\end{definition}

\begin{theorem}[\cref{thm:counting}]\label{pf:rows:counting}
Assume $\sfname{VolumeStep}$. For every $C_g\in\R$ there is $n_0$ such that for all counting data satisfying
$\rowsCH(C_g,n_0,d,p,q,b)$ the space of the count is finite-dimensional and
\[
\bigl|\rowsDc-\rowsDst\bigr|\le e^{-\frac1{16}nf\log n}\,\rowsDst .
\]
\formallabel{GlobalTheory.CountFinal.countingTheorem_of_volumeStep}
\formallabel{GlobalTheory.T.unitaryAutomorphicCounting_of_compactDualVolumeRatio}
\end{theorem}

\paragraph{Analytic inputs from other sections.}
We use the following results proved elsewhere in Part~II. Write $\rowsdHS(g)=\bigl(\sum_j(\frac12\log s_j)^2\bigr)^{1/2}$ for
the eigenvalues $s_j$ of $g^*g$ (the distance to the maximal compact subgroup) and $\rowsdOp(g)=\log\|g\|$.
\begin{enumerate}[label=(S\arabic*)]
\item\label{pf:rows:S1} \emph{Decay.} If $\sigma$ is a discrete series representation of $\Uni(\eps)$ with parameter $\lambda$ of
gap $h\ge120\,n\,r^{3/2}$, $r=\min(p,q)$, then for $u,u'$ in the lowest $K$-type,
$|\langle u',\sigma(g)u\rangle|\le\|u\|\|u'\|\exp\bigl(-\frac h8\min(\rowsdHS(g)^2,\rowsdHS(g)/\sqrt{2r})\bigr)$ and the
same with $\min(2\rowsdOp(g)^2,\rowsdOp(g)/\sqrt r)$ (\lref{GlobalTheory.DecGapBound.decGap_of_formulaGap}, from the
compression formula for the lowest $K$-type \lref{GlobalTheory.Jet.FormGapFree.formulaGap}).
\item\label{pf:rows:S2} \emph{Pointwise bound} (\lref{GlobalTheory.Deep.ptBound_of_decGap}). For discrete series data
of signature $(p,q)$ with gap $h\ge120\,n\min(p,q)^{3/2}$: $|\varphi(y)|\le(\prod d_v)\Re e(1)$, and if $q\ge1$,
$|\varphi(y)|\le(\prod d_v)\Re e(1)\prod_v\exp\bigl(-\frac h8\min(\rowsdHS(y_v)^2,\rowsdHS(y_v)/\sqrt{2q})\bigr)$; moreover
$\varphi(1)=(\prod d_v)e(1)$, $e(1)$ is real and $\ge0$, and $\varphi(y)\ne0$ implies $y_{\mathrm f}\in K$.
\item\label{pf:rows:S3} \emph{Projection} (\lref{GlobalTheory.GlobProj.globalProj_holds}). If the archimedean part of
$\varphi$ is integrable, there is an operator $T$ with $\langle F,TF'\rangle=\int\varphi(y)\langle F,R(y)F'\rangle\,dy$; every such
$T$ is the orthogonal projection onto its range, which consists of the vectors $A(\otimes w)$ for the operators $A$ of the
count, and $\dim T=\rowsDc$.
\item\label{pf:rows:S4} \emph{Kernel} (\lref{GlobalTheory.T.kernelB'}, \lref{GlobalTheory.T.automorphicKernel_convergence_properties}).
Under the gap condition and if the space of the count is finite-dimensional, there are continuous $\Gamma$-invariant
functions $F_1,\dots,F_{\rowsDc}$ representing an orthonormal basis of the range of $T$ with
$K_\varphi(x,z)=\sum_iF_i(x)\overline{F_i(z)}$; the function $x\mapsto K_\varphi(x,x)$ is integrable on $\mathcal F$ and
$\int_{\mathcal F}K_\varphi(x,x)\,d\mu_{\A}=\rowsDc$; the series defining $K_\varphi$ converges absolutely.
\item\label{pf:rows:S5} \emph{Far terms} (\lref{GlobalTheory.FarTerms.farBound_of_ptBound}). See \cref{pf:rows:far}.
\end{enumerate}

\paragraph{Reduction theory.}
For $x\in\rowsGA$ and a non-zero $\OO_E$-submodule $R\subseteq\rowsLv_x$ of rank $r$ put
$\rowssize_x(R)=\bigl(2^{fr}|d_E|^{-r/2}\covol_x(R)\bigr)^{1/(2fr)}$, where $\covol_x(R)$ is the covolume of $R$ in its
real span for the euclidean structure given by the metrics $\|\cdot\|_{x,v}$; for saturated $R\subsetneq S$ the size
$\rowssizeQ_x(S/R)$ is defined in the same way with the covolume of $S/R$ and the rank of $S/R$. (The normalisation makes
the standard lattice $\OO_E^r$ have size $1$.) Let
\[
\rowsleast(x)=\inf_{0\ne R\subseteq\rowsLv_x}\rowssize_x(R),\qquad \tau_1(x)=\min\bigl(\tfrac12,\rowsleast(x)\bigr).
\]
For $t>0$ let $S_{\le t}(x)$ (resp.\ $S_{<t}(x)$) be the sum of all non-zero submodules $N\subseteq\rowsLv_x$ all of whose
quotients $N/S$ by saturated submodules $S\subsetneq N$ have size $\le t$ (resp.\ $<t$). They form a filtration analogous to
the canonical filtration of a hermitian lattice in the sense of Stuhler and Grayson \cite{Stuhler1976,Grayson1984}: they
are saturated,
$S_{<t}\subseteq S_{\le t}\subseteq S_{<t'}$ for $t<t'$, and $S_{<t}(x)=0$ if and only if $t\le\rowsleast(x)$
(\lref{GlobalTheory.T.canonicalLatticeFiltration_size_properties}); $\rowsleast$ is continuous and $\Gamma$-invariant
(\lref{GlobalTheory.LeastVar.least_variation_continuity}). Moreover, the member $S_{<1/2}(x)$ is totally isotropic
for the hermitian form, and its rank is at most $\min(p_v,q_v)$ for every $v$
(\lref{GlobalTheory.T.canonicalLatticeFlag_isotropic_selfDual}). For $a\ge1$, $j\in\Z_{\ge1}^a$ and $\tau$ in
$\Delta_a=\{\tau\in\R^a:0<\tau_1<\dots<\tau_a<\frac12\}$, the point $x$ has \emph{type} $(j,\tau)$ if the distinct
members of $\{S_{\le t}(x):t<\frac12\}$ form a flag $0=S_0\subsetneq S_1\subsetneq\dots\subsetneq S_a=S_{<1/2}(x)$
with $S_i=S_{\le\tau_i}(x)$, $\rowssizeQ_x(S_i/S_{i-1})=\tau_i$ and $\rank S_i-\rank S_{i-1}=j_i$. For a measurable
$B\subseteq\Delta_a$ let $\mathcal X(j,B)$ be the set of points of type $(j,\tau)$ for some $\tau\in B$.

Fix counting data satisfying $\rowsCH(C_g,n_0,d,p,q,b)$ and put
\begin{gather*}
X=n\,f\,r_d\,g\,b,\qquad \tau_*=X^{-11q},\qquad \eps_N=\frac{10^{-6}}{n^8f^2},\\
\rowsmod=\{x\in\mathcal F:\rowsleast(x)\ge\tau_*\},\qquad \rowsdeep=\mathcal F\smallsetminus\rowsmod .
\end{gather*}
An element $\gamma\in\Gamma$ is \emph{near} $x\in\rowsGA$ if $\varphi(x^{-1}\gamma x)\ne0$ and
$\rowsdHS((x^{-1}\gamma x)_v)\le\eps_N$ for all $v$; otherwise it is \emph{far} from $x$. Put
$\mathrm{nr}(x)=\sum_{\gamma\ne1\text{ near }x}|\varphi(x^{-1}\gamma x)|$ and
$\mathrm{fr}(x)=\sum_{\gamma\text{ far from }x}|\varphi(x^{-1}\gamma x)|$ (in $[0,\infty]$).

\paragraph{Near elements are unipotent.}

\begin{lemma}\label{pf:rows:nearunip}
Let $0<c\le10^{-5}$, $\eps=c/(n^8f^2)$ and $\delta=3\sqrt{n\eps}$. Assume $n\ge1$, $4n<g_0$ and that $g_0\OO_E$ is a
radical ideal. Let $x\in\rowsGA$ and $\gamma\in\Gamma$ with $\varphi(x^{-1}\gamma x)\ne0$ and $\rowsdHS((x^{-1}\gamma x)_v)\le\eps$
for all $v$. Then $\gamma$ is unipotent: $(\gamma-1)^n=0$ in $M_n(E)$ and $(\gamma_v-1)^n=0$ for all $v$; moreover
$(x^{-1}\gamma x)_{\mathrm f}\in K$ and $\|(x^{-1}\gamma x)_v-1\|_{\mathrm{HS}}\le\delta$ for all $v$. Finally $\delta<\frac14$ and
$\delta qnf\le\frac1{100}$ whenever $q^2\le n$.
\formallabel{GlobalTheory.NearElement.near_unipotent}
\end{lemma}

\begin{proof}
The numerical assertions follow from $n\eps\le c/n^4\le10^{-5}$ (\lref{GlobalTheory.NearElement.radius_bounds}).
Let $G\in\GL_n(E)$ be the rational matrix of $\gamma$ and $y=x^{-1}\gamma x$. Since $\varphi(y)\ne0$ we have $e(y_{\mathrm f})\ne0$,
so $y_{\mathrm f}\in K$ and $(y_{\mathrm f}-1)^n\hat L\subseteq g_0\hat L$; equivalently $G\rowsLv_x=\rowsLv_x$ and
$(G-1)^n\rowsLv_x\subseteq g_0\rowsLv_x$ (\lref{GlobalTheory.PointLattice.lamX_unipMod}).

\emph{Congruence.} For $j\ge0$, $\tr(G^j)-n\in g_0\OO_E$ (\lref{GlobalTheory.NearElement.trace_pow_cong}). Indeed
$\tr(G^j)\in\OO_E$ because $G$ preserves the lattice $\rowsLv_x$; for every prime $\mathfrak p\mid g_0$, $G-1$ induces a
nilpotent endomorphism of the $\OO_E/\mathfrak p$-vector space $\rowsLv_x/\mathfrak p\rowsLv_x$, so $G^j-1$ does too and
$\tr(G^j)\equiv n\pmod{\mathfrak p}$; as $g_0\OO_E$ is the intersection of the primes above $g_0$, the claim follows.

\emph{Archimedean bound.} Fix $v$ and write $y_v=k\exp(X_v)$ (polar decomposition, $k$ unitary, $X_v$ hermitian); then
$\|X_v\|_{\mathrm{HS}}=\rowsdHS(y_v)\le\eps$ (\lref{GlobalTheory.Count.Cartan.cartan_distances}). In the Hilbert--Schmidt
norm, $\|\exp(X_v)-1\|\le e^{\eps}-1$ and $\|y_v^j\|\le\sqrt n\,e^{j\eps}$; since $n\eps\le\frac14$, for $j\le n$ we get
$|\tr(y_v^j)|\le\sqrt n\|y_v^j\|\le ne^{1/4}\le\frac32n$ and $|\tr(y_v^j)-n|\le3n$
(\lref{GlobalTheory.NearUnitaryExp.near_arch}). Now
$\tr(y_v^j)=\phi_v(\tr G^j)$, because $y_v$ is conjugate to $P_v^{-1}G^{\phi_v}P_v$. Every complex embedding of $E$ is
$\phi_v$ or its conjugate for some $v$ (\lref{GlobalTheory.NearElement.embedding_eq_or_conj}), so the algebraic
integer $\tr(G^j)-n\in g_0\OO_E$ has all conjugates of absolute value at most $3n<g_0$. Its norm is then an integer
divisible by $g_0^{[E:\Q]}$ of absolute value $<g_0^{[E:\Q]}$, hence zero
(\lref{GlobalTheory.ringOfIntegers_eq_zero_of_natCast_dvd_of_norm_embedding_le}): $\tr(G^j)=n$ for $j\le n$.

\emph{Unipotence.} By Newton's identities, the elementary symmetric functions of the eigenvalues of $G$ are those of
$(1,\dots,1)$, so the characteristic polynomial of $G$ is $(t-1)^n$ and $(G-1)^n=0$
(\lref{GlobalTheory.unipotent_of_trace_pow_eq_size}). The same identity holds for the images of $G$ under all ring
homomorphisms, in particular $(\gamma_v-1)^n=0$. Finally, if $\tr(y_v)=n$, then $\|\exp X_v-1\|\le2\eps$ (as $\eps\le1$),
$\Re\tr k=n-\Re\tr(k(\exp X_v-1))\ge n-\sqrt n\|\exp X_v-1\|$ and $\|k-1\|_{\mathrm{HS}}^2=2n-2\Re\tr k\le4\sqrt n\,\eps\le4n\eps$
for the unitary $k$, so $\|y_v-1\|\le\|k-1\|+\|k(\exp X_v-1)\|\le2\sqrt{n\eps}+2\eps\le3\sqrt{n\eps}$.
\end{proof}

\paragraph{The definite case.}

\begin{lemma}\label{pf:rows:definite}
Let discrete series data have signature $(p,0)$ with $p\ge1$, $g\ge1$, $g_0$ a prime split completely in $E$ with
$4n<g_0$, and assume that the space of the count is finite-dimensional. Then $\rowsDc=\rowsDst$.
\formallabel{GlobalTheory.Dec.definite_count_of_finite}
\formallabel{GlobalTheory.T.definiteUnitaryGroup_count_eq_mainTerm}
\end{lemma}

\begin{proof}
Here every $\Uni(\eps_v)$ is the compact group $\Uni(n)$, so $\rowsdHS(g)=0$ for all $g\in\rowsGinf$. Since $g_0$ is unramified
in $E$, $g_0\OO_E$ is radical (\lref{GlobalTheory.T.isRadical_span_prime_of_ncard_primesOver_eq_finrank}). By
\cref{pf:rows:nearunip} every $\gamma\in\Gamma$ with $\varphi(x^{-1}\gamma x)\ne0$ is unipotent at a real place; a unipotent element of
the compact group $\Uni(n)$ is $1$ (\lref{GlobalTheory.unipotent_isometry_definite_eq_one}), and an element of $\Gamma$ is
determined by one archimedean component. Hence $K_\varphi(x,x)=\varphi(1)=(\prod_vd_v)e(1)$ for all $x$. The matrix coefficients
are bounded by $1$, so the archimedean part of $\varphi$ is integrable
(\lref{GlobalTheory.T.discrete_series_coefficient_integrable_majorant}), and \ref{pf:rows:S3},
\ref{pf:rows:S4} (whose gap conditions are empty for $q=0$) give
$\rowsDc=\int_{\mathcal F}K_\varphi(x,x)\,d\mu_{\A}=(\prod_vd_v)e(1)\mu_{\A}(\mathcal F)$. Since $e(1)$ is real, this is $\rowsDst$.
\end{proof}

\paragraph{The volume bounds.}
Let $\rowsRes>0$ be the residue at $s=1$ of the Dedekind zeta function of $E$, $\zeta_E(l)$ its values, and for
$j\in\Z_{\ge1}^a$ with $J=\sum_ij_i$ put
\begin{gather*}
\rowsdimN(n,j)=4\sum_{i<k}j_ij_k+2J(n-2J)+\sum_ij_i^2,\\
\rowskap(E,p,q,j)=\frac{v_E^{-\frac f2\rowsdimN(p+q,j)}\,2^{-f((p+q)J-J^2+\frac J2)}}{\prod_{M=p+q-2J+1}^{p+q}L_\eta(M)}
\Bigl(\frac{s_ps_q}{s_{p-J}s_{q-J}\prod_is_{j_i}}\Bigr)^{f}\,\rowsRes^{\,a}\prod_i\prod_{l=2}^{j_i}\zeta_E(l),
\end{gather*}
where $s_i=\prod_{l=1}^i2\pi^l/(l-1)!$, and on $\Delta_a$ the density
\[
\rowsdens(n,f,j,\tau)=\prod_{i<k}\tau_i^{4j_ij_kf}\ \prod_i\tau_i^{2j_i(n-2J)f+2j_i^2f}\ \prod_i\frac{2fj_i}{\tau_i}.
\]
Finally let
$c_{\mathrm{Iw}}(p,q,j)=2^{-\rowsdimN(p+q,j)/2}\rowsvU(p)\rowsvU(q)/\bigl(\prod_i\rowsvU(j_i)2^{j_i^2/2}\cdot\rowsvU(p-J)\rowsvU(q-J)\bigr)$
(the constant of an Iwasawa decomposition).

\begin{lemma}[Iterated volume recursion]\label{pf:rows:volratio}
Assume $\sfname{VolumeStep}$. Let $I$ have signature $(p,q)$ and let $p_0\le p$, $q_0\le q$ with $p_0+q_0\ge2$. Then there is
an instance $I_0$ over the same field (up to an isomorphism commuting with the conjugations) of signature $(p_0,q_0)$, and
every such $I_0$ satisfies, for all gauge constants and Haar measures,
\[
\rowsgv(I)=\rowsgv(I_0)\prod_{M=p_0+q_0+1}^{p+q}\rowsrf(E,M).
\]
\formallabel{GlobalTheory.FlagMeas.VolRatio.volumeRatio_of_volumeStep}
\formallabel{GlobalTheory.T.hermVolumeRatioSameKind_of_volumeRatioInduction}
\end{lemma}

\begin{proof}
Induction on $d=(p-p_0)+(q-q_0)$. If $d=0$, take $I_0=I$; for every $I_0$ over the same field with the same signature,
$\rowsgv(I)=\rowsgv(I_0)$, because the gauge volume depends only on the field and the signature
(\lref{GlobalTheory.T.hermitian_volume_depends_only_on_field_signature}: the two hermitian spaces are related by a
semilinear isometry, which maps the self-dual lattice of one to a self-dual lattice of the form $\rowsLv_{(1,x_{\mathrm f})}$
for the other, and this identifies the gauge volumes). If $d>0$, then $n=p+q\ge3$; apply
$\sfname{VolumeStep}$ with $s=1$ if $p>p_0$ and $s=-1$ otherwise. It gives $I'$ over the same field, of signature $(p-1,q)$
resp.\ $(p,q-1)$, with $\rowsgv(I)/\rowsvU(n)^f=\rowsmr(E,n)\rowsgv(I')/\rowsvU(n-1)^f$, that is
$\rowsgv(I)=\rowsgv(I')\rowsrf(E,n)$ by \cref{pf:rows:constants}. The induction hypothesis for $I'$ gives $I_0$ and
the formula for $I'$; multiplying by $\rowsrf(E,n)$ gives the formula for $I$, since $\rowsrf(E,M)$ depends on $E$ only up
to isomorphism (\lref{GlobalTheory.FlagMeas.VolRatio.rankFac_congr}).
\end{proof}

\begin{lemma}[Absolute volume]\label{pf:rows:absvol}
Assume $\sfname{VolumeStep}$. Let $I$ have signature $(p,q)$ with $q^2\le n$ and $n\ge8$, and let $X\ge2$ with
$r_d\le X$ and $2fr_d\le X^2$. Then $\rowsgv(I)\le X^{n^2f}$ for all gauge constants and Haar measures.
\formallabel{GlobalTheory.FlagMeas.AbsVol.absoluteVolume_of_volumeStep}
\formallabel{GlobalTheory.T.hermAbsoluteVolumeBound_of_volume_ratio_induction}
\end{lemma}

\begin{proof}
Since $q^2\le n$ and $n\ge8$ we have $p\ge2$. By \cref{pf:rows:volratio} with $(p_0,q_0)=(2,0)$ there is an instance
$I_0$ over the same field of signature $(2,0)$ with $\rowsgv(I)=\rowsgv(I_0)\prod_{M=3}^{n}\rowsrf(E,M)$. For a definite
instance of rank $r$, the gauge volume is at most $\rowsvU(r)^f$ times the number of isometry classes of positive definite
unimodular hermitian $\OO_E$-lattices of rank $r$ (\lref{GlobalTheory.DefiniteVolume.gaugeVol_le}; here
$\rowsGinf=\Uni(r)^f$ is compact and $\rowsvU(r)$ is the volume of $\Uni(r)$ for the Hilbert--Schmidt metric,
\lref{GlobalTheory.T.riemannianHaarGauge_volume_unitaryGroup_eq}). In rank $2$ this class number is at most
$4^f(2v_E)^{5f}(1+f\log(2v_E))^{5f}$ (\lref{GlobalTheory.RankTwoCount.rankTwoClasses}, by the existence of short
vectors in rank-two unimodular lattices and a count of gluing data). An elementary estimate of
$\rowsvU(2)^f4^f(2v_E)^{5f}(1+f\log 2v_E)^{5f}\prod_{M=3}^n\rowsrf(E,M)$ under the hypotheses on $X$ gives the bound
(\lref{GlobalTheory.T.absolute_hermitian_volume_bound_arith}).
\end{proof}

\paragraph{Unfolding along isotropic flags.}
For a group $G$ with a measure $\mu$, $\Gamma\subseteq G$ and $A\subseteq G$ put
$\covol_\Gamma(A)=\inf\{\mu(D):D\text{ measurable, and every }a\in A\text{ has }\gamma a\in D\text{ for some }\gamma\in\Gamma\}$
(so $\covol_\Gamma(\mu_{\A})=\covol_\Gamma(\rowsGA)$). For $x\in\rowsGf$ let
$\Gamma_x=\{\gamma_\infty:\gamma\in\Gamma,\ x^{-1}\gamma_{\mathrm f}x\in K\}\subseteq\rowsGinf$; this is the image of the
stabiliser of the lattice $\rowsLv_{(1,x)}$.

\begin{lemma}[Covolumes of unions of slices]\label{pf:rows:slices}
Let $G_1,G_2$ be groups with $\sigma$-finite measures $\mu_1,\mu_2$, $\mu_2$ left invariant, $\Gamma\subseteq G_1\times G_2$ a
subgroup, $K\subseteq G_2$ a measurable subgroup with $\mu_2(K)<\infty$, $A\subseteq G_1\times G_2$, and let $(x_c)_{c\in C}$ be a
countable family in $G_2$ and $\Gamma_c,A_c\subseteq G_1$ such that (i) for every $g\in\Gamma_c$ there is $\gamma\in\Gamma$ with
$\gamma_1=g$ and $x_c^{-1}\gamma_2x_c\in K$, and (ii) for every $y\in A$ there are $c\in C$, $\gamma\in\Gamma$ and $k\in K$ with
$\gamma_2y_2=x_ck$ and $\gamma_1y_1\in A_c$. Then
$\covol_\Gamma(A)\le\mu_2(K)\sum_c\covol_{\Gamma_c}(A_c)$ (for the product measure, with $0\cdot\infty=0$).
\formallabel{GlobalTheory.FlagMeas.covol_prod_le_tsum_slices}
\end{lemma}

\begin{proof}
We may assume the right-hand side finite. Let $\eps>0$ and $\eps_c>0$ with $\sum_c\eps_c\le\eps$. For each $c$ choose a
measurable $D_c\subseteq G_1$ such that every $a\in A_c$ has $ga\in D_c$ for some $g\in\Gamma_c$ and
$\mu_2(K)\mu_1(D_c)\le\mu_2(K)\covol_{\Gamma_c}(A_c)+\eps_c$ (take $D_c=G_1$ if $\mu_2(K)=0$). Let
$D=\bigcup_cD_c\times x_cK$. For $y\in A$ choose $c,\gamma,k$ as in (ii), $g\in\Gamma_c$ with $g\gamma_1y_1\in D_c$, and
$\gamma'\in\Gamma$ as in (i) for $g$; then $\gamma'\gamma y=(g\gamma_1y_1,\ x_c(x_c^{-1}\gamma'_2x_c)k)\in D_c\times x_cK$. Hence
$\covol_\Gamma(A)\le\sum_c\mu_1(D_c)\mu_2(x_cK)\le\mu_2(K)\sum_c\covol_{\Gamma_c}(A_c)+\eps$.
\end{proof}

\begin{lemma}\label{pf:rows:dcsum}
\begin{enumerate}[label=(\alph*)]
\item Let $G_1,G_2$ be second countable topological groups, $G_1$ locally compact and Hausdorff, with $\sigma$-finite left
invariant Borel measures $\mu_1,\mu_2$; let $\Gamma\subseteq G_1\times G_2$ be a countable subgroup which meets every compact set in a
finite set and such that $\gamma_1=1$ implies $\gamma=1$, and $K\subseteq G_2$ a compact open subgroup. Let $x$ run over representatives
of the classes of the relation $\gamma_2y=y'k$ ($\gamma\in\Gamma$, $k\in K$) on $G_2$, and $\Gamma_x=\{\gamma_1:\gamma\in\Gamma,\
x^{-1}\gamma_2x\in K\}$. Then $\mu_2(K)\sum_x\covol_{\Gamma_x}(G_1)\le\covol_\Gamma(G_1\times G_2)$.
\item Let $I$ be an automorphic instance with gauge constants $c$ and a Haar measure $\mu_{\mathrm f}$. Then
$\sum_x\prod_vc_v\,\covol_{\Gamma_x}(\rowsGinf)\le\rowsgv(I)$, where $\covol_{\Gamma_x}$ is taken for $\prod_v\mu_v$ and $x$ runs over
representatives of $\{\gamma_{\mathrm f}:\gamma\in\Gamma\}\backslash\rowsGf/K$.
\end{enumerate}
\formallabel{GlobalTheory.FlagUnfold.sum_covol_le}
\formallabel{GlobalTheory.FlagUnfold.sum_latGroup_le}
\end{lemma}

\begin{proof}
(a) There are countably many classes, as $K$ is open and $G_2$ second countable. Each $\Gamma_x$ is a countable subgroup of $G_1$
meeting every compact set in a finite set (since $K$ is compact), so it has a strict measurable fundamental domain $F_x$ (every
orbit meets $F_x$, and $g\in\Gamma_x$, $y,gy\in F_x$ imply $g=1$; \lref{GlobalTheory.FundExists.exists_strict_domain}). Then
$X=\bigcup_xF_x\times xK$ is a strict fundamental domain for $\Gamma$: if $y\in F_x\times xK$ and $\gamma y\in F_{x'}\times x'K$,
then $x,x'$ lie in the same class, so $x=x'$, $\gamma_1\in\Gamma_x$, hence $\gamma_1=1$ and $\gamma=1$; and every $y$ is moved into $X$
by first moving $y_2$ into $xK$ and then $y_1$ into $F_x$ by an element of $\Gamma_x$. Hence
$\covol_\Gamma(G_1\times G_2)=(\mu_1\times\mu_2)(X)=\mu_2(K)\sum_x\mu_1(F_x)\ge\mu_2(K)\sum_x\covol_{\Gamma_x}(G_1)$
(\lref{GlobalTheory.covol_eq_measure_inter_fund}).
(b) $\Gamma$ is countable, meets every compact set in a finite set, and $\gamma_\infty=1$ implies $\gamma=1$
(\lref{GlobalTheory.RatDiscrete.subgroup_countable_discrete}); $K$ is compact and open. Apply (a) with
$G_1=\rowsGinf$, $G_2=\rowsGf$, and multiply by $\prod_vc_v/\mu_{\mathrm f}(K)$.
\end{proof}

A \emph{saturated isotropic flag of type $j$} in a lattice $\Lambda\subseteq E^n$ is a chain of $\OO_E$-submodules
$0=S_0\subseteq S_1\subseteq\dots\subseteq S_a$ with $\rank S_i-\rank S_{i-1}=j_i$, $S_i=\Lambda\cap ES_i$, and $S_a$ totally
isotropic. For such a flag in $\rowsLv_{(1,x)}$ and $g\in\rowsGinf$ let $\sigma_{x,S}(g)\in\R^a$ be the vector of sizes
$\rowssizeQ_{(g,x)}(S_i/S_{i-1})$, and let $\Gamma_{x,S}\subseteq\Gamma_x$ be the subgroup of elements preserving the flag at every
real place. Let $I_0$ be a second instance with an isomorphism $e:E_0\to E$ commuting with the involutions.
\emph{Adapted data} for $(x,S)$ consist of lattices $\Lambda_i\subseteq E^{j_i}$, a point $x_0\in\rowsGf(I_0)$ and matrices
$T_i,T_i'\in M_{n\times j_i}(E)$, $T_0\in M_{n\times n_0}(E)$ such that $n=n_0+2J$,
\[
T_i^*JT_k=T_i'^*JT'_k=0,\quad T_i^*JT'_k=\delta_{ik},\quad T_i^*JT_0=T_i'^*JT_0=0,\quad T_0^*JT_0=J_0^e,
\]
$\rowsLv_{(1,x)}$ consists of the vectors $\sum_iT_iu_i+\sum_iT'_iu'_i+T_0e(y)$ with $u_i\in\Lambda_i$, $u'_i$ in the dual
lattice $\{u':t^*u'\in\OO_E\ \forall t\in\Lambda_i\}$ and $y\in\rowsLv_{(1,x_0)}(I_0)$, and $S_k=\{\sum_{i\le k}T_iu_i:u_i\in\Lambda_i\}$.
Thus adapted data are a Witt decomposition of the lattice adapted to the flag. We write $\mathrm{sl}(\Lambda)$ for the covolume of
the stabiliser of a lattice $\Lambda\subseteq E^j$ in $\GL_j(E)$, acting on $\prod_{w\mid\infty}\GL_j(\C)$ with its Haar measure
$|\det|^{-2j}dx$, restricted to the slab $1\le\prod_w|\det x_w|^2<e$, and $\vol(\GL_j)=\sum_\Lambda\mathrm{sl}(\Lambda)$, the sum over
the classes of lattices $\Lambda\subseteq E^j$ modulo $\GL_j(E)$.

\begin{lemma}[Unfolding along isotropic flags]\label{pf:rows:unfoldflag}
Let $I,I_0$ be automorphic instances with an isomorphism $E_0\to E$ commuting with the involutions, $I$ of signature $(p,q)$
and $I_0$ of signature $(p-J,q-J)$, where $j\in\Z_{\ge1}^a$, $a\ge1$ and $J=\sum_ij_i\le\min(p,q)$. Let $c,c_0$ be gauge constants
and $\mu_{\mathrm f},\mu_0$ Haar measures of $I,I_0$, and $B\subseteq\Delta_a$ measurable. Then
\begin{multline*}
\prod_vc_v\frac{\covol_\Gamma(\mathcal X(j,B))}{\mu_{\mathrm f}(K)}\\
\le c_{\mathrm{Iw}}(p,q,j)^f(2v_E)^{\frac f2\rowsdimN(n,j)}\prod_i2^{fj_i^2}\vol(\GL_{j_i})\cdot\rowsgv(I_0)\int_B\rowsdens(n,f,j,\tau)\,d\tau.
\end{multline*}
\formallabel{GlobalTheory.T.isotropicFlag_unfolding_bound}
\formallabel{GlobalTheory.FlagUnfold.unfoldBound}
\end{lemma}

\begin{proof}
\emph{Classes.} Let $C$ be the set of tuples $t=((\Lambda_i)_i,x_0)$ of classes of lattices $\Lambda_i\subseteq E^{j_i}$ under
$\GL_{j_i}(E)$ and of double cosets $x_0\in\Gamma_0\backslash\rowsGf(I_0)/K_0$; it is countable (there are countably many lattices,
and $K_0$ is open in a second countable group). Fix representatives, and for every $t$ for which adapted data with these
representatives exist, fix such data $(x_t,S_t,\dots)$; let $C'\subseteq C$ be the set of these $t$.

\emph{Covering.} Let $y\in\mathcal X(j,B)$. The distinct members $S_{\le\tau_i}(y)$ of the canonical filtration form a saturated
isotropic flag $S$ of type $j$ in $\rowsLv_{(1,y_{\mathrm f})}$ with $\sigma_{y_{\mathrm f},S}(y_\infty)\in B$
(\lref{GlobalTheory.FlagSat.exists_isSatFlag}). The pair $(y_{\mathrm f},S)$ has adapted data
(\lref{GlobalTheory.T.exists_adaptedBasis_saturatedIsotropicFlag}: a saturated isotropic flag in the self-dual lattice of
a point of $I$ admits adapted data with respect to $I_0$), and these may be chosen with the fixed representatives of the classes
of their lattices and point (\lref{GlobalTheory.T.exists_adapted_basis_change_models}: adapted data can be transported
along $\GL_{j_i}(E)$ and along double cosets of $I_0$). Its class $t$ therefore lies in $C'$, and since $(x_t,S_t)$ and
$(y_{\mathrm f},S)$ have adapted data with the same lattices and point, there are $\gamma\in\Gamma$ and $k\in K$ with
$\gamma_{\mathrm f}y_{\mathrm f}=x_tk$ and $\sigma_{x_t,S_t}(\gamma_\infty g)=\sigma_{y_{\mathrm f},S}(g)$ for all $g$
(\lref{GlobalTheory.T.flags_conjugate_of_same_adapted_models}); in particular $\gamma_\infty y_\infty\in
A_t=\{g:\sigma_{x_t,S_t}(g)\in B\}$.

\emph{Slices.} \Cref{pf:rows:slices} with $G_1=\rowsGinf$, $G_2=\rowsGf$, $\Gamma_t=\Gamma_{x_t,S_t}$ (condition (i) holds because
$\Gamma_{x_t,S_t}\subseteq\Gamma_{x_t}$) gives $\covol_\Gamma(\mathcal X(j,B))\le\mu_{\mathrm f}(K)\sum_{t\in C'}\covol_{\Gamma_t}(A_t)$.

\emph{One class.} For $t\in C'$ with lattices $\Lambda_i$ and point $x_0$,
\[
\prod_vc_v\,\covol_{\Gamma_t}(A_t)\le c_{\mathrm{Iw}}^f(2v_E)^{\frac f2\rowsdimN}\prod_i2^{fj_i^2}\mathrm{sl}(\Lambda_i)\cdot
\prod_vc_{0,v}\covol_{\Gamma_{0,x_0}}(\rowsGinf(I_0))\cdot\int_B\rowsdens
\]
(\lref{GlobalTheory.T.flagClass_measure_bound_lattice_stabiliser}: the integration over the flag stabiliser, by the
Iwasawa decomposition at the real places, the covolume of the integral points of the unipotent radical and the covolume of
the Levi factor, which is the product of the groups $\GL_{j_i}$ and of the unitary group of $I_0$).

\emph{Summation.} Summing over $t\in C'\subseteq C$, the sum factors: $\sum_{\Lambda_i}\mathrm{sl}(\Lambda_i)=\vol(\GL_{j_i})$ by
definition, and $\sum_{x_0}\prod_vc_{0,v}\covol_{\Gamma_{0,x_0}}(\rowsGinf(I_0))\le\rowsgv(I_0)$ by \cref{pf:rows:dcsum}. Dividing the
slice inequality by $\mu_{\mathrm f}(K)\in(0,\infty)$ gives the assertion.
\end{proof}

\begin{proposition}[Flag measure bound]\label{pf:rows:flagmeasure}
Assume $\sfname{VolumeStep}$. Let $I$ have signature $(p,q)$ with $2q+2\le n$, $\mu_{\mathrm f}$ a Haar measure and
$\mathcal F$ a fundamental domain. Let $a\ge1$, $j\in\Z_{\ge1}^a$ with $J=\sum_ij_i\le q$, and $B\subseteq\Delta_a$
measurable. Then $\mathcal X(j,B)$ is measurable, $0<\mu_{\A}(\mathcal F)<\infty$, and
\[
\frac{\mu_{\A}(\mathcal X(j,B)\cap\mathcal F)}{\mu_{\A}(\mathcal F)}\le\rowskap(E,p,q,j)\int_B\rowsdens(n,f,j,\tau)\,d\tau .
\]
\formallabel{GlobalTheory.FlagMeas.Bound.flagMeasureBound}
\formallabel{GlobalTheory.T.isotropicFlagProbabilityBound_of_volumeRatio}
\end{proposition}

\begin{proof}
Measurability and $\Gamma$-invariance of $\mathcal X(j,B)$ are
\lref{GlobalTheory.T.pointsOfFlagTypeSet_measurable_and_invariant}, and $0<\mu_{\A}(\mathcal F)<\infty$ follows from
the volume recursion (\lref{GlobalTheory.T.unitary_adelic_quotient_volume_pos_lt_top}). Let $I_0$ be an instance
over the same field of signature $(p-J,q-J)$ (it exists since $p+q-2J\ge2$); by \cref{pf:rows:volratio},
$\rowsgv(I)=\rowsgv(I_0)\prod_{M=p+q-2J+1}^{p+q}\rowsrf(E,M)$. The core of the proof is the unfolding of the
measure of $\mathcal X(j,B)$ along isotropic flags, \cref{pf:rows:unfoldflag}, which bounds
$\prod_vc_v\covol_\Gamma(\mathcal X(j,B))/\mu_{\mathrm f}(K)$ by
$c_{\mathrm{Iw}}(p,q,j)^f(2v_E)^{\frac f2\rowsdimN(n,j)}\prod_i2^{fj_i^2}\vol(\GL_{j_i})\,\rowsgv(I_0)\int_B\rowsdens$; here
$\covol_\Gamma(\mathcal X(j,B))=\mu_{\A}(\mathcal X(j,B)\cap\mathcal F)$ since $\mathcal X(j,B)$ is measurable and $\Gamma$-invariant
(\lref{GlobalTheory.T.leftQuotientCovolume_eq_fundamentalDomain_inter}).
The volume of $\GL_j$ equals $\rowsRes v_E^{fj^2}\prod_{l=2}^j\zeta_E(l)$
(\lref{GlobalTheory.T.genLin_tamagawa_volume}). Dividing by $\rowsgv(I)=\prod_vc_v\,\mu_{\A}(\mathcal F)/\mu_{\mathrm f}(K)$
and simplifying the product of the five factors (\lref{GlobalTheory.T.prod_five_factors_hermitian_volume}) gives
$\rowskap(E,p,q,j)$.
\end{proof}

\begin{lemma}[The explicit constant]\label{pf:rows:kappa}
Let $E$ be as in \cref{pf:rows:instance}, $p+q=n$, $j\in\Z_{\ge1}^a$ with $J=\sum j_i\le q$, and assume $\log n\ge15$,
$q^2\le n$, $\log(2v_E)\le e^n$ and $\log f\le\sqrt n$. Then $\rowskap(E,p,q,j)\le e^{-\frac12Jnf\log n}$.
\end{lemma}

This is \lref{GlobalTheory.T.flag_explicit_constant_le_exp_neg_half}. Its proof bounds the factors of $\rowskap$
separately: $v_E^{-\frac f2\rowsdimN}\le e^{fJn/3}$, the power of $2$ by $1$, the ratio of the $s_i$ by
$\exp\bigl(fJ(\log2\pi-(n-2\sqrt n)(\log n-\log2\pi-1)+3.49+\frac{\sqrt n}4\log n)\bigr)$, the residue by
$\rowsRes\le(16\log 2v_E)^{2f}$ (\lref{GlobalTheory.T.dedekindZeta_residue_le_sixteen_log_rootDiscr_pow}; compare
\cite{Louboutin2000}), and the zeta and $L$-values by elementary estimates. The ratio of the $s_i$ contributes
$e^{-fJn\log n(1+o(1))}$ and dominates the other factors under the hypotheses.

\paragraph{Rare flags.}

\begin{lemma}\label{pf:rows:rare}
Assume $\sfname{VolumeStep}$. If $I$ has signature $(p,q)$ with $\log n\ge15$, $q^2\le n$, $\log r_d\le e^n$ and
$\log f\le\sqrt n$, then
$\mu_{\A}(\{x\in\mathcal F:S_{<1/2}(x)\ne0\})\le e^{-\frac13nf\log n}\mu_{\A}(\mathcal F)$. Consequently there is $n_3$ such that
$\mu_{\A}(\{x\in\mathcal F:\rowsleast(x)<\frac12\})\le e^{-\frac14nf\log n}\mu_{\A}(\mathcal F)$ whenever moreover $n\ge n_3$.
\formallabel{GlobalTheory.NonzeroFlag.nonzeroFlagBound}
\formallabel{GlobalTheory.T.nonzeroFlagProbabilityBound_of_volumeRatioInduction}
\formallabel{GlobalTheory.FlagRareProof.flagRare_of_nonzeroFlagBound}
\end{lemma}

\begin{proof}
Put $\Xi=nf\log n\ge15$; also $n\ge16$ and $2q+2\le n$. If $S_{<1/2}(x)\ne0$, then $x$ has a type $(j,\tau)$ with
$\tau\in\Delta_a$, where $j$ is a composition of $J=\rank S_{<1/2}(x)$ and $1\le J\le q$ by the isotropy of $S_{<1/2}(x)$
(\lref{GlobalTheory.NonzeroFlag.mem_typeSet}). For such $j$, \cref{pf:rows:flagmeasure} with $B=\Delta_a$,
$\int_{\Delta_a}\rowsdens\le1$ (\lref{GlobalTheory.T.lintegral_flagSizeDensity_le_one}) and \cref{pf:rows:kappa} give
$\mu_{\A}(\mathcal X(j,\Delta_a)\cap\mathcal F)\le e^{-\frac12J\Xi}\mu_{\A}(\mathcal F)$. There are $2^{J-1}$ compositions of $J$,
so the measure in question is at most $\sum_{J=1}^q2^{J-1}e^{-\frac12J\Xi}\mu_{\A}(\mathcal F)\le e^{-\frac13\Xi}\mu_{\A}(\mathcal F)$
for $\Xi\ge15$ (\lref{GlobalTheory.T.sum_two_pow_pred_mul_exp_neg_le_exp_neg_third}). For the second assertion take
$n_3=\lceil e^{15}\rceil$ and note that $\rowsleast(x)<\frac12$ if and only if $S_{<1/2}(x)\ne0$.
\end{proof}

\paragraph{Far terms.}

\begin{lemma}\label{pf:rows:far}
Assume the conclusion of \cref{pf:rows:flagmeasure} and the far bound \ref{pf:rows:S5} stated in the proof. Let $q\ge1$ and
assume $\rowsCH(C_g,8,d,p,q,b)$. Then $\int_{\rowsmod}\mathrm{fr}\,d\mu_{\A}\le e^{-n^2f}\rowsDst$.
\formallabel{GlobalTheory.Deep.farModerate_of_farBound}
\formallabel{GlobalTheory.T.farTermsModerateRegion_of_flagBound}
\end{lemma}

\begin{proof}
By \lref{GlobalTheory.FarTerms.farBound_of_ptBound} (a consequence of \ref{pf:rows:S2} and a lattice point count),
for $0<\eps\le(2q)^{-1/2}$ with $2n^2\eps\le\log2$, $64n^2\le h\eps$, $1\le h\eps^2/16$ and $h\ge120nq^{3/2}$, the far terms
(with respect to $\eps$) at every $x$ satisfy
\[
\mathrm{fr}(x)\le\Bigl(\prod_vd_v\Bigr)\Re e(1)\Bigl(\frac{3\sqrt n}{\tau_1(x)^2}\Bigr)^{2n^2f}2f\,3^f\,e^{-h\eps^2/16}.
\]
For $\eps=\eps_N$ the conditions hold (\lref{GlobalTheory.Deep.epN_pos_and_le_inv_sqrt_and_mul_le_log_two} and \lref{GlobalTheory.gap_threshold_pays_far_terms}).
For $x\in\rowsmod$ we have $\tau_1(x)\ge\tau_*$ (as $\tau_*<\frac12$), and the elementary inequality
$(3\sqrt n/\tau_*^2)^{2n^2f}2f3^fe^{-h\eps_N^2/16}\le e^{-n^2f}$, which follows from the size of the gap $h$
(\lref{GlobalTheory.gap_threshold_pays_far_terms}), gives $\mathrm{fr}(x)\le(\prod d_v)\Re e(1)e^{-n^2f}$. Integrating over
$\rowsmod\subseteq\mathcal F$ gives the claim, since $(\prod d_v)\Re e(1)\mu_{\A}(\mathcal F)=\rowsDst$ and $\mu_{\A}(\mathcal F)<\infty$.
\end{proof}

\paragraph{Near terms.}

\begin{lemma}\label{pf:rows:near}
Assume the conclusions of \cref{pf:rows:flagmeasure,pf:rows:kappa,pf:rows:nearunip} and \ref{pf:rows:S2}. For every
$C_g\ge0$ there is $n_1$ such that if $q\ge1$ and $\rowsCH(C_g,n_1,d,p,q,b)$ holds, then
$\int_{\rowsmod}\mathrm{nr}\,d\mu_{\A}\le e^{-\frac14nf\log n}\rowsDst$.
\formallabel{GlobalTheory.NearMod.nearModerate}
\end{lemma}

\begin{proof}
The analysis of the near terms for a single flag type is
\lref{GlobalTheory.T.nearTerms_moderate_points_one_type_bound}: under $\rowsCH(C_g,8,d,p,q,b)$ with
$q\ge1$, (a) $\mathrm{nr}(x)=0$ if $\rowsleast(x)\ge\frac12$, and (b) for every composition $j$ of some $J\le q$ into $a$
parts,
\[
\int_{\mathcal X(j,\Delta_a)\cap\mathcal F}\one[\rowsleast\ge\tau_*]\,\mathrm{nr}\,d\mu_{\A}\le
\rowsDst\cdot2\rowskap(E,p,q,j)\bigl(2fq\cdot11q\log X\bigr)^a .
\]
Both rest on \cref{pf:rows:nearunip}: near elements $\gamma\ne1$ lie in the set $N_\delta(x)$ of $\gamma\in\Gamma$ with
$(x^{-1}\gamma x)_{\mathrm f}\in K$ and $\|(x^{-1}\gamma x)_v-1\|_{\mathrm{HS}}\le\delta$ for all $v$, where $\delta=3\sqrt{n\eps_N}<\frac14$.
For (a), $N_\delta(x)=\{1\}$ when $\rowsleast(x)\ge\frac12$ (\lref{GlobalTheory.NearLower.nearSet_eq}). For (b), at a
point of type $(j,\tau)$ the number of elements of $N_\delta(x)$ is at most
$\prod_{i<k}\bigl(\vartheta(\tau_i/\tau_k)\vartheta(\tau_i\tau_k)\bigr)^{2j_ij_kf}\prod_i\vartheta(\tau_i)^{2j_i(n-2J)f}\vartheta(\tau_i^2)^{j_i^2f}$
with $\vartheta(t)=1+12\delta/t$ (\lref{GlobalTheory.NearCount.near_count}), each term is at most $(\prod_vd_v)\Re e(1)$ by
\ref{pf:rows:S2}, and this count is integrated against the flag measure bound of \cref{pf:rows:flagmeasure} over
$\tau_i\ge\tau_*$.
By (a) and the decomposition of $\{\rowsleast<\frac12\}$ into flag types
(\lref{GlobalTheory.TypeDecomp.lintegral_le_sum_types}), $\int_{\rowsmod}\mathrm{nr}$ is at most the sum of the
left-hand sides of (b) over $J\in[1,q]$ and the compositions $j$ of $J$. Choose $n_1\ge3^{15}$ such that, for $n\ge n_1$,
$Y:=2fq\cdot11q\log X\le e^{\frac18nf\log n}$ and $2\cdot2^qe^{-\frac38nf\log n}\le e^{-\frac14nf\log n}$
(\lref{GlobalTheory.moderate_log_factor_absorbed}; this uses $\log b\le e^{\frac1{16}nf\log n}$, $g\le b^{C_g}$,
$\log r_d\le e^n$ and $\log f\le\sqrt n$). Then $\log n\ge15$, and \cref{pf:rows:kappa} gives
$2\rowskap Y^a\le2e^{-\frac12Jnf\log n}e^{\frac a8nf\log n}\le2e^{-\frac38nf\log n}$ since $a\le J$
(\lref{GlobalTheory.NearMod.one_type_arith}). As there are fewer than $2^q$ compositions of integers $J\le q$
(\lref{GlobalTheory.sum_card_composition_lt}), the sum is at most $2^q\cdot2e^{-\frac38nf\log n}\rowsDst\le
e^{-\frac14nf\log n}\rowsDst$.
\end{proof}

\paragraph{The deep region.}
For $0<\tau$ put $N=n^2f$, $\rho(\tau)=\frac14n^2\sqrt f\,g\,r_d(2\tau)^{1/q}$, $\beta=\sqrt nb^2+\frac32n^{3/2}+3(8q+1)N$ and
\[
\Omega(\tau)=\bigl(2\rho(\tau)\sqrt f\beta\bigr)^{16qN}(6\beta)^Ne^f\Bigl(\frac{17\sqrt n}{\tau^2}\Bigr)^{2N}.
\]

\begin{lemma}\label{pf:rows:deepkernel}
Assume \ref{pf:rows:S2} and that the space of the count is finite-dimensional whenever $\rowsCH(C_g,n_0,d,p,q,b)$
holds with $n_0\ge8$. Let $q\ge1$, $\rowsCH(C_g,n_0,d,p,q,b)$ with $n_0\ge8$, $c$ gauge constants, and $x$ with
$\rowsleast(x)<\tau_*$. Then
$K_\varphi(x,x)\le\Omega(\rowsleast(x))\,\rowsDc\,\prod_vc_v/\mu_{\mathrm f}(K(g))$.
\formallabel{GlobalTheory.DeepKernel.deep_kernel_bound_of_ptBound}
\formallabel{GlobalTheory.T.deepPointKernelBound_of_pointwiseBound}
\end{lemma}

\begin{proof}
Choose isometric intertwiners $T_1,\dots,T_{\rowsDc}$ of the count whose images are cuspidal and fixed by $K(g)$ and such
that the vectors $T_i(\otimes w)$ are an orthonormal basis of the range of the projection of \ref{pf:rows:S3}
(\lref{GlobalTheory.T.exists_orthonormal_cuspidal_isometric_intertwiners}), and continuous representatives $F'_i$
of the $T_i(\otimes w)$ (\lref{GlobalTheory.T.exists_continuousRowRepresentatives_lowestKType}). The expression
$\sum_i|F_i(x)|^2$ of \ref{pf:rows:S4} does not depend on the choice of orthonormal basis
(\lref{GlobalTheory.Dec.kerDiag_basis_indep}), so $K_\varphi(x,x)=\sum_i|F'_i(x)|^2$. The bound for this sum at points
deep in the cusp, for any family of cuspidal isometric intertwiners, is
\lref{GlobalTheory.T.deepPoint_bound_cuspidalIsometricIntertwinerFamily}. We do not reproduce its proof; it combines a
Sobolev bound on a unipotent neighbourhood of $x$ of radius $\rho(\rowsleast(x))$, on which $\tau_1\ge\rowsleast(x)/2$, a
step bound for the derivatives of the lowest $K$-type vectors
(\lref{GlobalTheory.T.discreteSeries_stepBound_of_countingHypotheses}), a covering of the unipotent orbit of $x$
(\lref{GlobalTheory.DeepCover.unipotent_orbit_cover}) and an iteration which uses the cuspidality of the images
(\lref{GlobalTheory.DeepFamilyGap.iterate}).
\end{proof}

\begin{lemma}\label{pf:rows:deep}
Assume the conclusion of \cref{pf:rows:absvol} and the conclusion of \cref{pf:rows:deepkernel} for all data. If $q\ge1$ and
$\rowsCH(C_g,8,d,p,q,b)$ holds, then $\int_{\rowsdeep}K_\varphi(x,x)\,d\mu_{\A}\le\rowsDc\,e^{-n^2f}$.
\formallabel{GlobalTheory.DeepRegion.deepBound}
\formallabel{GlobalTheory.T.deepRegion_kernelDiagonal_bound}
\end{lemma}

\begin{proof}
Let $c$ be gauge constants. We have $X\ge16$, $r_d\le X$ and $2fr_d\le X^2$, so \cref{pf:rows:absvol} gives
$\prod_vc_v\,\mu_{\A}(\mathcal F)\le X^N\mu_{\mathrm f}(K)\le X^Ng^{2N}\mu_{\mathrm f}(K(g))$, using
$\mu_{\mathrm f}(K)\le g^{2n^2f}\mu_{\mathrm f}(K(g))$ (\lref{GlobalTheory.CongIndex.measure_Kset_le}). For $x\in\rowsdeep$,
$0<\rowsleast(x)\le\tau_*=X^{-11q}$, and the elementary inequality $\Omega(\tau)g^{2N}X^N\le e^{-N}$ for $0<\tau\le X^{-11q}$
(\lref{GlobalTheory.deep_region_exponents}) together with \cref{pf:rows:deepkernel} gives
$K_\varphi(x,x)\le e^{-N}\rowsDc\prod_vc_v/(g^{2N}X^N\mu_{\mathrm f}(K(g)))$. Integrating over $\rowsdeep\subseteq\mathcal F$ and using the
volume bound gives $\int_{\rowsdeep}K_\varphi(x,x)\le\rowsDc e^{-N}$.
\end{proof}

\paragraph{Finite-dimensionality.}

\begin{lemma}\label{pf:rows:findim}
Assume the decay estimate \ref{pf:rows:S1} and the conclusion of \cref{pf:rows:absvol}. If
$\rowsCH(C_g,n_0,d,p,q,b)$ holds with $n_0\ge8$, the space of the count is finite-dimensional.
\formallabel{GlobalTheory.RowsFinMult.rows_finiteDimensional}
\formallabel{GlobalTheory.T.finiteDimensional_rowsSpace_discreteSeriesCount}
\end{lemma}

\begin{proof}
By \cref{pf:rows:absvol}, $V=\mu_{\A}(\mathcal F)\in(0,\infty)$. Let $C$ be a constant such that the continuous
representatives $F$ of the vectors $T(\otimes w)$, for intertwiners $T$ of norm at most $1$ with $K(g)$-fixed images, satisfy
$|F(x)|\le C\tau_1(x)^{-2n^2f}$ (\lref{GlobalTheory.EvalBound.exists_contRep_bound_of_isKTypeOf}). Put
$m=\min(\frac12,\tau_*)$, $M=m^{-2n^2f}$ and $\epsilon=e^{-n^2f}$; then $\epsilon\le\frac12$. Let $S$ be a finite-dimensional
subspace of the space of the count, $k=\dim S$, and choose isometric intertwiners $T_1,\dots,T_k$ spanning $S$ with
orthonormal vectors $T_i(\otimes w)$ (\lref{GlobalTheory.IntertwinerFrame.exists_orthonormal_frame_of_le}); their images
are cuspidal and $K(g)$-fixed (\lref{GlobalTheory.T.idempotent_intertwiner_images_cuspidal_of_level}). For continuous
representatives $F_i$ of $T_i(\otimes w)$ we claim $\sum_i|F_i(x)|^2\le k\epsilon/V+C^2M^2$ for all $x$. If $q\ge1$ and
$\rowsleast(x)<\tau_*$, this follows from the deep point bound
(\lref{GlobalTheory.T.deepPoint_bound_cuspidalIsometricIntertwinerFamily}) and the inequality
$\Omega(\rowsleast(x))\prod_vc_v V/\mu_{\mathrm f}(K(g))\le\epsilon$
(\lref{GlobalTheory.T.measure_fundamentalDomain_lt_of_deepPoint_bound}). Otherwise $\tau_1(x)\ge m$ (if $q=0$, then
$\rowsleast(x)\ge\frac12$ for all $x$, \lref{GlobalTheory.RowsFinMult.half_le_least}), and Bessel's inequality for the
orthonormal family gives $\sum_i|F_i(x)|^2\le(C\tau_1(x)^{-2n^2f})^2\le C^2M^2$
(\lref{GlobalTheory.RowsFinMult.bessel}). Integrating over $\mathcal F$,
$k=\sum_i\|T_i(\otimes w)\|^2\le k\epsilon+C^2M^2V$, so $k\le2C^2M^2V$. A space all of whose finite-dimensional subspaces
have dimension at most $\lceil2C^2M^2V\rceil$ is finite-dimensional.
\end{proof}

\paragraph{Assembly.}

\begin{lemma}\label{pf:rows:split}
Let $\mu$ be a finite measure on a measurable space $Y$, $M\subseteq Y$ measurable, $\mathcal K\ge0$ integrable, $\mathrm{nr}$ and
$\mathrm{fr}$ functions with values in $[0,\infty]$ ($\mathrm{nr}$ measurable), and $c,a,b',e'\ge0$ with
$|\mathcal K-c|\le\mathrm{nr}+\mathrm{fr}$ on $M$, $\int_M\mathrm{nr}\le a$, $\int_M\mathrm{fr}\le b'$ and $\int_{Y\smallsetminus M}\mathcal K\le e'$.
Then $|\int_Y\mathcal K\,d\mu-c\,\mu(Y)|\le a+b'+c\,\mu(Y\smallsetminus M)+e'$.
\formallabel{GlobalTheory.abs_integral_sub_const_mul_le_of_split}
\end{lemma}

\begin{proof}
$\int_Y\mathcal K-c\mu(Y)=\int_M(\mathcal K-c)+\int_{Y\smallsetminus M}\mathcal K-c\mu(Y\smallsetminus M)$; the first term is bounded by
$a+b'$, and the remaining two terms lie in $[-c\mu(Y\smallsetminus M),e']$.
\end{proof}

\begin{lemma}\label{pf:rows:numerics}
Let $D,D_*\ge0$, $n\ge e^{15}$, $f\ge1$ and
$|D-D_*|\le(2e^{-\frac14nf\log n}+e^{-n^2f})D_*+e^{-n^2f}D$. Then $|D-D_*|\le e^{-\frac1{16}nf\log n}D_*$.
\formallabel{GlobalTheory.abs_sub_le_exp_neg_sixteenth_mul}
\end{lemma}

\begin{proof}
Put $\Xi=nf\log n$; then $\log n\ge15$, $\Xi\ge240$ and $\Xi\le n^2f$. Hence $e^{-\frac14\Xi}$ and $e^{-n^2f}$ are at most
$\frac18e^{-\frac1{16}\Xi}$ (as $\frac3{16}\Xi\ge\log8$), and $e^{-n^2f}\le\frac12$. Using $D\le D_*+|D-D_*|$,
\[
|D-D_*|\le\tfrac38e^{-\frac1{16}\Xi}D_*+\tfrac18e^{-\frac1{16}\Xi}D_*+\tfrac12|D-D_*|,
\]
so $|D-D_*|\le e^{-\frac1{16}\Xi}D_*$.
\end{proof}

\begin{lemma}\label{pf:rows:pointsplit}
For discrete series data and every $x\in\rowsGA$ at which $\sum_{\gamma\in\Gamma}|\varphi(x^{-1}\gamma x)|<\infty$,
$|K_\varphi(x,x)-\varphi(1)|\le\mathrm{nr}(x)+\mathrm{fr}(x)$; and $\mathrm{nr}$ is measurable.
\formallabel{GlobalTheory.NearMeas.nearAbs_measurable_and_split}
\end{lemma}

\begin{proof}
$K_\varphi(x,x)-\varphi(1)=\sum_{\gamma\ne1}\varphi(x^{-1}\gamma x)$, an absolutely convergent series; every $\gamma\ne1$ is either near
$x$ or far from $x$, and the triangle inequality gives the bound. $\Gamma$ is countable, and $\mathrm{nr}$ is a countable sum of
measurable functions, each the indicator of a measurable set times a continuous function.
\end{proof}

\begin{lemma}[Assembly]\label{pf:rows:assembly}
Assume: the finite-dimensionality statement of \cref{pf:rows:findim}; the conclusion of \cref{pf:rows:flagmeasure}; the
pointwise bound \ref{pf:rows:S2}; and the conclusions of \cref{pf:rows:near,pf:rows:far,pf:rows:rare,pf:rows:deep} (each
for every $C_g\ge0$ with some threshold for $n$). Then the conclusion of \cref{pf:rows:counting} holds.
\formallabel{GlobalTheory.CountSum.countingTheorem_of_finite_of_bounds}
\formallabel{GlobalTheory.T.discreteSeries_counting_of_finite_and_bounds}
\end{lemma}

\begin{proof}
Let $C_g\in\R$. If $C_g<0$, then $\rowsCH$ is never satisfied (it would give $1\le g\le b^{C_g}<1$), and $n_0=8$
works. Let $C_g\ge0$, let $n_1,n_2,n_3,n_4$ be the thresholds of the near, far, rare and deep statements, and put
$n_0=8+n_1+n_2+n_3+n_4+\lceil e^{15}\rceil$. Assume $\rowsCH(C_g,n_0,d,p,q,b)$. The space of the count is
finite-dimensional. If $q=0$, \cref{pf:rows:definite} gives $\rowsDc=\rowsDst$. Let $q\ge1$. Elementary consequences of the
hypotheses (\lref{GlobalTheory.Dec.countHyp_numerics}) are: $p\ge q$, $2q+2\le n$, $\prod d_v>0$, $X\ge16$,
$0<\tau_*<\frac12$, and the gap $h$ satisfies $h\ge120nq^{3/2}$, $8nq<h$ and the conditions of \ref{pf:rows:S4}. By
\ref{pf:rows:S2} and the majorant of \lref{GlobalTheory.Dec.PhMajProof.phMaj}, the archimedean part of $\varphi$ is
integrable, so \ref{pf:rows:S3} and \ref{pf:rows:S4} apply: $\rowsDc=\int_{\mathcal F}K_\varphi(x,x)\,d\mu_{\A}$, with
$K_\varphi(x,x)=\sum_i|F_i(x)|^2\ge0$ and the series converging absolutely. By \cref{pf:rows:pointsplit},
$|K_\varphi(x,x)-c|\le\mathrm{nr}(x)+\mathrm{fr}(x)$ with $c=\varphi(1)=(\prod_vd_v)\Re e(1)$. Apply \cref{pf:rows:split} to
$\mu=\mu_{\A}|_{\mathcal F}$ (finite by \cref{pf:rows:flagmeasure}), $M=\{\rowsleast\ge\tau_*\}$ (measurable, $\rowsleast$ being
continuous), $a=e^{-\frac14nf\log n}\rowsDst$ (\cref{pf:rows:near}), $b'=e^{-n^2f}\rowsDst$ (\cref{pf:rows:far}) and
$e'=\rowsDc e^{-n^2f}$ (\cref{pf:rows:deep}). Since $\tau_*<\frac12$, \cref{pf:rows:rare} gives
$c\,\mu_{\A}(\mathcal F\smallsetminus M)\le c\,e^{-\frac14nf\log n}\mu_{\A}(\mathcal F)=e^{-\frac14nf\log n}\rowsDst$. Therefore
\[
|\rowsDc-\rowsDst|\le\bigl(2e^{-\frac14nf\log n}+e^{-n^2f}\bigr)\rowsDst+e^{-n^2f}\rowsDc,
\]
and \cref{pf:rows:numerics} (with $n\ge e^{15}$) concludes.
\end{proof}

\begin{proof}[Proof of \cref{pf:rows:counting}]
Assume $\sfname{VolumeStep}$. Then \cref{pf:rows:absvol,pf:rows:flagmeasure} hold; \ref{pf:rows:S1} and \ref{pf:rows:S2}
hold unconditionally; \cref{pf:rows:findim} follows from \ref{pf:rows:S1} and \cref{pf:rows:absvol};
\cref{pf:rows:near} from \cref{pf:rows:flagmeasure,pf:rows:kappa,pf:rows:nearunip} and \ref{pf:rows:S2};
\cref{pf:rows:far} from \cref{pf:rows:flagmeasure} and \ref{pf:rows:S5}; \cref{pf:rows:rare} from the volume recursion;
and \cref{pf:rows:deep} from \cref{pf:rows:absvol,pf:rows:deepkernel}. \Cref{pf:rows:assembly} concludes.
\end{proof}

\begin{remark}\label{pf:rows:countremark}
(1) In the definite case the count equals the main term exactly. (2) The measures $\mu_v$ on $\Uni(\eps_v)$ are a fixed choice
of Haar measures; the formal degrees and $\rowsDst$ refer to them, and $\rowsDst$ does not depend on the choice. (3) Limit multiplicity theorems for discrete series \cite{DeGeorgeWallach1978,Clozel1986,Savin1989} and the Plancherel
equidistribution theorems of \cite{Shin2012,ShinTemplier2016} concern a fixed group with varying level (and weight), with
power-saving error terms in \cite{ShinTemplier2016}; \cref{pf:rows:counting} is uniform in the rank, the field and the
level, with relative error $e^{-\frac1{16}nf\log n}$, at the price of the gap condition.
\end{remark}

% ---------------------------------------------------------------------------
\subsection{The counting row}\label{sec:pf:rows:countrow}

The row $\mathsf{c30}$ asks for counting constants with which the counting statement holds at all points of the
designed families, and for an irreducibility count. The counting statement of the families is an ``unbundled'' form of
\cref{pf:rows:counting}, in which the data are not packaged and the hypotheses are slightly rearranged.

\begin{lemma}[Unbundled counting]\label{pf:rows:unbundled}
Assume the conclusion of \cref{pf:rows:counting}. For every $A\in\R$ there are $n_0\in\N$ and $C_0\ge0$ such that for all
$n_1\ge n_0$ and $C\ge C_0$ the following holds. Let $I$ be an instance with $n\ge n_1$ of signature $(p,q)$ with $q^2\le
n$, let $b\ge2$ with $\log\log r_d\le n$ and $\log\log b\le\frac1{16}nf\log n$, $\log f\le\sqrt n$, let
$\lambda_v$ be Harish-Chandra parameters with $|\lambda_v(a)|\le b^2$, $e$ an admissible idempotent of level $g$ and type
$g_0$ with $1\le g\le b^A$, $g_0$ a prime with $4n<g_0$ such that every prime of $\OO_E$ above $g_0$ has norm $g_0$ and
is unramified, and let the remaining counting data be given, with discrete series $\sigma_v$ of parameter $\lambda_v$ and
unit vectors $w_v$ of their lowest $K$-types. If all $\lambda_v$ have gap at least $h$ and $Cn^{19}f^5\log(br_d)\le h$, then
$|\rowsDc-\rowsDst|\le e^{-\frac1{16}nf\log n}\rowsDst$, where $\rowsDst$ is formed with the formal degrees of the $\sigma_v$.
\formallabel{GlobalTheory.CountBridge.count_unbundled}
\formallabel{GlobalTheory.T.unitaryAutomorphicCounting_unbundled}
\end{lemma}

\begin{proof}
Take $n_0$ from \cref{pf:rows:counting} for $C_g=A$ and $C_0=\max(0,8\cdot10^{14}(A+4))$. The formal degree $d_v$ of
$\sigma_v$ exists by Schur orthogonality for square-integrable irreducible unitary representations
(\lref{GlobalTheory.schur_orthogonality_unitary}) and is the one computed from $w_v$. A prime $g_0$ all of whose prime
divisors in $\OO_E$ have norm $g_0$ and are unramified splits completely
(\lref{GlobalTheory.T.ncard_primesOver_eq_finrank_of_unramified_split}). The conditions on $\log\log r_d$ and
$\log\log b$ are equivalent to $\log r_d\le e^n$ and $\log b\le e^{\frac1{16}nf\log n}$, and since $\log(br_d)\ge0$ the
gap is at least $8\cdot10^{14}(A+4)n^{19}f^5\log(br_d)$. Thus $\rowsCH(A,n_0,d,p,q,b)$ holds and
\cref{pf:rows:counting} applies.
\end{proof}

\begin{lemma}[Counting constants]\label{pf:rows:countconst}
Assume the Siegel--Weil formula in the convergent range. There are a map $c_{30}$ from finite lists of real numbers to
systems of counting constants $(C_{\mathrm{res}},n_0,A,C_b,C,c,c_b)$ and a constant $c_{Ic}>0$ such that for every family
$\mathcal F$ with list of output constants $\ell$: $c_{30}(\ell).C\ge0$, $c_{30}(\ell).A\ge2(C_{YL}(\mathcal F)+1)$,
$c_{Ic}\le c_{30}(\ell).c$, and for every additional predicate on the data, at every point $D_{x,i}$ of $\mathcal F$ the
counting statement holds with the constants $(n_0,A,C,c,c_b)$ of $c_{30}(\ell)$: for all counting data satisfying the
hypotheses of \cref{pf:rows:unbundled} with $n\ge n_0$, $g\le b^A$, gap at least $h\ge Cn^{19}f^5\log(br_d)$ and
$\log\log b\le c_b\,nf\log n$, one has $|\rowsDc-\rowsDst|\le e^{-c\,nf\log n}\rowsDst$.
\formallabel{GlobalTheory.CountConstants.count_constants_of_siegelWeil}
\formallabel{GlobalTheory.CountBridge.count_typed_of_volumeStep}
\end{lemma}

\begin{proof}
By \cref{pf:rows:volstep,pf:rows:counting,pf:rows:unbundled} there are, for every $A$, numbers $n_0(A)$ and $C_0(A)\ge0$
such that the counting statement holds with $(n_1,A,C,\frac1{16},\frac1{16})$ for all $n_1\ge n_0(A)$, $C\ge C_0(A)$ and
every datum (the datum plays no role in the statement). Put $A(\ell)=2(\ell_7+1)$, where $\ell_7$ is the seventh entry of
$\ell$ ($0$ if there is none; for the output constants of a family it is $C_{YL}$), and
$c_{30}(\ell)=(0,n_0(A(\ell)),A(\ell),0,C_0(A(\ell)),\frac1{16},\frac1{16})$, $c_{Ic}=\frac1{16}$.
\end{proof}

\begin{proposition}[Irreducibility row]\label{pf:rows:irr}
Assume the Siegel--Weil formula in the convergent range and the finite-dimensionality of row spaces
(\lref{GlobalTheory.rowsFinDim_link}: the spaces of automorphic forms attached to the rows are finite-dimensional).
There is $c_{Ic}>0$ such that for every class $\mathcal C$ all of whose members satisfy $\sfname{TuSize}$ the following holds.
If, closedly on $\mathcal C$ and eventually in $x$, every point $D_{x,i}$ carries label data (a Hida setting $X$ for
$D_{x,i}$, a good operator and an isometry between the row space of $D_{x,i}$ and that of $X$, with automorphic,
row-dictionary and label ties), then, closedly on $\mathcal C$ and eventually, every $D_{x,i}$ has a system of all rows $\mathbf A$
(labels at the two type primes by polynomials over $\F_{g_k}$, idempotents of $\C[\GL_r(\F_{g_k})]$ and cylinder
idempotents $e_T$) whose good labels are injective, which satisfies the join property, and whose cylinder counts satisfy
\[
|\dim_T-\rowsDst(e_T)|\le e^{-c_{Ic}\,rf\log r}\,\rowsDst(e_T)\qquad\text{for all cylinders }T,
\]
with $0<\rowsDst(e)$, $r$ and $f$ the rank and degree of $D_{x,i}$.
\formallabel{GlobalTheory.IrrRowConcrete.irr_row_closed}
\formallabel{GlobalTheory.T.irreducibilityRow_of_concreteAutomorphicData}
\end{proposition}

\begin{proof}
Take $c_{30},c_{Ic}$ from \cref{pf:rows:countconst}. Closedly on $\mathcal C$, eventually, the size hypotheses of the
counting statement with the constants $c_{30}$ hold at all points (\lref{GlobalTheory.IrrCountEv.countClauses_closed},
from $\sfname{TuSize}$), and so does the gap condition (\lref{Formalism.L2.SG.GapClosed.gapEv_closed_of}). Intersect
these closed statements with the hypothesis (taking the minimum of the exponents and the larger thresholds, and
$x_0$ the maximum of the three thresholds and $x_1+1$). At every point beyond $x_0$, the system of all rows is
constructed from the label data, and the counting statement applied to each cylinder idempotent gives the counts
(\lref{GlobalTheory.AllRowsAtRecord.allRows_at_record}).
\end{proof}

\begin{proposition}[Counting row]\label{pf:rows:countrow}
Assume the Siegel--Weil formula in the convergent range. Let $\mathrm{KP}$ be a base predicate, $Q,Q'$ slots and $\rowsLk$
a linkage of bundles. If the label row holds under $\rowsLg(\mathrm{KP},Q,Q',\cdot)$, then so does the row
$\mathsf{c30}$: there are $c_{30}$ and $c_{Ic}>0$ as in \cref{pf:rows:countconst} such that, for the same auxiliary slot,
$\rowsRow(S,\rowsLk,H)$ holds, where $H(\mathcal A,\sigma)$ is the conjunction of the counting statement at all points of
$\rowsDiii(\mathcal A)$ with the constants $c_{30}$ (closedly) and of the conclusion of \cref{pf:rows:irr} for
$\rowsDiii(\mathcal A)$ with $c_{Ic}$.
\formallabel{GlobalTheory.CountRowByName.row_counts_of_labels_alt}
\formallabel{GlobalTheory.CountRow.row_counts_of_ties_alt}
\end{proposition}

\begin{proof}
Write the label row as $\rowsLBC(\mathrm{KP},Q,R)$ and $\rowsRow(Q'\wedge R,\rowsLk,\rowsAtAll(\text{label data}))$; we
show that $\rowsRow(S,\rowsLk,\rowsAtAll(\text{label data}))$ implies the counting row for every slot $S$. The first
conjunct holds for all families (\cref{pf:rows:countconst}), hence closedly. For the second, let $\sigma$ be a selection
for the row of label data over $\mathcal A$. Every member of $\rowsDiii(\mathcal A)$ is a member of $\mathcal A$ and
satisfies $\sfname{TuSize}$, and the label data exist at all points of the members of $\mathcal A$ (with $x_0=0$), so the
hypothesis of \cref{pf:rows:irr} holds for $\mathcal C=\rowsDiii(\mathcal A)$; its conclusion is the second conjunct.
\end{proof}

% ---------------------------------------------------------------------------
\subsection{Period entries and level valuations}\label{sec:pf:rows:periods}

The six period rows hold for all points over \emph{true column choices}: the data assembled from a column choice whose
columns are given by the explicit theta and kernel constructions (true columns), with laws and later stages. We say
that a slot $Q$ lies \emph{over columns} if every core satisfying $Q$ arises from a column choice $w$ and a later tuple
$Z$ such that the twisting datum, the norm choice and the row integrality of the core are those determined by $(w,Z)$
(the ``column view'').

\begin{proposition}[\cref{prop:families:periods}]\label{pf:rows:periods}
For every slot system $\mathcal S$ and every linkage $\rowsLk$ of bundles, the following six rows hold under
$\rowsLg(\rowsKPx,\rowsBuilt(\mathcal S),\rowsGood(\mathcal S),\cdot)$, with requirements of the form
$\rowsEvC(R)(\mathcal A)$ (\cref{pf:rows:defrow}; for (5), in the form in which the condition $x>x_1$ is omitted):
\begin{enumerate}[label=(\arabic*)]
\item $\mathsf{aPt0}$: for every word $u$, the reference-chain evaluation statement: the evaluation map of the reference
chain of $u$ is injective, the stabiliser maps are bijective, the period points of the forms of level $m$ are integral,
and the period lattice and its majorants are bounded by those of the forms of level $m$ up to an explicit factor;
\item $\mathsf{aPtTw}$: the twisted period points are integral;
\item $\mathsf{aPer0}$: for every word $u$, every point $x$ of the period lattice and every form $F$ of the space
$V_0(\mu_0)$, the period of $x$ against $F$ is the value of some $y\in\mathit{LL}$ with $y=0$ or $v_w(y)$ bounded below
explicitly at every place $w$;
\item $\mathsf{aPerTw}$: the same for the twisted periods, at the places outside the twisting primes;
\item $\mathsf{tie}$: the real-place second-layer conditions (the real ties of the datum and symmetric column forms with a
finite character at step zero);
\item $\mathsf{aP10}$: for every word $u$, every place $w$ of $\mathit{LL}$ above the level prime $T_u$ and every
$c\in\mathit{LL}$ with $v_w(c)\ge e_w\,m\,E_T(u,w)$, $c\cdot\mathrm{col}_T(u,\varphi)$ lies in the lattice $S_m(w)$ for all
$\varphi$ in the relevant range; and the twisted columns at $T_u$ are bounded below by $-m\,E_{\mathrm{at}T}(u,w)$ for the used
tuples $\mu\ne\mu_0$ whose prime set does not contain $T_u$.
\end{enumerate}
\formallabel{GlobalTheory.T.period_entries_tie_and_level_valuations_at_built}
\formallabel{GlobalTheory.Fc.six_rows_at_member_alt}
\end{proposition}

The statement (6) is a divisibility at the level primes of order proportional to the chain length $m$; it carries most
of the weight. We prove the proposition modulo the six \emph{closing statements} for true column choices
(\cref{pf:rows:closing}).

\begin{lemma}[Closing statements]\label{pf:rows:closing}
Let $B$ be a base, $I$ a true column choice for $B$ (with the properties ``ok'', jets in range and pins), $l$ a law of
$I$, $Z$ a later tuple and $D$ the datum assembled from $(I,l,Z)$, with the twisting data determined by them. Then $D$
satisfies the properties (1)--(6) of \cref{pf:rows:periods} (for (5), under the real ties of $D$; for (6), for every
division datum of $D$ and under the sign conditions of the base).
\end{lemma}

The first five assertions are, respectively, \lref{GlobalTheory.T.aux_of_ok_of_pins},
\lref{GlobalTheory.T.twistedPointRow_admissible_trueColumns}, \lref{GlobalTheory.T.periodClassSumWeyl_of_rowInt},
\lref{GlobalTheory.T.admissible_period_twist_of_integral_rows} and
\lref{GlobalTheory.T.exists_symmetricColumnForms_stepZeroFiniteCharacter} (together with the real ties); we do not
reproduce their proofs. The sixth, the largest part of the region, is the closing theorem \cref{pf:rows:closingsharp}, proved in
\cref{sec:pf:rows:level} through the sharp bound for the last step of the true column chains, the five laws of the true chains,
the integrality of the moments of kernel jets at the level places, and the moments of the translates of the unit periods of the
CM theta functions, which use the CM theory of the theta functions recalled after \cref{prop:families:periods}.

\begin{lemma}[True-column view]\label{pf:rows:view}
Let $R(D,O)$ be a property of data with objects which holds for $(D,O)=(D(w,Z),O)$ whenever $w$ is a column choice, $Z$
a later tuple, $O$ an object over $D(w,Z)$ and the column view, the sign conditions of the base and the real ties hold.
Then $R(\beta.D,\beta.O)$ holds for every bundle point $\beta$ over columns.
\formallabel{GlobalTheory.T.property_at_constructedPoints_of_trueColumnsView}
\formallabel{GlobalTheory.Fc.Cv.view_at_built}
\end{lemma}

\begin{proof}
By induction on the position of the core of $\beta$: the position arises from a column choice $w$ and a later tuple $Z$
with the column view; the datum of $\beta$ is assembled from $(w,Z')$, where $Z'$ differs from $Z$ only by the division
stage determined by the level of the ledger point of $\beta$, and the twisting datum of $\beta$ is the transport of that of
$(w,Z)$. The sign conditions of the base hold at all bundle points
(\lref{GlobalTheory.Fc.Sg.baseSigns_of_bPt}), and the real ties are part of the core.
\end{proof}

\begin{proof}[Proof of \cref{pf:rows:periods}]
Cores at good positions of $\mathcal S$, also intersected with any slot, lie over columns
(\lref{GlobalTheory.T.goodTwistCores_over_trueColumns}: a good position is in particular a built position, and built
positions come from column choices). Take the auxiliary slot $R=\mathrm{True}$ in $\rowsLg$. Let $Q$ be any slot over columns
and $R_k$ the property (k) of \cref{pf:rows:periods} ($k=1,\dots,6$). By \cref{pf:rows:closing,pf:rows:view}, $R_k$ holds at
every bundle point $\beta\models Q$ of every context, so the hypothesis of \cref{pf:rows:calculus}(e) holds trivially, and
$\rowsRowZ(Q,\rowsEvC(R_k))$; by \cref{pf:rows:calculus}(d) this gives $\rowsRow(Q,\rowsLk,\rowsEvC(R_k))$ for every bundle type
and linkage. For (5) the condition $x>x_1$ is removed by raising the threshold to $x_1+1$. These are the six row
statements over columns
(\lref{GlobalTheory.T.referenceChainEvaluationAtBundlePoints_of_overCols}, \lref{GlobalTheory.T.twistedChainEvaluation_rational_integral_bundlePoints}, \lref{GlobalTheory.T.referencePeriodIntegral_bundlePoints_of_overColumns}, \lref{GlobalTheory.T.twistedPeriods_integral_at_bundlePoints}, \lref{GlobalTheory.T.realSideColumnsCompatible_bundlePoints_overColumns}, and \lref{GlobalTheory.T.levelPrimeColumnValuations_bundlePoints_overColumns}).
\end{proof}

% ---------------------------------------------------------------------------
\subsection{Level-prime valuations over true columns}\label{sec:pf:rows:level}

In this subsection we prove assertion (6) of \cref{pf:rows:closing}, and with it the row $\mathsf{aP10}$ over columns. The
argument has four layers. The row follows from a \emph{closing theorem} for true column choices by the transfer lemma
(\cref{pf:rows:levelrow}). The closing theorem reduces the valuation bounds at $\mu_0$ and at the twisted tuples to a sharp
bound for every twisted chain (\cref{pf:rows:closingsharp}). The sharp bound follows from the stability of lattice column
chains which are realised by \emph{built data} carrying kernels with five laws (\cref{pf:rows:laststep}), and the five laws
hold for the true chains (\cref{pf:rows:fivelaws}). Two of the laws, the integrality of the kernels at the places above the
level prime, rest on the integrality of the moments of kernel jets (\cref{pf:rows:jetmoments}), which in turn comes from
$p$-adic measures attached to the translates of the theta functions of CM elliptic curves (\cref{pf:rows:unitperiods}).

\paragraph{Chains and the sharp bound.}
Let $D$ be a datum and $u$ a word, with level prime $T_u$. A \emph{chain} $X$ for $(D,u)$ consists of a column chain with
column map $\mathrm{col}_X$ from the forms of level $0$ to the forms of level $m$, columns at a prime $T_X$, a period map, and
member forms: for every level $i\le m$ and every place $w$ of $\mathit{LL}$ an $\OO_w$-lattice $S^X_i(w)$ of forms of level $i$,
together with the embeddings of these lattices into the spaces of forms. The \emph{own chain} of $u$ is the chain given by the
member of $D$ at $u$ (with $T_X=T_u$); twisting forms $\mathrm{Tw}$ of $D$ give, for every tuple $\mu$, the \emph{twisted
chain} $\mathrm{tw}(\mathrm{Tw},u,\mu)$, with twisted column map $\mathrm{col}^{\mathrm{tw}}_\mu$ and an initial map $e_{0,\mu}$
from the forms of level $0$ of the own chain to those of the twisted chain. For a place $w$ above $T_X$ let
$E_T(X,w)=E_\pi(X)$ if $w$ is \emph{seen} by the member forms of $X$ and $E_T(X,w)=0$ otherwise, where
\[
E_\pi(X)=2\sum_\tau W_\tau(X)+fN_m+\tfrac14\,|\mathrm{tops}|\,(t+1)(m+1)
\]
and $W_\tau(X)\in\N$ is the maximum over the levels $i\le m$ and the indices $a'$ of $|\lambda_{i,\tau,a'}-i(t+1)/2|$ for the weights
$\lambda_{i,\tau}$ of the member forms of $X$ (the index $a'=0$ omitted at the top embeddings $\tau\in\mathrm{tops}$). The twisted
chains have the same weights as the own chain, hence the same $E_\pi$. Let $N_\lambda(X,w)$ be the lattice of forms of level $0$
obtained from the ramified lattice at $w$ by applying $T_X^{e(X,w)}U^m$ (the \emph{$\lambda$-lattice}). For the own chain we write
$E_{\mathrm{at}T}(u,w)=E_T(\mathrm{own},w)$.

\begin{definition}\label{pf:rows:defsharp}
\begin{enumerate}[label=(\alph*)]
\item A chain $X$ satisfies the \emph{sharp bound} if for every place $w$ above $T_X$, every $c\in\mathit{LL}$ with
$v_w(c)\ge e_w\,m\,E_T(X,w)$ and every $x\in N_\lambda(X,w)$, the form $c\cdot\mathrm{col}_X(x)$ is the image of an element of
$S^X_m(w)$. The datum $D$ satisfies the sharp bound at $u$ if its own chain does.
\item Twisting forms $\mathrm{Tw}$ satisfy the \emph{twisted column bound at $T$} for a word $u$ and a division datum $Y$ if for
every place $w$ above $T_u$, every used tuple $\mu\ne\mu_0$ such that the residue characteristic of $w$ is not in the prime set
$Q_p(\mu)$ of $Y$, and every $x\in N_\lambda(\mathrm{own},w)$, the twisted column $\mathrm{col}^{\mathrm{tw}}_\mu(e_{0,\mu}x)$
equals $a\cdot z$ with $z$ in the localisation at $w$ of the level-$m$ lattice of the twisted forms and $a\in\mathit{LL}$ with
$a=0$ or $v_w(a)\ge-e_w\,m\,E_{\mathrm{at}T}(u,w)$.
\item The \emph{step-zero tie} holds for $(\mathrm{Tw},u,\mu)$ if at every place $w$ above $T_u$ the twisted chain sees $w$ if and
only if the own chain does, and $e_{0,\mu}$ maps $N_\lambda(\mathrm{own},w)$ into $N_\lambda(\mathrm{tw}(\mathrm{Tw},u,\mu),w)$.
\end{enumerate}
\end{definition}

Thus (6) of \cref{pf:rows:periods} is the conjunction, for all words, of (a) for $D$ and (b) for the twisting forms of the
object. \emph{True column data} consist of a base $B$ with a pre-stage $X$, a true column choice $I$ for $(B,X)$ which is ok
and has its jets in range, a law $l$ of $I$ with an initial law $l_0$, and a later tuple $Z$; they determine the datum
$D=D(I,l,Z)$ and its twisting forms $\mathrm{Tw}=\mathrm{Tw}(I,l,l_0,Z)$. The \emph{sign conditions} of $B$ are the conditions on
the signs of the weights of the base used in \cref{pf:rows:view}.

\begin{lemma}[Twisted columns]\label{pf:rows:twcol}
Let $D$ be a datum with twisting forms $\mathrm{Tw}$, $u$ a word, $\mu$ a tuple and $w$ a place above $T_u$.
\begin{enumerate}[label=(\alph*)]
\item $mE_{\mathrm{at}T}(u,w)\in\Z$.
\item If the step-zero tie holds for $(\mathrm{Tw},u,\mu)$ and $\mathrm{tw}(\mathrm{Tw},u,\mu)$ satisfies the sharp bound, then for
every $c\in\mathit{LL}$ with $v_w(c)\ge e_wmE_{\mathrm{at}T}(u,w)$ and every $x\in N_\lambda(\mathrm{own},w)$ there is $y$ in the
level-$m$ lattice of the twisted forms at $w$ whose image is $c\cdot\mathrm{col}^{\mathrm{tw}}_\mu(e_{0,\mu}x)$.
\item If the conclusion of (b) holds, then for every $x\in N_\lambda(\mathrm{own},w)$ the twisted column
$\mathrm{col}^{\mathrm{tw}}_\mu(e_{0,\mu}x)$ satisfies the bound of \cref{pf:rows:defsharp}(b).
\end{enumerate}
\formallabel{GlobalTheory.Proof.ColumnWeightTwisted.exists_int_m_mul_EatT}
\formallabel{GlobalTheory.LvT.TwSharp.tupleChainSharpColumnValuation_of_twisted}
\formallabel{GlobalTheory.Proof.ColumnWeightTwisted.colGeTw_of_tupleChainSharpColumnValuation}
\end{lemma}

\begin{proof}
(a) At the places not seen, $E_{\mathrm{at}T}=0$. Otherwise write $m=2a$ and $t=2b+1$; then
$mE_{\mathrm{at}T}(u,w)=m(2\sum_\tau W_\tau(\mathrm{own})+fN_m)+|\mathrm{tops}|(b+1)(m+1)a$.
(b) By the tie, $w$ is seen by the twisted chain if and only if it is seen by the own chain; as the two chains have the same
$E_\pi$, $E_T(\mathrm{tw}(\mathrm{Tw},u,\mu),w)=E_{\mathrm{at}T}(u,w)$. Also $e_{0,\mu}x\in N_\lambda(\mathrm{tw}(\mathrm{Tw},u,\mu),w)$. The
sharp bound of the twisted chain gives $y$ whose image is $c$ times the column of $e_{0,\mu}x$ in the twisted chain, and this
column is $\mathrm{col}^{\mathrm{tw}}_\mu(e_{0,\mu}x)$, since the embedding of the twisted forms is compatible with the block
embedding of the twisted columns.
(c) Let $k=mE_{\mathrm{at}T}(u,w)\in\Z$ and choose $c\in\mathit{LL}$ with $v_w(c)=e_wk$ (the valuation is surjective); then $c\ne0$.
With $y$ as in (b), $\mathrm{col}^{\mathrm{tw}}_\mu(e_{0,\mu}x)=c^{-1}y$ and $v_w(c^{-1})=-e_wk$.
\end{proof}

\begin{theorem}[Closing theorem]\label{pf:rows:closingsharp}
Let $(B,X,I,l,l_0,Z)$ be true column data satisfying the sign conditions of $B$, with datum $D$ and twisting forms
$\mathrm{Tw}$. For every division datum $Y$ of $D$ and every word $u$, $D$ satisfies the sharp bound at $u$ and $\mathrm{Tw}$
satisfies the twisted column bound at $T$ for $u$ and $Y$.
\formallabel{GlobalTheory.LvT.Built.TrueCols.closing_true_cols}
\formallabel{GlobalTheory.T.trueColumns_row_sharp_bounds}
\formallabel{GlobalTheory.LvT.Built.TrueCols.row_true}
\end{theorem}

\begin{proof}
The own chain of $u$ is, by construction, the twisted chain $\mathrm{tw}(\mathrm{Tw},u,\mu_0)$. By \cref{pf:rows:laststep} with
the tuple $\mu_0$ it satisfies the sharp bound, which is the first assertion.

Let $w$ lie above $T_u$ and let $\mu$ be a tuple. The step-zero tie holds for $(\mathrm{Tw},u,\mu)$
(\lref{GlobalTheory.T.twistedChain_initialLattice_compatible_trueColumns}: for an ok true column choice the twisted
chains are compatible with the initial lattices at every tuple), and $\mathrm{tw}(\mathrm{Tw},u,\mu)$ satisfies the sharp bound by
\cref{pf:rows:laststep}. By \cref{pf:rows:twcol}(b),(c), the twisted column bound holds at $w$ for $\mu$ and every
$x\in N_\lambda(\mathrm{own},w)$. In particular it holds for the tuples required in \cref{pf:rows:defsharp}(b); the
conditions on $\mu$ there (used, $\mu\ne\mu_0$, $Q_p(\mu)$) are not used.
\end{proof}

\paragraph{Built data and the five laws.}
Let $D$ be a datum and $u$ a word. \emph{Built data} $\mathcal B$ for $(D,u)$ consist of a CM field $E'$ with embeddings
$\mathrm{emb}_\tau:E'\to\mathit{LL}$ ($\tau\in\mathrm T$) compatible with the fixed embedding of $\mathit{LL}$, and of chain data
over $E'$: for every level $i$ a finite collection of \emph{decorated} hermitian $\OO_{E'}$-lattices $d$ in $E'^{N_i}$, organised in
blocks $b$ with groups $G_b$ acting on the decorations, with representations $\rho_b$ on the spaces $V_i$ of $\mathit{LL}$-valued
algebraic vectors of weight $\lambda_i$ (the weight of the word at level $i$); data at a prime $T=T(\mathcal B)$ with a set
$S_T$ of primes of $E'$ above $T$ and, at every $P\in S_T$, up-, down- and diamond operators on the decorations; period data;
member-form data at the places of $\mathit{LL}$; and the laws relating them. Built data determine a chain
$\mathrm{ch}(\mathcal B)$ for $(D,u)$. A \emph{kernel} at level $i<m$ is a family of $\mathit{LL}$-linear maps
$k(d,d'):V_i\to V_{i+1}$ indexed by decorations $d$ of level $i$ and $d'$ of level $i+1$ (in blocks $b,b'$). The five
\emph{laws} of a kernel are:
\begin{enumerate}[label=(K\arabic*)]
\item (equivariance) $k(gd,g'd')=\rho_{b'}(g')\,k(d,d')\,\rho_b(g)^{-1}$ for $g\in G_b$, $g'\in G_{b'}$, and the $K$-maps of the
chain are given on evaluations by $k$ at representatives of the classes;
\item (dilation) $k(\langle a\rangle_Pd,\langle a\rangle_Pd')=k(d,d')$ for the diamond operators at $P\in S_T$;
\item (Hecke relation) for $P\in S_T$: $\sum_{d_1}k(d,d_1)=T^{(t+1)/2}\sum_{d_0}k(d_0,d')$, where $d_1$ runs over the
decorations obtained from $d'$ by an up-operator at $P$ and $d_0$ over those from which $d$ is so obtained;
\item (conjugate primes) at every place $w$ above $T$ not seen by the member forms, $k(d,d')$ maps the localisation at $w$ of the
lattice of $d$ into $T^{f(t+1)/2}$ times the localisation of the lattice of $d'$;
\item (level primes) at every seen place $w$ above $T$, $T^{fN_i+\sum_\tau s_{\mathrm{div}}(i,\tau)}k(d,d')$ maps the localisation
at $w$ of the lattice of $d$ into the localisation of the lattice of the decoration of $d'$.
\end{enumerate}
A \emph{split} of the data at level $i$ is an equivalence relation on the decorations (``same lattice away from $S_T$'')
compatible with the group actions and the operators, which separates decorations with the same decoration data, together with
the conditions $T\in P$, $T\notin P^2$ for $P\in S_T$, and that no two primes of $S_T$ are conjugate. The data \emph{see the
primes} if, for every place $w$ above $T$ and $P\in S_T$, an element $x\in E'$ satisfies $v_P(x)\ge0$ if and only if its image
under the embedding attached to $(w,P)$ (applied to $x$ or to $x^*$, according to whether $w$ is seen) is integral at $w$.

\begin{proposition}[The last step]\label{pf:rows:laststep}
Let $(B,X,I,l,l_0,Z)$ be true column data satisfying the sign conditions of $B$, with datum $D$ and twisting forms $\mathrm{Tw}$.
For every word $u$ and every tuple $a$ of $I$, the twisted chain $\mathrm{tw}(\mathrm{Tw},u,a)$ satisfies the sharp bound.
\formallabel{GlobalTheory.LvT.Built.TrueCols.trueColumns_chain_last_step_sharp_bound}
\end{proposition}

\begin{proof}
By \cref{pf:rows:fivelaws} there are built data $\mathcal B$ for $(D,u)$ with $\mathrm{ch}(\mathcal B)=\mathrm{tw}(\mathrm{Tw},u,a)$,
splits at all levels $i\le m$, which see the primes, and with kernels satisfying the five laws at all levels $i<m$. The field of
$\mathcal B$ is the host field $E$, its prime is $T_u$, and $[E:\Q]N_m+1<T_u$: indeed $[E:\Q]=2f$, and the pre-stage $X$ has a
parameter $a_X$ with $2fN_m+1<a_X<T_u$. The weight signs of every chain of the assembled datum are equivalent to the sign conditions of $B$
(\lref{GlobalTheory.LvT.WtSign.wtSigns_assemble_iff}). The stability theorem for lattice column chains
(\lref{GlobalTheory.LvT.Built.Assemble.sharpBound_stability_latticeColumnChain}: built data with splits, seeing the primes,
with kernels satisfying the five laws at all levels, with $[E':\Q]N_m+1<T$ and with the weight signs, have a chain satisfying the
sharp bound and with stable forms) gives the assertion. (That theorem realises every level by a point model whose value
lattices are Weyl lattices, transfers the five laws from the coefficient kernels to point kernels, and deduces the stability of
the forms, the commutation of the chain with the $U$-operators and the sharp bound.)
\end{proof}

\begin{proposition}[The five laws for true chains]\label{pf:rows:fivelaws}
Let $(B,X,I,l,l_0,Z)$ be true column data, with datum $D$ and twisting forms $\mathrm{Tw}$. For every word $u$ and tuple $a$ of $I$,
the built data $\mathcal B=\mathcal B(I,l,Z,u,a)$ of the true chain satisfy $\mathrm{ch}(\mathcal B)=\mathrm{tw}(\mathrm{Tw},u,a)$,
admit splits at all levels $i\le m$, see the primes, and at every level $i<m$ admit a kernel satisfying the five laws.
\formallabel{GlobalTheory.T.fivePointLaws_trueChains}
\formallabel{GlobalTheory.LvT.Built.TrueLaws.laws_true_fam}
\end{proposition}

\begin{proof}
Since $I$ is ok, the column inclusion provides theta data: a curve ideal $\mathfrak c$, a twelfth root $\delta$ of the
discriminant of the lattice $\mathfrak c$ and a character $\chi$ (\lref{GlobalTheory.KerT.hyp_of_colIn}). For these, the
normalised theta family of the column choice satisfies the equivariance laws
(\lref{GlobalTheory.T.normalisedFamily_satisfies_equivarianceLaws}) and its vectors lie in the theta lattices
(\lref{GlobalTheory.T.normalisedThetaFamily_vectors_lattices_thetaZero}); hence the chain theta data are ok at every tuple and
word (\lref{GlobalTheory.ChainTheta.chTh_ok_of_colIn'}). We show that the true kernel $k_i$ of level $i<m$ satisfies the five
laws. Now:
\begin{itemize}
\item $\mathrm{ch}(\mathcal B)=\mathrm{tw}(\mathrm{Tw},u,a)$ holds by construction;
\item splits exist at every level (\lref{GlobalTheory.LvT.Built.TrueLaws.split}, from the existence of split points for
the selected primes), and $\mathcal B$ sees the primes by the laws of the column choice
(\lref{GlobalTheory.LvT.Built.TrueLaws.sees});
\item (K1): the equivariance of $k_i$ follows from the matrix and coupled laws of the chain theta data
(\lref{GlobalTheory.LvT.Built.TrueLaws.kerEq}), and the $K$-maps of the chain are given by $k_i$ at representatives
(\lref{GlobalTheory.T.trueChainKernel_eq_pointKernel_at_representatives}, using that the vectors lie in the lattices);
\item (K2): $k_i$ is invariant under the diamond operators (\lref{GlobalTheory.LvT.Built.TrueLaws.kerDil});
\item (K3): the Hecke relation follows from the move step of the theta family
(\lref{GlobalTheory.LvT.Built.TrueLaws.kernelCondition_at_row_family});
\item (K4) and (K5): the bounds at the conjugate primes and at the level primes
(\lref{GlobalTheory.T.trueChainKernels_bound_conjugatePrimePlaces} and \lref{GlobalTheory.T.trueChainKernels_bound_levelPrimePlaces})
are deduced from the integrality of the moments of kernel jets, \cref{pf:rows:jetmoments}, and from the simplicity of the
induced modules of $\GL_N$ in characteristic $p$ whose highest weight $\lambda$ satisfies $\lambda_i-\lambda_j+(j-i)\le p$ for
$i<j$ (closed lowest alcove; \cite[II, \S5]{Jantzen2003}, \lref{GlobalTheory.T.weylModule_simple_closedLowestAlcove}).
\qedhere
\end{itemize}
\end{proof}

\paragraph{CM theta data and $p$-adic moments.}
Let $D$ be a datum. Its host contains $K=\Q(\sqrt{-d})$, with $d$ prime and $d\equiv7\pmod{16}$
(\lref{GlobalTheory.T.cmBaseField_integers_and_curveIdeal_facts}); let $\varphi_0:E\to\C$ be the fixed embedding and
$\tau\in\mathrm T$ the embeddings $E\to\C$ of the CM type. \emph{Theta data} $(\mathfrak c,\delta,B,\chi)$ consist of a non-zero ideal
$\mathfrak c$ of $\OO_K$ with $\Nm\mathfrak c$ prime to $2d$; $\delta\in\C$ with $\delta^{12}=\Delta(\varphi_0(\mathfrak c))$, so that
$\Lambda=\delta\varphi_0(\mathfrak c)$ is a lattice with discriminant $1$; a \emph{base} $B$ (a lattice structure on $\Lambda$ with a
marked $2$-torsion point and an identification of the $K$-structures); and a function $\chi$ on the fractional ideals of
$\OO_E$, multiplicative on the ideals $\mathfrak a$ with $\mathfrak a\bar{\mathfrak a}=\OO_E$, with
$\chi(z\OO_E)=\prod_\tau\tau(z)^{-1}$ for $z\in E$ with $zz^*=1$. For even $N,N'$ and self-dual hermitian lattices
$L\subseteq E^N$, $L'\subseteq E^{N'}$ (an \emph{admitted pair}) let $V=(M_{N'\times N}(\C))^{\mathrm T}$ with the hermitian form
$H(z,w)=\sum z\bar w/\mathrm{area}(\Lambda)$, let $\iota_{\mathrm T}(A)=(\tau(A))_\tau$ for $A\in M_{N'\times N}(E)$, and let
$\mathrm{La}(L,L')\subseteq V$ be the group generated by the $w\,\iota_{\mathrm T}(A)$ with $w\in\Lambda$ and $A$ in the
$\OO_E$-module $Y(L,L')$ spanned by $aA$, $a$ in a fixed ideal, $A\in\Hom(L,L')$. The theta data define the \emph{normalised
theta family} $e_\chi(N,N',L,L'):V\to\C$ (a $\chi$-factor times the translated theta function of $B$). On functions $G:V\to\C$
we use the translation operators $U_xG(z)=G(z+x)\exp(-\pi H(z,x)-\pi H(x,x)/2)$, the Weil operators $W_xG=w_p(x)U_xG$ with a factor
$w_p(x)$ (a sign if $p$ is odd), and the phases $\kappa(x,y)=\exp(-2\pi i\Im H(x,y))$.

Let $p$ be a prime and $\lambda:\Qbar\to\C_p$ an embedding. For $G$ with algebraic Taylor coefficients at $0$ let $\lambda G$ be the
power series with coefficients $\lambda$ of these, and for $y'_1,\dots,y'_k\in M_{N'\times N}(E)$ let $\mathrm{jet}(y';\lambda G)$ be the
constant term of the iterated derivative of $\lambda G$ in the directions $(\lambda\tau(y'_j))$. For $\Omega\in\C_p^\times$ and
$b=(b_1,\dots,b_r)$ put $\mathrm{pair}_\Omega(b,y)(a)=\Omega^{-1}\sum_ia_i\lambda\varphi_0(\Tr_{E/K}\langle y,b_i\rangle)$ for
$a\in\Z_p^r$.

\begin{definition}[The moment property]\label{pf:rows:defmoments}
The theta data have the \emph{moment property} at $(p,\lambda,\Omega)$ if for every admitted pair $(L,L')$ and every $t\in V$
with $nt\in\mathrm{La}(L,L')$ for some $n\ge1$ prime to $p$, there is $N_0\ge1$ prime to $p$ such that for every family
$b=(b_1,\dots,b_r)$ in $M_{N'\times N}(E)$ which is a basis of $Y(L,L')$ locally at $p$ (linearly independent over $K$, with
$sb_i\in Y(L,L')$ and $sy\in\sum_i\OO_Kb_i$ for every $y\in Y(L,L')$, for suitable integers $s$ prime to $p$) and satisfies
$\delta\iota_{\mathrm T}(b_i)\in\mathrm{La}(L,L')$, there is
a continuous linear form $\mu$ on $C(\Z_p^r,\C_p)$ with $\|\mu\|\le1$ such that for every $n\ge0$ and every $y$ in the group
generated by the $\frac{\alpha N_0}{p^n}x$ with $x\in\mathrm{La}(L,L')$ and $\alpha\in\OO_K$, $|\lambda(\bar\alpha)|\le p^{-2n}$:
\begin{enumerate}[label=(\alph*)]
\item $W_yU_te_\chi(N,N',L,L')$ has algebraic Taylor coefficients;
\item for all $y'_1,\dots,y'_k$ and every continuous $g:\Z_p^r\to\C_p$ such that for all $a\in\N^r$
\[
g(a)=\lambda\bigl(\kappa(\delta\iota_{\mathrm T}(\textstyle\sum_ia_ib_i),y)\bigr)\prod_j\mathrm{pair}_\Omega(b,y'_j)(a),
\]
one has $\mathrm{jet}\bigl(y';\lambda\,W_yU_te_\chi(N,N',L,L')\bigr)=\mu(g)$.
\end{enumerate}
\end{definition}

\begin{theorem}[Moments of the translates of the unit periods]\label{pf:rows:unitperiods}
Assume that (i) for all theta data the normalised theta family satisfies the Euler product identity of the base, and (ii) for
all theta data, admitted pairs $(L,L')$ and torsion points $t$ of $V/\mathrm{La}(L,L')$, the Taylor coefficients of $U_te_\chi$ are
algebraic. Then for all theta data, every prime $p\ne2$ which splits in $K$ and is prime to $\Nm\mathfrak c$, and every embedding
$\lambda:\Qbar\to\C_p$, there is $\Omega\in\C_p$ with $|\Omega|=1$ such that the moment property holds at $(p,\lambda,\Omega)$.
\formallabel{GlobalTheory.T.unitPeriodTranslates_moments_finite}
\formallabel{GlobalTheory.MomFin.moments_fin}
\end{theorem}

\begin{proof}
\emph{The period.} Let $\iota:\C\to\C_p$ be a field isomorphism extending $\lambda$
(\lref{GlobalTheory.PadicComplexIso.exists_ringEquiv_extending}). For the base field $K$ and the split prime $p$ there are
$\delta_0\in\C$ with $\delta_0^{12}$ the discriminant of $\OO_K$, a theta function $\theta$ of the CM curve and a $p$-adic period
$\Omega_0$ with $|\Omega_0|=1$, satisfying the theta laws at all isomorphisms $\C\to\C_p$
(\lref{GlobalTheory.MomN.exists_per}: the unit period of the multiplicative chart of the formal group at the split prime).
The ratio $\delta/\delta_0$ is algebraic and an $\iota$-adic unit (\lref{GlobalTheory.T.discOneScale_ratio_ideal_ring_algebraic_unit}),
since $\mathfrak c$ is prime to $p$. Put $\Omega=\Omega_0\cdot\Nm\mathfrak c\cdot\lambda(\delta/\delta_0)$; then $|\Omega|=1$.

\emph{Analytic setup.} Fix an admitted pair and $t$ of order prime to $p$, with $N_0t\in\mathrm{La}(L,L')$, $N_0$ prime to $p$.
There are $g$, a setup $S$ of dimension $g$ (a lattice $\mathrm{La}_S\subseteq\C^g$, a hermitian form $H_S$ and the integer $d$)
and a linear isomorphism $e:V\to\C^g$ carrying $\mathrm{La}(L,L')$ to $\mathrm{La}_S$ and $H$ to $H_S$, under
which the pairing of $\delta\iota_{\mathrm T}(x),\delta\iota_{\mathrm T}(y)$ becomes $\varphi_0(\Tr_{E/K}\langle x,y\rangle)/\Nm\mathfrak c$
and local bases at $p$ go to local bases (\lref{GlobalTheory.T.exists_analyticSetup_periodLatticePair}). In these
coordinates $F=e_\chi(N,N',L,L')$ is a theta function for $S$ (\lref{GlobalTheory.T.normalizedThetaVector_isometry_isTheta}),
and since its Taylor coefficients are algebraic by (ii) (with $t=0$), its jets in the algebraic directions are algebraic
(\lref{GlobalTheory.T.algebraicJets_of_algebraic_taylor}). The operators $U_x,W_x$ satisfy the Heisenberg relations
$U_xU_y=e^{-\pi i\Im H(x,y)}U_{x+y}$, $U_xU_y=\kappa(x,y)U_yU_x$, and $W_xW_y=W_{x+y}$ when the pairing of $x,y$ is integral
(\lref{GlobalTheory.HeisOps.heisenberg}), and they preserve smoothness.

\emph{The measure.} For the theta function $F\circ e^{-1}$ with algebraic jets and the prime-to-$p$ point $e(t)$, the moments of
the translates define a measure: there is a continuous linear form $\mu$ on $C(\Z_p^r,\C_p)$ with the moment property for the
translate $U_{e(t)}(F\circ e^{-1})$ (\lref{GlobalTheory.MomOne.Translates.moments_of_translates}, using the theta laws of
$\theta$ and $\Omega_0$). Its norm is at most $1$ if and only if the translate has integral lattice jets
(\lref{GlobalTheory.T.momentMeasure_norm_le_one_iff_latticeIntegral}); this integrality holds for $U_te_\chi$ on
$\mathrm{La}(L,L')$ by (i) (\lref{GlobalTheory.T.normalizedVector_translates_prime_order_integral}: translates of the normalised
vector by points of order prime to $p$ are integral) and transfers to $\mathrm{La}_S$
(\lref{GlobalTheory.T.latticeJetsAreBoundedMoments_pullback}). Hence $\|\mu\|\le1$.

\emph{The moments.} Let $n\ge0$ and $y$ as in \cref{pf:rows:defmoments}. Its image $e(y)$ is a good lift for $S$, and there are a
level $m\ge2n$, an element $\eps$ with the idempotence property of level $m$, and $y_t$ such that
$W_{e(y)}U_{e(t)}(F\circ e^{-1})=W_{\mathrm{lift}(y_t)}U_{e(t)}(F\circ e^{-1})$
and the phases $\kappa$ of $e(y)$ and of the lift agree on the $\N$-combinations of the $b_i$
(\lref{GlobalTheory.T.goodLift_eq_idempotent_lift}). Moreover $p^ny\in\mathrm{La}(L,L')$, the pairing of $y$ with itself is
integral, and the values $\Im H$ on $\mathrm{La}(L,L')$ are integers; hence the Taylor coefficients of $W_yU_tF$ are algebraic
(\lref{GlobalTheory.T.coefIn_weilOperator_translate_algebraic}), which is (a). For (b), the function $g$ agrees on $\N^r$
with the function prescribed by the moment property in the coordinates of $S$ (here the pairing identity of the setup is used,
\lref{GlobalTheory.T.pairing_scaledDirections_eq_period_pairing}), hence everywhere on $\Z_p^r$, by continuity and density
(\lref{GlobalTheory.T.eq_on_padicInt_of_eq_on_naturals_moment}); the directions $\delta_0\delta^{-1}e(\delta\iota_{\mathrm T}(y'_j))$
are algebraic (\lref{GlobalTheory.T.algebraicDirection_scaled_embedding}). The moment property of $\mu$ then gives
$\mathrm{jet}=\mu(g)$ in the coordinates of $S$, and transporting the jets back along $e$
(\lref{GlobalTheory.T.iteratedFDeriv_algebraic_directions_jet_transport}) gives (b).
\end{proof}

\paragraph{Jet sums at the level places.}
For the jet tables of the kernels we use the ambient theta family $A=A(s,B,\chi,c_D)$ attached to a seam $s$ from a datum $H$ to the
fields $(F_X,E_X,K_X)$ of a jet datum, theta data $(\mathfrak c,\delta,B,\chi)$ of $H$ and constants $c_D$; these are \emph{analytically
admissible} if the seam is ok with the reference automorphism of the jet datum, $\mathfrak c,\delta,\chi$ are as above, and $c_D$ is
generic. For a word $u$ with level prime $T=T_u$, a selection, a set $S$ of primes containing the host primes, a tuple and local
constants with $\Nm\mathfrak c$ prime to the row primes, a level $i<m$, decorated self-dual lattices $d,d'$ of ranks $N_i,N_{i+1}$
which are standard at $T$, and a function $\Xi_0$ on $\Hom_E(E^{N_i},E^{N_{i+1}})$ which is integral at $w$ and periodic on the
multiplied lattice (an \emph{off-datum}), let $\mathcal J(m)$ be the sum, over the divided lattice of homomorphisms modulo the
multiplied lattice, of $\Xi_0(x)\phi_T(x)$ times the jet coefficient $\mathrm{jet}_K(x,m)$ of $A$, for monomials $m$ in the coordinates
$(\tau,i,j)$; call $m$ \emph{rational} if for every $x$ in the support of $\Xi_0\phi_T$ the jet coefficient at $m$ of the induced theta
vector of $x$ is a $\chi$-factor times an element of $\sigma_0(K_X)$.

\begin{definition}[Integral jet moments]\label{pf:rows:defjm}
Let $\iota:K_X\to\C_T$ induce the place $w$. The \emph{jet moment statement} at $(u,w,\iota)$ asserts: there are
$\upsilon_\tau\in\C_T$ ($\tau\in\mathrm T$), $|\upsilon_\tau|=1$, such that for all data as above, with
$\mathrm{pt}(x,m)=\prod_{(\tau,i,j)}\bigl(\upsilon_\tau\,\iota\tau(x^\dagger(e_i)_j)\bigr)^{m_{\tau ij}}$ ($x^\dagger$ the adjoint of $x$):
\begin{enumerate}[label=(\alph*)]
\item if $w$ is seen-bar for $u$, then for every finite set $\mathcal M$ of rational monomials and every $\eps>0$ there are
$c_x\in\C_T$, $|c_x|\le1$, finitely many non-zero and supported on the set of $x$ mapping the lattices of $d$ at the primes
above $T$ into those of $d'$ compatibly with the decorations, such that $|\iota\,\mathcal J(m)-\sum_xc_x\mathrm{pt}(x,m)|<\eps$
for $m\in\mathcal M$;
\item if $w$ is seen for $u$, the same holds with $\mathrm{pt}(x,m)$ replaced by $T^{-fN_i}\mathrm{pt}(x,m)$ and the support
condition replaced by: $x^\dagger$ maps the vectors of $d'$ with vanishing decoration into the lattice of $d$.
\end{enumerate}
The \emph{jet moments} statement $\sfname{JM}$ asserts this for every ok jet datum, every analytically admissible
$(H,s,\mathfrak c,\delta,B,\chi,c_D)$, every word $u$ with $T_u$ prime and prime to $\Nm\mathfrak c$, and every place $w$ of $K_X$
above $T_u$ and embedding $\iota$ inducing $w$.
\end{definition}

\begin{theorem}[Integral moments of kernel jets]\label{pf:rows:jetmoments}
$\sfname{JM}$ holds.
\formallabel{GlobalTheory.T.kernelJets_levelPlace_integralMoments}
\formallabel{GlobalTheory.MomTop.jetMoments}
\formallabel{GlobalTheory.MomT.jetMoments_of_unitPeriod}
\formallabel{GlobalTheory.MomLat.latticeFactsAll}
\end{theorem}

\begin{proof}
The two hypotheses of \cref{pf:rows:unitperiods} hold (\lref{GlobalTheory.MomEq.eq_and_alg}: the Euler product identity of
the normalised theta family for curve ideals, and the algebraicity of the Taylor coefficients of its torsion translates), so
\cref{pf:rows:unitperiods} applies. We also use the lattice facts for the ambient family: its lattices and vectors are the
transports along the seam of the lattices $\mathrm{La}(L,L')$ and of the normalised theta family (for admitted pairs), and the
family is smooth (\lref{GlobalTheory.MomLat.latticeFactsAll}, from \lref{GlobalTheory.PeriodTransport.seam_lattices}).

Fix the data of \cref{pf:rows:defjm}. The embedding $\iota$ and the reference automorphism determine $\lambda:\Qbar\to\C_T$
(\lref{GlobalTheory.MomPlace.places}). The prime $T$ is not one of the small primes of the host, so $T\ne2$; and $T$ splits in
$K$, since $T=\pi_T\bar\pi_T$ with $\pi_T$ prime to $\bar\pi_T$ (\lref{GlobalTheory.MomT.SplitK.splitk}). By
\cref{pf:rows:unitperiods} at $p=T$ there is $\Omega$, $|\Omega|=1$, with the moment property. The constants $c_D(\tau)$ are
algebraic, and $\upsilon_\tau=\lambda(c_D(s(\tau)))/\Omega$ has absolute value $1$
(\lref{GlobalTheory.EigenGenerator.norm_generator_eq_one}: the genericity generators have $\lambda$-adic absolute value $1$
at primes which are neither small nor equal to $q$).

Now fix the selection data, $i$, $d$, $d'$ and $\Xi_0$. The ranks $N_i$ and $N_{i+1}$ are even. The ambient lattice ideal is a unit at
the primes above $T$ and at the row primes, so the data form a \emph{context} for the ambient family at $(\pi_T,T)$
(\lref{GlobalTheory.Proof.AbsNormSplitPlace.main}, \lref{GlobalTheory.ThetaMove.BadPrimes.bad_primes}, and \lref{GlobalTheory.PeriodIdealUnit.locP_AA_eq_one}),
and the $\chi$-factors of $d,d'$ are algebraic with $|\lambda(\cdot)|=1$
(\lref{GlobalTheory.MomT.ChiUnit.chi_det_unit}). The moment property yields point moments at $\bar\pi_T$ and at $\pi_T$: the
jets of the translates are moments of measures of norm $\le1$ against the point table
(\lref{GlobalTheory.MomT.PtMoments.ptMom}).
\begin{itemize}
\item If $w$ is seen-bar, (a) follows from the point moments at $\bar\pi_T$ (\lref{GlobalTheory.MomT.type_bar}: in a
context, point moments at the conjugate prime give the bar case of the jet moment statement).
\item If $w$ is seen, let $N''$ be the companion lattice of $d'$ at $\pi_T$ (\lref{GlobalTheory.MomT.MomentCompanion.companion};
$\pi_T$ is not divisible by the square of a prime above $T$). Then $\chi$ of the determinant ideal of $(d',N'')$ equals
$\sigma_0(k)$ with $k=\prod_\tau\tau(\pi_T/\bar\pi_T)^{-1}$ (\lref{GlobalTheory.MomT.ChiMove.chi_move}), and
$|\iota k|\le T^f$. The point moments at $\pi_T$, the movement of the ambient family along $N''$
(\lref{GlobalTheory.MomT.Move.move}, which uses the Euler product identity) and $|\lambda\sigma_0(k^{N_i/2})|\le T^{fN_i}$
give (b) (\lref{GlobalTheory.MomT.type_p}: the corresponding statement for the seen case).\qedhere
\end{itemize}
\end{proof}

\paragraph{The row.}

\begin{proposition}[Level-prime valuations at bundle points]\label{pf:rows:levelrow}
Let $Q$ be a slot over columns. For every bundle type and linkage, the row $\mathsf{aP10}$ holds: $\rowsRow(Q,\rowsLk,\rowsEvC(R_6))$,
where $R_6(D,O)$ asserts, for every word $u$, the sharp bound of $D$ at $u$ and the twisted column bound at $T$ for the twisting
forms and the division datum of $O$.
\formallabel{GlobalTheory.T.levelPrimeColumnValuations_bundlePoints_overColumns}
\formallabel{GlobalTheory.Fc.Cv.row_level_valuations}
\end{proposition}

\begin{proof}
By \cref{pf:rows:calculus}(d),(e) it suffices that $R_6$ holds at every bundle point $\beta\models Q$ of every context. By
\cref{pf:rows:view} it suffices that $R_6(D(w,Z),O)$ holds whenever $w$ is a column choice, $Z$ a later tuple, $O$ an object over
$D(w,Z)$, and the column view, the sign conditions of the base and the real ties hold. By the column view, the twisting forms of
$O$ are those of $(w,Z)$, and the column choice is ok with jets in range; so $R_6(D(w,Z),O)$ is \cref{pf:rows:closingsharp} for the
true column data of $w$ and the division datum of $O$.
\end{proof}

% ---------------------------------------------------------------------------
\subsection{The definite rows and the arithmetic row at generator places}\label{sec:pf:rows:rrows}

\paragraph{Definite rows.}
For a realisation $Q$ of a bundle (the indefinite real realisation of its face), constants $C,C_H$ and propositions
$h_N,h_O$, the \emph{constant bounds} are the nine statements: every place at which the first defect condition of $Q$
holds is a twisted place; the covolume bound $\sfname{L3}$; $\sfname{ON}$ if $h_N$; $\sfname{L6}$ if $h_O$; the isometry
condition; the multiplicity condition; $\sfname{FUN}$ if $h_N$; the pass condition; and the $\kappa$-rows. The row
$\mathsf{rRows}$ for a definite-side predicate $\mathrm{DR}$ requires: closedly on $\rowsDiii(\mathcal A)$, there is $C\ge0$
such that eventually along the selection the constant bounds hold for the realisation of $\sigma(\mathcal F,x,i)$ with
$(C,C_H)$, $h_N=\sfname{DIVinN}(\mathcal F)$ and $h_O$ the leaf uniqueness of the selected bundle, together with
$\mathrm{DR}(\sigma(\mathcal F,x,i),C,C_H)$.

A \emph{compatible selection} $\tau$ for a raw selection $\sigma$ chooses a comparison for every selected bundle; it
yields a selection of \emph{point summands} $\sigma_\tau$ (the selected bundle with the chosen comparison and its point).

\begin{lemma}[Assembly of the definite rows]\label{pf:rows:rrowsparts}
Let $\mathrm{DR}$ be a definite-side predicate, $Q$ a slot and $\rowsLk$ a linkage of bundles. Assume the rows asserting:
for every selection $\sigma$ there is a compatible selection $\tau$ such that, closedly on $\rowsWiii(\mathcal A)$ and for all
$C$ beyond some $C_i\ge0$, eventually along $\sigma_\tau$: (i) the small point conditions (twisted places, isometry, pass);
(ii) $\sfname{L3}$; (iii) $\sfname{ON}$ for members satisfying $\sfname{DIVinN}$; (iv) $\sfname{L6}$; (v) $\sfname{FUN}$ for members
satisfying $\sfname{DIVinN}$. Assume moreover the rows $\mathsf{mi}$, $\mathsf{kap}$ (eventually, the multiplicity condition and
the $\kappa$-rows) and $\mathsf{df}$ (for some $C_0\ge0$ and all $C\ge C_0$, eventually $\mathrm{DR}$). Then $\mathsf{rRows}$ holds
for $Q$ and $\rowsLk$.
\formallabel{GlobalTheory.Fc.row_rRows_of_parts_alt}
\end{lemma}

\begin{proof}
Combine the four rows into one (\lref{GlobalTheory.Fc.Row.and}). Given a selection, the eight closed statements on
$\rowsWiii(\mathcal A)$ combine into one, which restricts to $\rowsDiii(\mathcal A)\subseteq\rowsWiii(\mathcal A)$. For a member,
take $C$ the maximum of the six constants $C_i$ and $x_0$ the maximum of the eight thresholds; beyond $x_0$ the nine fields
of the constant bounds are read off from (i)--(v), $\mathsf{mi}$ and $\mathsf{kap}$, and $\mathrm{DR}$ from $\mathsf{df}$.
\end{proof}

\paragraph{Placement of realisations.}
The realisation $Q$ of a bundle comes with a Hida setting $X$ (a unitary group over the host field with a level subgroup
$K_2$ and a level character $\chi_2$, and a space $\mathcal E(X)$ of good eigenforms), an isomorphism $\iota_S$ between the adelic
group of the hermitian space $Y_0$ of the datum and that of $X$, an element $t_S$ of the adelic group of $Y_0$ (the
translation element of $Q$), a discrete series model $M$ for $Y_0$ with a measure and a fundamental set, and the space $E_Q$
of row maps (maps from the model to automorphic forms) at level $0$, with level group $J_Q$ and level character $\chi_Q$.
We say that $Q$ \emph{extends} if
\begin{enumerate}[label=(\roman*)]
\item there is $t$ in the standard open compact subgroup of the finite adelic group of $Y_0$ such that conjugation by $t$
carries $J_Q$ onto $\iota_S^{-1}(\{x:x_\infty=1,\ x_{\mathrm f}\in K_2\})$ and $\chi_Q$ to $\overline{\chi_2}\circ\iota_S$
(i.e.\ $k\in J_Q\Leftrightarrow t^{-1}kt\in\iota_S^{-1}(\dots)$, and then $\overline{\chi_2}(\iota_S(t^{-1}kt))=\chi_Q(k)$), and the
row levels match at $t$ (the three matching conditions between the level data of the row and of $X$, through an
isomorphism of the finite type groups and a reduction map); and
\item for some $c'\in\R$ the \emph{evaluation tie} holds at $t_S$: there are a vector $v_0$ of the model with components of norm
$1$ and a linear isomorphism $j:E_Q\to\mathcal E(X)$ such that, for every $T\in E_Q$, $j(T)(x)=\mathrm{ev}(T,v_0)(\iota_S^{-1}(x)\,t_S)$
for almost all $x$, and $\|j(T)\|^2=c'\,\Re\langle T,T\rangle_M$.
\end{enumerate}

\begin{proposition}[Placement]\label{pf:rows:placement}
Assume the classification of discrete series, (5) of \cref{pf:rows:arch}. Let $Q$ be a realisation as above which is
normalised (a list of conditions fixing the coefficient data, the phases, the group records and the translation element $t_S$;
among them a pin of the group records, $t_S$ lies in the standard open compact subgroup, and (i) below implies the same
conjugation and matching at $t_S$) and for which the embedding of the host field of its Hida setting is the identity. Then $Q$
extends.
\formallabel{GlobalTheory.T.indefinite_scale_theta_nonvanishing_alt}
\formallabel{GlobalTheory.RowPlacementReal.placement_of_realisation_alt}
\end{proposition}

\begin{proof}
\emph{Step 1: the idempotent of the type data has real coefficients.} For both type primes $g_k$, the Eisenstein series
attached by the type data to a monic polynomial $f$ over $\F_{g_k}$ with $f(0)\ne0$ satisfy: the complex conjugates of the
coefficients of the series of $f$ are the coefficients of the series of the reciprocal polynomial $f(0)^{-1}X^{\deg f}f(1/X)$
(\lref{GlobalTheory.T.star_coeff_eisenstein_series_reverse_poly}, applied to the equivariant label maps of the type
data). This implies that the coefficients of the idempotent $e_0$ of the type data are real
(\lref{GlobalTheory.GoodIdemReal.star_coeff_of_reverse}).

\emph{Step 2: a placement.} From the level, group and model data of the realisation and the host hypothesis there is $t$ in the
open compact subgroup of $Y_0$ which conjugates the level of the row to that of $X$ as in (i), together with a reordering of
frames (\lref{GlobalTheory.T.exists_finite_element_conj_rowLevel_character}); and for every such $t$ the row levels match
at $t$ (\lref{GlobalTheory.T.rowLevelMatchAt_of_levelConjugation}), by the pin of the group records contained in the
normalisation, the identity embedding of the host field, and Step 1. This is (i).

\emph{Step 3: the evaluation tie.} The normalisation contains the statement that (i) implies the same conjugation and matching
at $t_S$, and $t_S$ lies in the open compact subgroup. At an element of the open compact subgroup at which the levels are
conjugate and match, the good rows compare by evaluation: the evaluation tie holds for some $c'$
(\lref{GlobalTheory.T.goodRows_compare_by_levelMatching_evaluation}, which uses the classification of discrete series).
This is (ii).
\end{proof}

\begin{theorem}[Definite rows at stage zero]\label{pf:rows:rrows}
Assume the five statements of \cref{pf:rows:arch}, $\sfname{SurjectiveOntoX}$ and $\sfname{LambdaLowerBound}_2$. Let $Q$ be a
slot and $\rowsLk$ a linkage of bundles for which the rows $\mathsf{hPA}$, $\mathsf{hDivG}$, $\mathsf{good}$, $\mathsf{geo}$, the
covolume row, the functional row, $\mathsf{mi}$, $\mathsf{kap}$ and $\mathsf{df}$ (for $\sfname{DefiniteSideCondition}$) hold. Then
$\mathsf{rRows}$ holds for $Q$ and $\rowsLk$ with $\mathrm{DR}=\sfname{DefiniteSideCondition}$.
\formallabel{GlobalTheory.Fc.StageZeroRow.row_rRows_stage0_alt}
\end{theorem}

\begin{proof}
By \cref{pf:rows:rrowsparts} it suffices to produce, from the conjunction of the rows $\mathsf{hPA}$, $\mathsf{hDivG}$,
$\mathsf{good}$, $\mathsf{geo}$, covolume and functional, the statements (i)--(v). Let $\sigma$ be a selection over $\mathcal A$
for this conjunction. The row $\mathsf{geo}$ provides, at every selected bundle, a comparison whose point summands have the
iota-point property, a class-sum pin, the compatibility conditions and the ``more'' property; let $\tau$ be the
compatible selection choosing them. The members of $\mathcal A$ satisfy the seven duties
(\lref{GlobalTheory.EX.DQ.familyClass_sevenDuties_and_stableLattices}). The row $\mathsf{hDivG}$ gives the division
conditions of the point summands (\lref{GlobalTheory.Fc.StageZeroRow.p_div_at_piece_of_div3g}), and $\mathsf{hPA}$ their
primitivity. Then:
\begin{itemize}
\item by (1)--(4) of \cref{pf:rows:arch}, primitivity, the division conditions, goodness and the class-sum pins give,
eventually, the row conditions of the point summands (\lref{GlobalTheory.Fc.Ind.evW_pRows_of_rows_alt});
\item $\sfname{ON}$ follows from $\sfname{SurjectiveOntoX}$, the row conditions and the ``more'' property
(\lref{GlobalTheory.T.eventually_norm_bound_of_more_points});
\item $\sfname{L6}$ follows from $\sfname{LambdaLowerBound}_2$ and the row conditions
(\lref{GlobalTheory.T.eventually_L6_bound_of_lamLow});
\item the small point conditions follow from primitivity and the fact that every normalised realisation whose host
embedding is the identity extends (\cref{pf:rows:placement}), by
\lref{GlobalTheory.T.eventually_small_of_seeEx_primitive};
\item $\sfname{L3}$ and $\sfname{FUN}$ are the covolume and functional rows.\qedhere
\end{itemize}
\end{proof}

\paragraph{The arithmetic row at the generator places.}
The requirement of $\mathsf{aAtG}$ is: closedly on $\rowsWiii(\mathcal A)$ there is $C_1$ such that, eventually, at the places
$w$ above the two type primes and away from the level prime of a word, the period values of the used tuples are images of
elements $y\in\mathit{LL}$ with $y=0$ or $v_w(y)$ bounded below by a sum of local valuations minus $C_1$ times the error term
of the point, and such that the size bound at the type prime holds with $C_1$.

\begin{lemma}\label{pf:rows:haspart}
Let $\mathcal S$ be a slot system. Every core arising from a column choice through a built record of $\mathcal S$ has the
part property, and every core at a built position of $\mathcal S$ is at a position with the part property.
\formallabel{GlobalTheory.Svd.builtC_hasPart}
\end{lemma}

\begin{proof}
A built record determines a part (the part of the built record), which gives the first assertion; the second follows
because positions are monotone in the property defining them.
\end{proof}

\begin{proposition}[The arithmetic row from row pins]\label{pf:rows:atgpin}
Let $Q$ be a slot and $\rowsLk$ a linkage for a bundle type such that every core satisfying $Q$ lies at a position with the part
property at which all row pins hold. Then $\mathsf{aAtG}$ holds for $Q$ and $\rowsLk$.
\formallabel{GlobalTheory.T.holds_at_bundlePoints_of_trueColumns}
\formallabel{GlobalTheory.AtG.row_aAtG_of_rowPin}
\formallabel{GlobalTheory.AtG.row_atG_pinned}
\formallabel{GlobalTheory.T.atGRows_bundlePoints_of_jetSplitAmbient}
\end{proposition}

Here the \emph{row pin} at a tuple $a$ of a column choice $w$ with later tuple $Z$ asserts that for every word $u$ there is a linear
isomorphism $e$ from the space of the period data of $(w,u,a)$ onto the ambient space of the datum at $a$, equal to the standard
identification (the twisted one if $a\ne\mu_0$), such that every generator of the lattice $\widetilde M(a)$ is the image under $e$
of a row form, and a vector lies in the integral row lattice if and only if $e$ maps its row form to an
$\OO_{\mathit{LL}}$-combination of these generators.

\begin{proof}
Three ingredients are used.

(1) \emph{Jet tables at split primes} (\lref{GlobalTheory.T.jetTableBound_splitPrime}). For the ambient theta family $A$
attached to a seam, a curve ideal $\mathfrak c$, a twelfth root $\delta$ of the discriminant of the lattice $\mathfrak c$, a
character $\chi$ and constants $c_D$, for even ranks $a,b$, self-dual lattices $N,N'$ standard at $p$, a prime $p\ge5$ prime to
$d$ and to the discriminant of $E$, $|\iota(c_D(\tau))|\le1$, the lattice hypothesis, and $\pi\in\OO_K$ with $\pi\bar\pi=p$ and the
splitting condition: for $c'>0$ prime to $p$, every $x$ in the divided lattice $c'^{-1}\mathrm{Hom}_S(N,N')$ and every
monomial $m$, $|\iota\,\sigma_0(\mathrm{jet}_{pc'}(x,m))|\cdot|\iota(\bar\pi)|^{\deg m}\cdot|p|^{|\mathrm T|ab}\le1$.

(2) \emph{Period valuations at the type primes} (\lref{GlobalTheory.T.typePrimePeriodRow_of_constructedPointChains}).
Assume (1). If the ties of the concrete chains with the datum at the tuples $\mu_0$ and $\mu_w$ (for the places $w$) hold
eventually at linked bundles, then the row holds whose requirement is: closedly on $\rowsWiii(\mathcal A)$ there is $C_1$ such
that eventually the used tuples are bounded by $C_1$ times the error term of the point, and the size bound at the type prime
holds with $C_1$. This requirement is $\mathsf{aAtG}$.

(3) \emph{Ties from row pins} (\lref{GlobalTheory.T.concreteChainAgreesAtTuple_of_normalisedRows}). At every bundle
point whose core lies at a position with the part property at which all row pins hold, the ties hold at every tuple.

By the hypothesis on $Q$ and (3), the ties hold at all bundle points of $Q$, in particular eventually at linked bundles, and
(2) applies by (1).
\end{proof}

\begin{proposition}\label{pf:rows:atg}
Let $\mathcal S$ be a slot system whose row position is the fixed position $R_X$, and $\rowsLk$ a linkage of bundles. Then
$\mathsf{aAtG}$ holds under $\rowsLg(\rowsKPx,\rowsBuilt(\mathcal S),\rowsGood(\mathcal S),\cdot)$.
\formallabel{GlobalTheory.Fc.row_aAtG_at_built_alt}
\end{proposition}

\begin{proof}
Take the auxiliary slot $\mathrm{True}$. Every core at a good position of $\mathcal S$ is at a built position, hence comes from
a column choice with the part property (\cref{pf:rows:haspart}); since the row position is $R_X$, its
row forms match the period map (\lref{GlobalTheory.T.rowFormsMatchPeriodMap_of_furtherProperties}), so all row pins
hold. For slots at whose cores all row pins hold, the row $\mathsf{aAtG}$ holds by \cref{pf:rows:atgpin}.
\end{proof}

% ---------------------------------------------------------------------------
\subsection{The remaining rows}\label{sec:pf:rows:small}

\begin{proposition}[Small rows]\label{pf:rows:small}
Let $Q$ be a slot and $\rowsLk$ a linkage (in (a)--(c) for an arbitrary bundle type, otherwise for bundles).
\begin{enumerate}[label=(\alph*)]
\item $\mathsf{hn4}$ holds: closedly on $\rowsWi(\mathcal A)$, at all points, if the datum is large, the type primes split
and the level prime $T_u$ of the word $u$ splits, then the number of classes at level $0$ attached to $u$ equals
$\mathrm{mass}_k\cdot[\text{index at }g_0,g_1]\cdot(T_u^k-1)^f$. The same holds for the $\lambda$-row, which asserts that
eventually the logarithmic size parameter of $L_{x,i}$ is non-negative.
\item If the mass ratio identity holds for all data, then $\mathsf{hInstR}$ holds.
\item $\mathsf{aC2}$ holds.
\item If $\mathsf{hPA}$ holds, then $\mathsf{jOn}$ holds.
\item Under (1)--(4) of \cref{pf:rows:arch}, $\mathsf{hRB}$ holds.
\item If the comparison row $\mathsf{tie}'$ holds, then $\mathsf{hTw}$ holds.
\item $\mathsf{zL}$ holds.
\end{enumerate}
\formallabel{GlobalTheory.Fc.R0.classCountBound_paramsLambda_constructedPoints}
\formallabel{GlobalTheory.row_hInstR_of_hU}
\formallabel{GlobalTheory.Fc.ClassTwoEntries.divisionSecondComplement_constructedPoints}
\formallabel{GlobalTheory.Fc.JOnBundle.row_jOn_of_row_hPA_alt}
\formallabel{GlobalTheory.Fc.R1.row_hRB_univ_alt}
\formallabel{GlobalTheory.Fc.row_hTw_of_tie_alt}
\formallabel{GlobalTheory.Fc.R1.ZLAt.zL_of_hPA}
\end{proposition}

\begin{proof}
(a) The class number formula holds at every datum with every level datum (\lref{GlobalTheory.ColClassNumber.classNumber_all});
hence it holds closedly on every class (with $\eps_1=1$ and an arbitrary threshold), and the row follows by
\lref{GlobalTheory.Fc.Row.of_sel}. The second row holds in the same way, since the logarithmic size parameter is the
logarithm of a product of two quantities which are at least $1$ at all points with $x>x_1$
(\lref{GlobalTheory.famP_La_nonneg}).
(b) The mass ratio identity gives, for every datum, the volume--signature independence and the class items in the rank of
$Y_0$ (\lref{GlobalTheory.T.volume_signature_indep_and_classNumberItems}); conclude by \cref{pf:rows:calculus}(d).
(c) The members of $\rowsWiii(\mathcal A)$ satisfy $\sfname{TuSize}$ and $\sfname{DGamSize}$, and for $\eps<\frac1{52}$ these size
conditions imply the second-complement inputs (\lref{GlobalTheory.T.admissibleComplementTwo_along_family}); so the
inputs hold closedly on $\rowsWiii(\mathcal A)$ with $\eps_1=\frac1{100}$, and \cref{pf:rows:calculus}(d) applies.
(d) A primitive, type-separated realisation of the selected bundle provides the data of
\lref{GlobalTheory.JOnRow.row_jOn_of_realisations_alt}, which deduces that $\iota_{\mu_0}(V_0(\mu_0))$ is the whole space.
(e) For every bundle $b$, the realisation of $b$ satisfies the four statements of
\lref{GlobalTheory.T.indefinite_scale_theta_equality_alt} (existence at the base point, the bar condition for the
$J$-eigenvalues, the geometric action and the injectivity of evaluation), which use (1)--(4) of \cref{pf:rows:arch};
together with rapid decay and the size condition $n_V\ge8fr+4$ of the family parameters,
\lref{GlobalTheory.T.indefinite_scale_theta_bound_right_alt} gives bounded rows. The row holds for all selections.
(f) A comparison for the selected bundle and the twist-support property of the family give the tie of the twisting
records (\lref{Formalism.L2.Top4.tie_tw_rec_of_real_rec}).
(g) For every bundle $b$, the level projections $\pi_k$ at $\mu_0$ are surjective
(\lref{GlobalTheory.Fc.R1.ZLAt.lPrj_surj_at_alt}, from the surjectivity of the lifted projections and the column and
division data of $b$); hence the requirement holds closedly with $\eps_1=1$ and $x_0=0$ for all selections.
\end{proof}

% ---------------------------------------------------------------------------
\subsection{Weierstrass curves over ultrametric fields, CM invariants and unitary approximation}\label{sec:pf:rows:aux}

This subsection contains six auxiliary results of the formal proof of the row package which are used in the proof of the
integrality of the CM kernels (Weierstrass curves with good reduction, the group law of the Weierstrass function, the
integrality of singular $j$-invariants, and approximation in unitary groups).

\paragraph{Weierstrass curves.}
Let $K$ be a field with an ultrametric absolute value $|\cdot|$. A Weierstrass curve
$W:\ y^2+a_1xy+a_3y=x^3+a_2x^2+a_4x+a_6$ over $K$ is \emph{integral with unit discriminant} if $|a_i|\le1$ for all $i$ and
$|\Delta(W)|=1$. Its points are the origin $O$ and the non-singular solutions $(x,y)\in K^2$, with the chord--tangent group
law; for a point $P=(x,y)\ne O$ we write $x(P)=x$, $y(P)=y$ and
\[
\psi(P)=2y+a_1x+a_3,\qquad \kappa(P)=3x^2+2a_2x+a_4-a_1y
\]
(the partial derivatives of the Weierstrass equation, up to sign), and we call $P$ \emph{integral} if $|x|,|y|\le1$. Recall
that $-P=(x,-y-a_1x-a_3)$; if $x(P)\ne x(Q)$, the line through $P,Q$ has slope $\lambda=(y(P)-y(Q))/(x(P)-x(Q))$ and
$x(P+Q)=\lambda^2+a_1\lambda-a_2-x(P)-x(Q)$; if $P=Q$ and $\psi(P)\ne0$, the tangent has slope $\lambda=\kappa(P)/\psi(P)$
and $x(2P)=\lambda^2+a_1\lambda-a_2-2x(P)$.

\begin{lemma}\label{pf:rows:quad}
If $t,b,c\in K$ satisfy $t^2+bt=c$ and $|b|,|c|\le1$, then $|t|\le1$.
\end{lemma}

\begin{proof}
If $|t|>1$, then $|t|^2=|c-bt|\le\max(|c|,|b||t|)\le|t|$, a contradiction.
\end{proof}

\begin{lemma}[Non-singular reduction]\label{pf:rows:ns}
Let $W$ be integral with unit discriminant and let $(x,y)\in K^2$ satisfy the Weierstrass equation, with $|x|,|y|\le1$. Then
$|2y+a_1x+a_3|=1$ or $|3x^2+2a_2x+a_4-a_1y|=1$.
\formallabel{Formalism.L2.Th1.E.nonsingular_norm}
\formallabel{Formalism.L2.Th1.E.nonsingNorm}
\end{lemma}

\begin{proof}
Write $\psi=2y+a_1x+a_3$, $\kappa=3x^2+2a_2x+a_4-a_1y$, $b=a_1^2+4(a_2+3x)$ and $u=2\kappa+a_1\psi$. A direct computation shows
that
\begin{equation}\label{eq:pf:rows:disc}
\Delta(W)=\psi\bigl(-b^2(-a_1\kappa+(a_2+3x)\psi)-8a_1u^2-27\psi^3+9bu\psi\bigr)+\kappa\bigl(b^2\kappa-16u^2\bigr)
\end{equation}
for every solution $(x,y)$ of the Weierstrass equation: the difference of the two sides is the product of
$y^2+a_1xy+a_3y-x^3-a_2x^2-a_4x-a_6$ with a polynomial with integer coefficients in $a_1,\dots,a_6,x,y$. All quantities in
\eqref{eq:pf:rows:disc} are integral, and $|\psi|,|\kappa|\le1$. If both were $<1$, the ultrametric inequality would give
$|\Delta(W)|<1$.
\end{proof}

\begin{lemma}[Integral torsion points]\label{pf:rows:unitpts}
Let $W$ be integral with unit discriminant, and let $P,Q\ne O$ be integral points.
\begin{enumerate}[label=(\alph*)]
\item If $2P=O$, then $\psi(P)=0$ and $|\kappa(P)|=1$.
\item If $2P\ne O$ and $x(2P)$ is integral, then $|\psi(P)|=1$.
\item If $P+Q\ne O$, $P-Q\ne O$ and $x(P+Q)$, $x(P-Q)$ are integral, then $|x(P)-x(Q)|=1$.
\end{enumerate}
\formallabel{Formalism.L2.Th1.E.UnitKap_of}
\formallabel{Formalism.L2.Th1.E.UnitPsi2_of}
\formallabel{Formalism.L2.Th1.E.UnitDiff_of}
\formallabel{Formalism.L2.Th1.E.unitKap}
\formallabel{Formalism.L2.Th1.E.unitDiff}
\end{lemma}

\begin{proof}
Each part is deduced from the conclusion of \cref{pf:rows:ns} alone (formally, (a)--(c) are stated as consequences of
that statement). Let $P=(x_1,y_1)$, $Q=(x_2,y_2)$, $\psi=\psi(P)$, $\kappa=\kappa(P)$.

(a) If $y_1\ne-y_1-a_1x_1-a_3$, the tangent at $P$ is not vertical and $2P\ne O$; hence $y_1=-y_1-a_1x_1-a_3$, i.e.\
$\psi=0$. By \cref{pf:rows:ns}, $|\psi|=1$ or $|\kappa|=1$, and the first is impossible.

(b) Since $2P\ne O$, $\psi\ne0$ and the tangent slope $\lambda=\kappa/\psi$ satisfies
$\lambda^2+a_1\lambda=x(2P)+a_2+2x_1$, whose right-hand side is integral; by \cref{pf:rows:quad}, $|\lambda|\le1$. If
$|\psi|<1$, then $|\kappa|=1$ by \cref{pf:rows:ns}, while $|\kappa|=|\lambda||\psi|<1$. Hence $|\psi|=1$ (as $|\psi|\le1$).

(c) First $x_1\ne x_2$: otherwise subtracting the two Weierstrass equations gives $(y_1-y_2)(y_1+y_2+a_1x_2+a_3)=0$, so
$P=Q$ or $P=-Q$, contrary to the hypotheses. Let $\lambda_\pm$ be the slopes of the lines through $P$ and $\pm Q$:
\[
\lambda_+(x_1-x_2)=y_1-y_2,\qquad \lambda_-(x_1-x_2)=y_1+y_2+a_1x_2+a_3 .
\]
Since $x(P\pm Q)=\lambda_\pm^2+a_1\lambda_\pm-a_2-x_1-x_2$ is integral, \cref{pf:rows:quad} gives $|\lambda_\pm|\le1$.
Subtracting the Weierstrass equations,
$(y_1-y_2)(y_1+y_2+a_1x_2+a_3)=(x_1-x_2)\bigl(x_1^2+x_1x_2+x_2^2+a_2(x_1+x_2)+a_4-a_1y_1\bigr)$, so
$\lambda_+\lambda_-(x_1-x_2)=x_1^2+x_1x_2+x_2^2+a_2(x_1+x_2)+a_4-a_1y_1$, and therefore
\[
\psi=(x_1-x_2)(\lambda_++\lambda_-+a_1),\qquad \kappa=(x_1-x_2)(\lambda_+\lambda_-+2x_1+x_2+a_2).
\]
If $|x_1-x_2|<1$, both $|\psi|<1$ and $|\kappa|<1$, contradicting \cref{pf:rows:ns}. Hence $|x_1-x_2|=1$.
\end{proof}

\paragraph{The group law of the Weierstrass function.}
Let $L\subset\C$ be a lattice with invariants $g_2,g_3$ and Weierstrass function $\wp=\wp_L$, let $\ell$ be a prime and
$\iota:\C\to\C_\ell$ a field isomorphism. Let $W^\iota_L$ be the curve $y^2=x^3-\frac{\iota(g_2)}4x-\frac{\iota(g_3)}4$ over
$\C_\ell$ and define $P(z)=O$ for $z\in L$ and $P(z)=(\iota\wp(z),\frac12\iota\wp'(z))$ for $z\notin L$ (a non-singular point of
$W^\iota_L$, by the differential equation of $\wp$ and $g_2^3-27g_3^2\ne0$,
\lref{Formalism.L2.Th1.E.AddThm.nonsing}).

\begin{lemma}\label{pf:rows:ptofadd}
For all $s,t\in\C$, $P(s+t)=P(s)+P(t)$ in the group of points of $W^\iota_L$.
\formallabel{Formalism.L2.Th1.E.PtOfAdd_of}
\formallabel{Formalism.L2.Th1.E.ptOf_add}
\formallabel{Formalism.L2.Th1.E.ptOf_add_generic}
\formallabel{Formalism.L2.Th1.E.ptOf_add_lattice}
\formallabel{Formalism.L2.Th1.E.ptOf_lattice}
\formallabel{Formalism.L2.Th1.E.ptOf_neg}
\end{lemma}

\begin{proof}
Since $\wp$ is even and $L$-periodic and $\wp'$ is odd and $L$-periodic, $P(t+\omega)=P(t)$ for $\omega\in L$ and
$P(-t)=-P(t)$.

\emph{Generic case: $s,t,s\pm t\notin L$.} Then $\wp(s)\ne\wp(t)$ (\lref{Formalism.L2.Th1.E.AddThm.P_sub_ne}), and the
addition theorem gives
$\wp(s+t)=\frac14\bigl(\frac{\wp'(s)-\wp'(t)}{\wp(s)-\wp(t)}\bigr)^2-\wp(s)-\wp(t)$
(\lref{Formalism.L2.Th1.E.AddThm.P_add}) and the collinearity relation
$(\wp(s)-\wp(t))\wp'(s+t)=-(\wp'(s)-\wp'(t))(\wp(s+t)-\wp(s))-\wp'(s)(\wp(s)-\wp(t))$
(\lref{Formalism.L2.Th1.E.AddThm.collinear}). Apply $\iota$. The chord through $P(s)$ and $P(t)$ has slope
$\lambda=\iota\frac{\wp'(s)-\wp'(t)}{2(\wp(s)-\wp(t))}$; the first identity says $\iota\wp(s+t)=\lambda^2-\iota\wp(s)-\iota\wp(t)$, and
the second, divided by $2(\wp(s)-\wp(t))$, says $\frac12\iota\wp'(s+t)=-\lambda(\iota\wp(s+t)-\iota\wp(s))-\frac12\iota\wp'(s)$. These are
the formulas for $P(s)+P(t)$.

\emph{The other cases.} If $s\in L$ or $t\in L$, the assertion follows from periodicity and $P(\omega)=O$. If $s+t\in L$, then
$P(t)=P(-s)=-P(s)$ and both sides are $O$. If $s-t\in L$ (and $s,t,s+t\notin L$), then $P(s)=P(t)$, $P(s+t)=P(2t)$ and
$2t\notin L$. There is $u$ with $2u,t+u,t-u\notin L$ (one of $\omega_1/3$, $\omega_1/5$ for a basis vector $\omega_1$ of $L$;
\lref{Formalism.L2.Th1.E.AddThm.exists_finite_part_one_idele_split_congruence_quotient}). The generic case applies to
the pairs $(t+u,t-u)$, $(t,u)$ and $(t,-u)$, and gives
$P(2t)=P(t+u)+P(t-u)=P(t)+P(u)+P(t)+P(-u)=2P(t)=P(s)+P(t)$.
\end{proof}

\paragraph{Singular $j$-invariants.}
For a lattice $L$ put $j(L)=1728g_2^3/(g_2^3-27g_3^2)$. For a squarefree $d\equiv3\pmod4$ let $\omega_d=(1+i\sqrt d)/2$.

\begin{lemma}\label{pf:rows:jint}
\begin{enumerate}[label=(\alph*)]
\item If $\ell$ is a prime and $x\in\C_\ell$ is a root of a monic polynomial with integer coefficients, then $|x|\le1$.
\item Let $d$ be squarefree with $d\equiv3\pmod 4$ and $L$ a lattice with $\omega_dL\subseteq L$. Then $|\iota(j(L))|\le1$ for
every prime $\ell$ and every field isomorphism $\iota:\C\to\C_\ell$.
\end{enumerate}
\formallabel{ReciprocityTheory.ChartRestriction.L1.norm_le_one_of_root}
\formallabel{ReciprocityTheory.ChartRestriction.L1.jInt}
\end{lemma}

\begin{proof}
(a) Let $P=X^n+\sum_{i<n}c_iX^i$; $n\ge1$ since $P$ has a root. If $|x|>1$, then
$|x|^n=|\sum_{i<n}c_ix^i|\le\max_i|c_i||x|^i\le|x|^{n-1}<|x|^n$, because $|c_i|\le1$.
(b) $j(L)$ is an algebraic integer (\lref{ReciprocityTheory.ChartRestriction.L1.cm1}; classically
\cite[Ch.~II, Thm.~6.1]{Silverman1994}), so $\iota(j(L))$ is a root of a monic integer polynomial; apply (a).
\end{proof}

\paragraph{Approximation in unitary groups.}
Let $E$ be a number field with an involution $x\mapsto x^*$ such that $E$ is totally complex, every embedding $E\to\C$
transforms ${}^*$ into complex conjugation, and for every maximal ideal $\mathfrak P$ of $\OO_E$ there is $\theta\in\OO_E$ with
$\theta-\theta^*\notin\mathfrak P$. Let $F$ be the fixed field, $h(x,y)=\sum_ix_i^*y_i$ on $E^n$ and $\Uni(h)$ its isometry
group. For a maximal ideal $\mathfrak P$ let $S_{\mathfrak P}$ be the multiplicative monoid generated by the
$s\in\OO_E$ with $s\notin\mathfrak P$ and $s^*\notin\mathfrak P$, and for an $\OO_E$-submodule $M\subseteq E^n$ let
$M_{(\mathfrak P)}=\{x\in E^n:sx\in M\text{ for some }s\in S_{\mathfrak P}\}$.

\begin{proposition}\label{pf:rows:wa}
Let $p$ be a prime and $L\subseteq E^n$ an $\OO_E$-submodule. Suppose that for every maximal ideal $\mathfrak P\ni p$ of $\OO_E$
there is $g_{\mathfrak P}\in\Uni(h)$ with $g_{\mathfrak P}\bigl((\OO_E^n)_{(\mathfrak P)}\bigr)=L_{(\mathfrak P)}$. Then there is
$\gamma\in\Uni(h)$ with $\gamma\bigl((\OO_E^n)_{(\mathfrak P)}\bigr)=L_{(\mathfrak P)}$ for all maximal ideals $\mathfrak P\ni p$.
\formallabel{Formalism.L2.KerInert.weak_approximation_places}
\end{proposition}

\begin{proof}
We use the following statement on unitary groups (\lref{unitary_approximation_finite}): if $\sigma$ is an
automorphism of order $2$ of $E$ with fixed field $F$, $J$ an invertible $\sigma$-hermitian matrix, $T$ a finite set of
finite places of $F$, $g_v$ ($v\in T$) elements of the unitary group of $J$ over $E\otimes_FF_v$ and $0\ne D\in\OO_F$, then
there is $\gamma$ in the unitary group of $J$ over $E$ with $(g_v)_{ij}-\gamma_{ij}\in D\cdot R_v$ for all $v\in T$ and all
$i,j$, where $R_v\subseteq E\otimes_FF_v$ is the subring generated by $\OO_E\otimes1$ and $1\otimes\OO_{F_v}$.

\emph{Step 1 (localisations).} For $s\in\OO_E$ the norm $ss^*$ lies in $\OO_F$, and $s,s^*\notin\mathfrak P$ if and only if
$ss^*\notin\mathfrak P\cap\OO_F$. Hence $S_{\mathfrak P}$, and with it $M_{(\mathfrak P)}$, depends only on $\mathfrak P\cap\OO_F$.

\emph{Step 2 (local integrality).} For a prime $v$ of $\OO_F$ call $e\in E$ \emph{$v$-integral} if $se\in\OO_E$ for some
$s\in\OO_F\smallsetminus v$. The $v$-integral elements form a subring stable under ${}^*$ (as $s^*=s$). If a matrix $A$
has $v$-integral entries and $\mathfrak P$ lies above $v$, then $A$ maps $(\OO_E^n)_{(\mathfrak P)}$ into itself: if $sx\in\OO_E^n$
with $s\in S_{\mathfrak P}$ and $tA\in M_n(\OO_E)$ with $t\in\OO_F\smallsetminus v$, then $tsAx\in\OO_E^n$ and
$t=t^*\notin\mathfrak P$, so $ts\in S_{\mathfrak P}$.

\emph{Step 3 (traces).} For $b\in\OO_E$ the map $\tau_b(e)=eb+(eb)^*$ is $F$-linear $E\to F$ and maps $\OO_E$ into $\OO_F$;
its $F_v$-linear extension maps $R_v$ into $\OO_{F_v}$. Let $\theta\in\OO_E$ and $\delta=\theta-\theta^*$. If $x\in E$ and
$x\otimes1\in R_v$, then $\tau_b(x)\in F\cap\OO_{F_v}$ is $v$-integral for every $b$, and
$\delta x=\tau_\theta(x)-\tau_1(x)\theta^*$ is $v$-integral.

\emph{Step 4 (construction).} Since $p$ is not a unit in $\OO_E$ there is a maximal ideal $\mathfrak P_0\ni p$; choose
$\theta$ with $\delta=\theta-\theta^*\notin\mathfrak P_0$, so $\delta\ne0$. Let $T$ be the set of primes of $\OO_F$ containing
$p$. For $v\in T$ choose a maximal ideal $\mathfrak P_v\ni p$ of $\OO_E$ above $v$ and let $G_v$ be the matrix of
$g_{\mathfrak P_v}$, so $G_v^*G_v=1$ (if there is no such ideal put $G_v=1$). Choose $0\ne c\in\Z$ with $c\delta^{-1}\in\OO_E$
and $c\,G_v\in M_n(\OO_E)$ for all $v\in T$. The approximation statement with $J=1$ and $D=c^2$ gives $\gamma\in\Uni(h)$ with
$(G_v-\gamma)_{ij}\in c^2R_v$.

\emph{Step 5 (local step).} Fix $v\in T$ and put $Y=c^{-1}(G_v-\gamma)$. Then $Y_{ij}=(c\delta^{-1})\cdot\delta x_{ij}$ with
$x_{ij}=c^{-2}(G_v-\gamma)_{ij}$, $x_{ij}\otimes1\in R_v$, so $Y$ is $v$-integral by Step 3. Since $\gamma=G_v-cY$ and
$G_v^*G_v=1$,
\[
G_v^*\gamma=1-(cG_v)^*Y,\qquad \gamma^*G_v=1-Y^*(cG_v)
\]
have $v$-integral entries.

\emph{Step 6.} Let $\mathfrak P\ni p$ be maximal and $v=\mathfrak P\cap\OO_F\in T$. By Step 1,
$G_v\bigl((\OO_E^n)_{(\mathfrak P)}\bigr)=L_{(\mathfrak P)}$. By Steps 2 and 5,
$\gamma(\OO_E^n)_{(\mathfrak P)}=G_v(G_v^*\gamma)(\OO_E^n)_{(\mathfrak P)}\subseteq L_{(\mathfrak P)}$ and
$L_{(\mathfrak P)}=\gamma(\gamma^*G_v)(\OO_E^n)_{(\mathfrak P)}\subseteq\gamma(\OO_E^n)_{(\mathfrak P)}$.
\end{proof}

% ---------------------------------------------------------------------------
\subsection{Remarks on the formal statements}\label{sec:pf:rows:remarks}

\begin{remark}\label{pf:rows:fidelity}
\begin{enumerate}[label=(\arabic*)]
\item Six of the twenty-six rows ($\mathsf{hPA}$, $\mathsf{hLamP}$, $\mathsf{hDivG}$, $\mathsf{hDcPtG}$, $\mathsf{selJ}$, $\mathsf{nRest}$)
are hypotheses of \cref{pf:rows:main}, and in $\mathsf{rRows}$ the definite-side condition, $\sfname{L3}$, $\sfname{FUN}$, the
multiplicity condition and the $\kappa$-rows are hypotheses as well.
\item The rows $\mathsf{hn4}$, $\mathsf{aC2}$, $\mathsf{hRB}$, $\mathsf{zL}$, and, given the Siegel--Weil formula, $\mathsf{hInstR}$,
$\mathsf{sLp}$, $\mathsf{sO}$, $\mathsf{hTF}$, hold for every slot and every linkage, and the six period rows for every slot system;
for them the open parameter statements play no role. The hypothesis $\mathsf{hPA}$ in the formal statement of
\cref{pf:rows:small}(g) is not used.
\item A row only constrains classes whose members are built, point by point, from linked bundles. In the transfer
(\cref{pf:rows:calculus}(e)) the closedness is witnessed with $\eps_1=1$, an arbitrary threshold and $x_0=0$: the
property holds at all points with $x>x_1$, not only eventually.
\item In the counting statement used by the families, the datum $D_{x,i}$ at which the statement is ``evaluated'' plays no
role: the statement quantifies over all automorphic instances. Similarly, the additional predicate on the data in
\cref{pf:rows:countconst} is not used in the conclusion.
\item In the definite case ($q=0$) of \cref{pf:rows:counting} the count equals the main term exactly
(\cref{pf:rows:definite}).
\item The statement $\sfname{VolumeStep}$ exposes only the hermitian matrices and signatures of $I'$ (and not the relation
between the lattices of $I$ and $I'$), and holds for all gauge constants and Haar measures.
\item The bound for the residue of $\zeta_E$ used in \cref{pf:rows:kappa}, $\rowsRes\le(16\log r_d)^{2f}$, is weaker than the
explicit bounds of Louboutin \cite{Louboutin2000}, which are of the form $\rowsRes\le(e\log d_E/(2[E:\Q]-2))^{[E:\Q]-1}$.
\item In \cref{pf:rows:closingsharp} the twisted column bound is obtained for every tuple $\mu$; the restrictions in its
definition (used tuples, $\mu\ne\mu_0$, residue characteristic outside $Q_p(\mu)$) are not used.
\item The three parts of \cref{pf:rows:unitpts} are stated formally as consequences of the non-singular reduction statement
of \cref{pf:rows:ns}, and \cref{pf:rows:unitperiods} takes the Euler product identity and the algebraicity of the Taylor
coefficients of torsion translates as hypotheses; these are discharged by \cref{pf:rows:ns} and
\lref{GlobalTheory.MomEq.eq_and_alg} respectively.
\item In \cref{pf:rows:wa}, of the hypotheses on $E$ only the existence, for one maximal ideal $\mathfrak P_0\ni p$, of
$\theta\in\OO_E$ with $\theta-\theta^*\notin\mathfrak P_0$ is used; that $E$ is totally complex and that the involution induces
complex conjugation are not used.
\end{enumerate}
\end{remark}
